\documentclass[12pt,a4paper]{book}

\usepackage[a4paper,top=25mm,bottom=25mm,inner=28mm,outer=22mm]{geometry}
\usepackage{amsmath,amssymb}
\usepackage{graphicx}
\usepackage{xcolor}
\usepackage{booktabs}
\usepackage{longtable}
\usepackage{enumitem}
\usepackage{array}
\usepackage{tikz}
\usetikzlibrary{positioning}
\usepackage[most]{tcolorbox}
\usepackage{fancyhdr}
\usepackage[colorlinks=true,linkcolor=blue!45!black,urlcolor=blue!45!black,citecolor=blue!45!black]{hyperref}
\usepackage{xepersian}

\settextfont[
  Path={../assets/fonts/},
  UprightFont={IRLotusICEE.ttf},
  BoldFont={IRLotusICEE_Bold.ttf},
  BoldItalicFont={IRLotusICEE_BoldIranic.ttf},
  ItalicFont={IRLotusICEE_Iranic.ttf},
  FontFace={b}{n}{IRLotusICEE_Bold.ttf},
  FontFace={bx}{n}{IRLotusICEE_Bold.ttf},
  FontFace={b}{it}{IRLotusICEE_BoldIranic.ttf},
  FontFace={bx}{it}{IRLotusICEE_BoldIranic.ttf},
  FontFace={m}{sl}{IRLotusICEE_Iranic.ttf},
  Scale=1.2
]{IRLotusICEE.ttf}
\setlatintextfont[
  Path={../assets/fonts/},
  BoldFont={LiberationSerif-Bold.ttf},
  BoldItalicFont={LiberationSerif-BoldItalic.ttf},
  ItalicFont={LiberationSerif-Italic.ttf},
  Scale=1
]{LiberationSerif-Regular.ttf}
\setdigitfont[
  Path={../assets/fonts/},
  Scale=1.2
]{IRLotusICEE.ttf}
\defpersianfont\titlefont[
  Path={../assets/fonts/},
  FontFace={b}{n}{IRTitr.ttf},
  FontFace={bx}{n}{IRTitr.ttf},
  FontFace={m}{it}{IRTitr.ttf},
  Scale=1
]{IRTitr.ttf}
\defpersianfont\nast[Path={../assets/fonts/}, Scale=2]{IranNastaliq.ttf}

\linespread{1.22}
\setlength{\parskip}{0.45em}
\setlength{\parindent}{0pt}
\setcounter{tocdepth}{2}
\setlist[itemize]{topsep=0.2em,itemsep=0.15em}
\setlist[enumerate]{topsep=0.2em,itemsep=0.15em}

\definecolor{mainblue}{RGB}{24,75,125}
\definecolor{softblue}{RGB}{235,243,252}
\definecolor{softgreen}{RGB}{238,248,241}
\definecolor{softyellow}{RGB}{255,248,226}
\definecolor{softred}{RGB}{254,238,238}
\definecolor{boxgray}{RGB}{245,246,248}

\setlength{\headheight}{16pt}
\pagestyle{fancy}
\fancyhf{}
\fancyhead[LE,RO]{\small\thepage}
\fancyhead[RE]{\small\nouppercase{\leftmark}}
\fancyhead[LO]{\small\nouppercase{\rightmark}}
\renewcommand{\headrulewidth}{0.45pt}
\renewcommand{\footrulewidth}{0pt}
\fancypagestyle{plain}{
  \fancyhf{}
  \fancyfoot[C]{\small\thepage}
  \renewcommand{\headrulewidth}{0pt}
  \renewcommand{\footrulewidth}{0pt}
}

\newtcolorbox{definitionbox}[1]{
  enhanced,
  breakable,
  colback=softblue,
  colframe=mainblue,
  title={#1},
  fonttitle=\bfseries,
  boxrule=0.7pt,
  arc=1mm,
  right=2mm,left=2mm,top=1mm,bottom=1mm
}

\newtcolorbox{assignmentbox}[1]{
  enhanced,
  breakable,
  colback=softyellow,
  colframe=orange!65!black,
  title={#1},
  fonttitle=\bfseries,
  boxrule=0.7pt,
  arc=1mm,
  right=2mm,left=2mm,top=1mm,bottom=1mm
}

\newtcolorbox{correctionbox}[1]{
  enhanced,
  breakable,
  colback=softgreen,
  colframe=green!45!black,
  title={#1},
  fonttitle=\bfseries,
  boxrule=0.7pt,
  arc=1mm,
  right=2mm,left=2mm,top=1mm,bottom=1mm
}

\newtcolorbox{warningbox}[1]{
  enhanced,
  breakable,
  colback=softred,
  colframe=red!55!black,
  title={#1},
  fonttitle=\bfseries,
  boxrule=0.7pt,
  arc=1mm,
  right=2mm,left=2mm,top=1mm,bottom=1mm
}

\newcommand{\eng}[1]{\lr{#1}}
\newcommand{\term}[2]{\textbf{#1} \eng{(#2)}}
\newcommand{\course}{علوم شناختی محاسباتی}
\newcommand{\coursefigure}[3]{%
\begin{figure}[htbp]
  \centering
  \includegraphics[width=#1\textwidth]{../assets/figures/#2}
  \caption{#3}
\end{figure}
}

\title{{\titlefont\Huge \course}\\[0.5em]
{\Large از نظریه‌های شناختی تا مدل‌های محاسباتی زبان و گفتار}}
\author{دکتر زینب آقاهادی\\[0.4em]
دانشگاه علم و صنعت ایران}
\date{بهار ۱۴۰۵}

\begin{document}
\raggedbottom
\frontmatter
\maketitle

\begingroup
\thispagestyle{empty}
\begin{latin}
\centering
\vspace*{0.16\textheight}
{\Huge\bfseries Computational Cognitive Science\par}
\vspace{3em}
{\Large From Cognitive Theories to Computational Models of Language and Speech\par}
\vfill
{\Large\bfseries Zeinab Aghahadi\par}
\vspace{1.2em}
{\large Assistant Professor\par}
{\large School of Computer Engineering\par}
{\large Iran University of Science and Technology (IUST)\par}
\vfill
{\large Spring 2026\par}
\vspace{0.5em}
{\large Version 1.0\par}
\end{latin}
\clearpage
\endgroup

\begingroup
\thispagestyle{empty}
{\titlefont\LARGE چکیده\par}
\vspace{1.5em}
علوم شناختی محاسباتی با ساخت مدل‌های ریاضی، فرایندهای ذهنی و عصبی را مطالعه می‌کند و نشان می‌دهد چگونه می‌توان نظریه‌های شناختی را به مدل‌های محاسباتی آزمون‌پذیر برای زبان، گفتار و تصمیم‌گیری تبدیل کرد. مباحث از مبانی علوم شناختی و معماری‌های شناختی آغاز می‌شوند و سپس حافظه، بار شناختی، پیش‌بینی، سوگیری و توجه را در سامانه‌های انسانی و مصنوعی بررسی می‌کنند. در ادامه، مدل‌سازی ادراک شنیداری و گفتار، سازوکارهای پاداش و یادگیری تقویتی، و شبکه‌های عصبی اسپایکی و محاسبات نورومورفیک معرفی می‌شوند. فصل‌های پایانی نیز به مدل‌های اسپایکی نوین، کدگذاری و رمزگشایی عصبی، هم‌ترازی مغز و مدل، و شباهت خطاهای استدلالی انسان و مدل‌های زبانی می‌پردازند. هدف اصلی، ایجاد پلی میان مفاهیم نظری علوم شناختی و ابزارهای محاسباتی معاصر و فراهم‌کردن چارچوبی برای تحلیل، طراحی و ارزیابی مدل‌های شناختی است.

\vfill
\textbf{کلیدواژه‌ها:} علوم شناختی محاسباتی، مدل‌سازی شناختی، زبان و گفتار، معماری شناختی، یادگیری تقویتی، شبکه‌های عصبی اسپایکی، هم‌ترازی مغز و مدل
\clearpage
\endgroup

\begingroup
\thispagestyle{empty}
\begin{latin}
{\LARGE\bfseries Abstract\par}
\vspace{1.5em}
Computational cognitive science studies mental and neural processes by constructing mathematical models and explains how cognitive theories can be translated into testable computational models of language, speech, and decision-making. The discussion begins with the foundations of cognitive science and cognitive architectures, then examines memory, cognitive load, prediction, bias, and attention in human and artificial systems. Subsequent chapters introduce computational approaches to auditory and speech processing, reward mechanisms and reinforcement learning, spiking neural networks, and neuromorphic computing. The final chapters discuss modern spiking models, neural encoding and decoding, brain-model alignment, and similarities between human reasoning errors and those of language models. The overall aim is to connect theoretical concepts in cognitive science with contemporary computational tools and to provide a framework for analysing, designing, and evaluating cognitive models.

\vfill
\textbf{Keywords:} computational cognitive science, cognitive modelling, language and speech, cognitive architecture, reinforcement learning, spiking neural networks, brain-model alignment
\end{latin}
\clearpage
\endgroup

\tableofcontents

\mainmatter

\chapter{مبانی علوم شناختی محاسباتی}

\section{چارچوب مفهومی و مسئلهٔ اصلی}

\subsection{علوم شناختی چیست؟}
علوم شناختی فقط مطالعهٔ مغز به‌معنای زیستی و کالبدشناختی آن نیست. پرسش اصلی علوم شناختی این است که انسان چگونه می‌تواند \term{کنش}{Act} داشته باشد: چگونه می‌بیند، می‌شنود، حرف می‌زند، فکر می‌کند، تصمیم می‌گیرد، قضاوت می‌کند، یاد می‌گیرد، خطا می‌کند و گاهی دچار اختلال می‌شود. بنابراین موضوع علوم شناختی، فهم سازوکارهایی است که رفتار و تجربهٔ انسانی را ممکن می‌کنند.

در این نگاه، «انسان» فقط فرد بزرگسال سالم و متوسط نیست. نظریهٔ شناختی باید بتواند دربارهٔ کودک، سالمند، فرد بسیار توانمند، فرد دارای ناتوانی یا اختلال، و فردی با عملکرد عادی توضیحی منسجم بدهد. اگر نظریه‌ای فقط برای یک گروه محدود جواب بدهد، هنوز نظریهٔ عمومی خوبی دربارهٔ شناخت انسان نیست.

برای توضیح این نکته می‌توان از تعبیر «ماژول شناختی» استفاده کرد. منظور از ماژول در اینجا یک قطعهٔ سخت‌افزاری جدا و دقیقاً مرزبندی‌شده نیست؛ بلکه مجموعه‌ای از فرایندها و سازوکارهاست که باعث می‌شود انسان بتواند ادراک، زبان، حافظه، تصمیم‌گیری و عمل را انجام دهد. علوم شناختی می‌کوشد این سازوکارها را توصیف، آزمایش و در نهایت مدل کند.

\subsection{از علوم شناختی تا علوم شناختی محاسباتی}
در علوم شناختی، ممکن است یک نظریه در سطح کلامی یا مفهومی توضیح دهد که انسان چگونه یک مهارت را انجام می‌دهد. علوم شناختی محاسباتی یک گام جلوتر می‌رود و می‌پرسد آیا می‌توان آن نظریه را به شکل صوری، محاسباتی و قابل پیاده‌سازی بیان کرد یا نه.

پس هدف فقط این نیست که بگوییم «انسان از زمینه استفاده می‌کند» یا «حافظهٔ کاری محدود است». هدف این است که بتوانیم این ادعاها را به مدل، الگوریتم، معماری، آزمایش و پیش‌بینی تبدیل کنیم. یک نظریهٔ شناختی زمانی برای علوم شناختی محاسباتی مفیدتر می‌شود که بتوان آن را در قالبی اجراپذیر یا دست‌کم آزمون‌پذیر بازنویسی کرد.

به همین دلیل، بحث از نظریه‌های کلاسیک شناختی شروع می‌شود، اما در آنجا متوقف نمی‌ماند. مسیر اصلی این است که هر نظریهٔ شناختی به یک مسئلهٔ محاسباتی وصل شود: آیا این نظریه می‌تواند ضعف یک مدل هوش مصنوعی را توضیح دهد؟ آیا می‌تواند راهی برای طراحی معماری بهتر، دادهٔ بهتر یا ارزیابی بهتر پیشنهاد کند؟

\subsection{نسبت علوم شناختی محاسباتی با هوش مصنوعی}
مسیر دانشی این حوزه را می‌توان به‌صورت تقریبی چنین دید:
\[
\begin{aligned}
\text{\lr{Artificial Intelligence}}
&\rightarrow
\text{\lr{Machine Learning}}
\rightarrow
\text{\lr{Deep Learning}}\\
&\rightarrow
\text{\lr{Computational Cognitive Science}}.
\end{aligned}
\]

در هوش مصنوعی کلاسیک، با روش‌هایی مانند جست‌وجو، \eng{brute force}، \eng{BFS}، الگوریتم ژنتیک و \eng{simulated annealing} روبه‌رو می‌شویم. در یادگیری ماشین، مسئلهٔ یادگیری از داده و استخراج ویژگی برجسته می‌شود. در یادگیری عمیق، مدل‌ها بازنمایی‌های پیچیده‌تر و معماری‌های عمیق‌تر می‌سازند. در علوم شناختی محاسباتی، پرسش تازه‌ای اضافه می‌شود: کدام مهارت شناختی در این مدل‌ها غایب است و چگونه می‌توان آن را وارد مدل کرد؟

در بسیاری از دوره‌های رایج علوم شناختی محاسباتی، نمونهٔ اصلی یا \eng{case study} از حوزهٔ بینایی و پردازش تصویر انتخاب می‌شود. این انتخاب تاریخی و آموزشی قابل فهم است، چون هم علوم اعصاب و روان‌شناسی شناختی داده‌های کلاسیک زیادی دربارهٔ بینایی دارند و هم هوش مصنوعی در دوره‌ای رشد بسیار مهمی را از بینایی ماشین تجربه کرده است. با این حال، در این منبع تمرکز اصلی روی زبان، گفتار، استدلال و مدل‌های زبانی است؛ یعنی حوزه‌هایی مانند:
\begin{itemize}
  \item \term{پردازش زبان طبیعی}{Natural Language Processing / NLP}
  \item \term{استدلال}{Reasoning}
  \item \term{پردازش گفتار}{Speech Processing}
  \item \term{مدل‌های زبانی}{Language Models}
  \item \term{ترنسفورمر}{Transformer}
\end{itemize}

برای دنبال‌کردن این بحث‌ها، آشنایی پایه با \eng{NLP} مفید است، اما بسیاری از مدل‌ها و معماری‌های لازم در متن مرور می‌شوند. چون وزن اصلی مثال‌ها روی زبان و گفتار است، خواننده‌ای که از حوزه‌ای دیگر وارد می‌شود بهتر است نسبت مفاهیم مطرح‌شده را با مسئلهٔ پژوهشی یا کاربردی خود صریح کند.

\subsection{تفاوت پرسش شناختی و پرسش یادگیری ماشین}
در یادگیری ماشین، مسئله معمولاً به شکل زیر صورت‌بندی می‌شود:
\[
\text{برای این ورودی، چگونه خروجی درست را پیش‌بینی کنیم؟}
\]
در چنین چارچوبی، معیارهایی مانند \eng{accuracy}، \eng{precision}، \eng{recall} یا معیارهای مشابه اهمیت زیادی دارند. اگر مدل ورودی را بگیرد و خروجی درست بدهد، از نگاه مهندسی ممکن است موفق تلقی شود؛ حتی اگر ندانیم دقیقاً در درون مدل چه اتفاقی افتاده است.

اما در علوم شناختی، پرسش اصلی کمی متفاوت است:
\[
\text{انسان چگونه این کار را انجام می‌دهد؟}
\]
بنابراین درست بودن خروجی کافی نیست. یک مدل شناختی باید تا حدی دربارهٔ فرایند، زمان واکنش، الگوی خطا، تفاوت‌های فردی و اختلالات نیز توضیح بدهد. اگر مدلی در یک آزمون امتیاز بالا بگیرد، اما خطاهایش هیچ شباهتی به خطاهای انسانی نداشته باشد، شاید ابزار مهندسی خوبی باشد، اما الزاماً مدل شناختی خوبی نیست.

همین تفاوت باعث می‌شود علوم شناختی محاسباتی پلی میان دو دغدغه باشد: از یک طرف می‌خواهد مدل محاسباتی بسازد، و از طرف دیگر نمی‌خواهد فقط به خروجی نهایی قانع شود. مدل باید از نظر رفتاری و در صورت امکان از نظر فرایندی با شواهد انسانی قابل مقایسه باشد.

\subsection{نظریهٔ شناختی، مدل محاسباتی و اعتبارسنجی}
برای توضیح تعدد نظریه‌های علوم شناختی می‌توان از استعارهٔ فیل استفاده کرد. گویی پدیدهٔ شناخت انسان مانند فیلی بزرگ است و هر نظریه فقط بخشی از آن را لمس می‌کند: یک نظریه شنیدن را توضیح می‌دهد، نظریه‌ای دیگر حافظه را، نظریه‌ای دیگر توجه را، و نظریه‌ای دیگر تصمیم‌گیری یا خطا را. هیچ نظریه‌ای به‌تنهایی تمام ذهن انسان را توضیح نمی‌دهد.

یک نظریهٔ شناختی خوب باید چند ویژگی داشته باشد:
\begin{itemize}
  \item رفتار درست انسان را توضیح دهد.
  \item خطاهای انسان را نیز توضیح دهد، نه اینکه فقط پاسخ‌های درست را مدل کند.
  \item فقط به یک سن، یک وضعیت یا یک گروه محدود وابسته نباشد.
  \item دربارهٔ تفاوت‌های فردی و اختلالات شناختی حرفی برای گفتن داشته باشد.
  \item پیش‌بینی تولید کند و با آزمایش‌های مستقل سنجیده شود.
\end{itemize}

اعتبارسنجی در علوم شناختی معمولاً با یک مثال قانع‌کننده تمام نمی‌شود. باید آزمایش طراحی شود، دادهٔ انسانی جمع‌آوری شود، تحلیل آماری انجام شود و نظریه در موقعیت‌های مستقل دوباره سنجیده شود. اگر نظریه‌ای می‌گوید انسان در تشخیص گفتار ابتدا چند نامزد واژگانی را فعال می‌کند، باید بتواند توضیح دهد چرا در یک موقعیت نویزی یا مبهم، فرد به‌جای یک کلمه، کلمهٔ دیگری را می‌شنود.

نکتهٔ مهم این است که نظریه‌های شناختی معمولاً به شکل صفر و یکی کنار گذاشته نمی‌شوند. بسیاری از آن‌ها بخشی از پدیده را خوب توضیح می‌دهند، اما بعداً نظریه‌های کامل‌تر یا دقیق‌تر می‌آیند و بخش‌های جاافتاده را اضافه می‌کنند. بنابراین تکامل نظریه‌ها در علوم شناختی بیشتر شبیه اصلاح، گسترش و دقیق‌تر کردن است تا صرفاً رد کامل نظریهٔ قبلی.

\subsection{مغز، نواحی کارکردی و احتیاط در تفسیر}
نظریه‌های شناختی باید با شواهد مغزی نیز تا حدی سازگار باشند، اما نباید مغز را به مجموعه‌ای از جعبه‌های کاملاً جدا و بدون هم‌پوشانی تبدیل کرد. قشر مغز نواحی مختلفی دارد که با کارکردهایی مانند بینایی، شنوایی، گفتار، تعادل، هماهنگی، تصمیم‌گیری و استدلال مرتبط‌اند. برای نمونه، نواحی پس‌سری بیشتر با پردازش بینایی، نواحی گیجگاهی با بخش‌هایی از پردازش شنیداری و زبان، و نواحی پیشانی با کارکردهایی مانند برنامه‌ریزی، قضاوت و کنترل شناختی ارتباط دارند.

با این حال، این نگاشت‌ها مرزهای کاملاً تیز ندارند. کارکردهای شناختی معمولاً حاصل شبکه‌هایی از نواحی مغزی‌اند، نه فقط یک نقطهٔ منفرد. پس وقتی می‌گوییم «این ناحیه با گفتار مرتبط است» یا «این ناحیه با استدلال مرتبط است»، منظور یک رابطهٔ غالب و تقریبی است، نه مالکیت انحصاری آن ناحیه بر یک توانایی.

\subsection{کاربرد علوم شناختی و محدودهٔ این منبع}
یکی از کاربردهای مهم علوم شناختی، فهم اختلالات و بیماری‌هاست. اگر بفهمیم یک اختلال شناختی یا رشدی، مانند بخشی از طیف اوتیسم، با چه سازوکارهایی در ادراک، زبان، تعامل اجتماعی یا پیش‌بینی محیط مرتبط است، می‌توانیم مداخله‌های بهتر و انسانی‌تری طراحی کنیم. البته این منبع جایگزین آموزش تخصصی درمان اختلالات یا نوروسایکولوژی بالینی نیست.

هدف این منبع ساختن پلی میان علوم شناختی و هوش مصنوعی است. بنابراین از نظریه‌های شناختی استفاده می‌کنیم تا ضعف‌های مدل‌های محاسباتی را بهتر بفهمیم، آزمایش‌های دقیق‌تری طراحی کنیم و در صورت امکان، مدل‌هایی بسازیم که از نظر شناختی معنادارتر باشند.

\subsection{شکست مدل‌های بزرگ در مهارت‌های شناختی}
مدل‌های زبانی بزرگ در بسیاری از وظایف زبانی و دانشی عملکرد چشمگیری دارند، اما این موفقیت به‌تنهایی نشان نمی‌دهد که آن‌ها مهارت‌های شناختی را به شکل انسانی و عمومی یاد گرفته‌اند. یکی از مسیرهای مهم ارزیابی این مدل‌ها، طراحی آزمون‌هایی است که شباهت سطحی به داده‌های آموزشی را کم می‌کند و از مدل می‌خواهد رابطهٔ انتزاعی را در موقعیتی تازه تشخیص دهد.

نمونهٔ مهم، آزمون‌های \term{استدلال قیاسی}{Analogical Reasoning} با تغییرات \term{پادواقعی}{Counterfactual} است. اگر مدل در قیاس‌های آشنا خوب عمل کند، مثلاً رابطهٔ جمع بستن در «\eng{book} به \eng{books}» و «\eng{girl} به \eng{girls}» را تشخیص دهد، هنوز معلوم نیست مفهوم انتزاعی رابطه را فهمیده یا فقط الگوهای آشنا را بازتولید کرده است. برای آزمون سخت‌تر، می‌توان الفبا، نمادها، رنگ‌ها یا ترتیب‌ها را عوض کرد و دید آیا مدل هنوز رابطهٔ «بعدی بودن»، «معکوس بودن» یا «هم‌رده بودن» را منتقل می‌کند یا نه.

مثال‌هایی مانند الفبای ساختگی، داستان‌هایی با زمینه‌های متفاوت اما \eng{theme} مشترک، و جایگزینی دسته‌ها با نمادهایی مانند ضربدر یا علامت تعجب همین مسئله را نشان می‌دهند. نکتهٔ مشترک همهٔ این مثال‌ها این است که مدل نباید فقط به شباهت ظاهری تکیه کند؛ باید رابطهٔ عمیق‌تر را تشخیص دهد. پژوهش‌های جدید دربارهٔ قیاس در مدل‌های زبانی دقیقاً در همین نقطه بحث دارند: برخی نتایج افت شدید مدل‌ها را در نسخه‌های پادواقعی نشان می‌دهند، و برخی کارهای دیگر ادعا می‌کنند مدل‌های قوی‌تر یا طراحی‌های دقیق‌تر می‌توانند بخشی از این تعمیم را انجام دهند. بنابراین صورت دقیق ادعا باید محتاطانه باشد: قیاس انتزاعی و پادواقعی هنوز یکی از آزمون‌های مهم و چالش‌برانگیز برای سنجش شناختی مدل‌های زبانی است.

\subsection{تغییرات \eng{Counterfactual} و \eng{Adversarial}}
وقتی ترتیب الفبا را عوض می‌کنیم، نمادهای تازه تعریف می‌کنیم، یا یک رابطهٔ آشنا را در محیطی ناآشنا می‌گذاریم، یک تغییر \eng{counterfactual} ساخته‌ایم. اگر تغییر طوری طراحی شود که مدل را از الگوهای سطحی و حفظ‌شده دور کند، می‌تواند نقش \eng{adversarial} نیز داشته باشد؛ یعنی برای آشکار کردن ضعف مدل به کار رود.

برای انسان، تعمیم یک رابطهٔ انتزاعی به محیط تازه معمولاً طبیعی‌تر است. البته انسان هم همیشه بی‌خطا نیست، اما اگر رابطه را فهمیده باشد، می‌تواند آن را از الفبای واقعی به الفبای ساختگی منتقل کند. مدل‌هایی که بیشتر بر همبستگی‌های آماری آشنا تکیه دارند، ممکن است با همین تغییر کوچک دچار افت شدید شوند. این افت همان جایی است که پرسش شناختی آغاز می‌شود: کدام مهارت انتزاعی یا سازوکار تعمیم در مدل ضعیف است؟

\subsection{نقش دادهٔ انسانی و همبستگی خطا}
برای ارزیابی شناختی مدل‌ها، فقط نباید بپرسیم آیا مدل \eng{task} را حل کرده است یا نه. باید بپرسیم الگوی عملکرد مدل چقدر با الگوی عملکرد انسان همبستگی دارد.

دو معیار را باید از هم جدا کرد:
\begin{enumerate}
  \item آیا مدل در انجام وظیفه موفق است؟
  \item آیا سختی‌ها و خطاهای مدل با سختی‌ها و خطاهای انسان \eng{correlation} دارد؟
\end{enumerate}

این تمایز به هدف ما بستگی دارد. اگر هدف ساختن ابزاری باشد که خطاهای انسان را جبران کند، شاید نخواهیم مدل همان خطاهای انسان را تکرار کند. اما حتی در این حالت هم دانستن نسبت خطای مدل و انسان مهم است: ابتدا باید بفهمیم آیا مدل و انسان در یک خطای شناختی مشترک‌اند یا نه، سپس می‌توانیم با \eng{fine-tuning}، \eng{few-shot learning}، تغییر داده، تغییر معماری یا روش‌های مشابه آن خطا را کاهش دهیم.

اگر هدف مدل‌سازی شناخت انسان باشد، همبستگی با خطای انسانی اهمیت بیشتری پیدا می‌کند. مدلی که فقط پاسخ‌های درست بیشتری می‌دهد، اما هیچ شباهتی به الگوی خطای انسان ندارد، الزاماً توضیح‌دهندهٔ شناخت انسان نیست. به همین دلیل، در علوم شناختی محاسباتی گاهی خطاها به اندازهٔ پاسخ‌های درست مهم‌اند.

در بحث \term{هوش مصنوعی عمومی}{Artificial General Intelligence / AGI} نیز همین مسئله دیده می‌شود. برخی چارچوب‌های جدید ارزیابی، به‌جای سنجش یک مهارت واحد، مجموعه‌ای از توانایی‌های شناختی را بررسی می‌کنند. یکی از چارچوب‌های روان‌سنجی مهم برای دسته‌بندی توانایی‌های شناختی انسان، \term{نظریهٔ کتل-هورن-کرول}{Cattell-Horn-Carroll / CHC} است که توانایی‌ها را در سطوح گسترده و محدود طبقه‌بندی می‌کند. استفاده از چنین چارچوب‌هایی برای ارزیابی مدل‌های هوش مصنوعی به این معناست که مدل فقط در یک بنچمارک خاص سنجیده نشود، بلکه پروفایل شناختی آن در چند توانایی مختلف بررسی شود.

\subsection{مسیر پژوهشی}
یک مسیر پژوهشی معمول در علوم شناختی محاسباتی چنین است:
\begin{enumerate}
  \item یک ضعف یا خطای شناختی در مدل پیدا کنیم.
  \item بررسی کنیم آیا این ضعف با خطای انسان همبستگی دارد یا نه.
  \item خطای انسانی را با یک نظریهٔ شناختی توضیح دهیم.
  \item داده، آزمایش یا توصیف مناسب برای سنجش آن طراحی کنیم.
  \item مدل را با روش‌هایی مانند \eng{prompting}، \eng{few-shot learning}، \eng{fine-tuning}، \eng{RAG} یا تغییر معماری اصلاح کنیم.
  \item نشان دهیم اصلاح مدل فقط افزایش سطحی عملکرد نیست، بلکه با تحلیل شناختی مسئله ارتباط دارد.
\end{enumerate}

اگر بدون مرحلهٔ شناختی فقط مدل را بزرگ‌تر کنیم یا روی دادهٔ بیشتری آموزش دهیم، ممکن است بهبود مهندسی به دست آید، اما سهم علوم شناختی روشن نیست. ارزش پژوهشی زمانی بیشتر می‌شود که ابتدا ضعف مدل در قالب یک مهارت شناختی توضیح داده شود و سپس راهکار محاسباتی دقیقاً همان ضعف را هدف بگیرد.

\chapter{معماری شناختی، حافظه و بار شناختی}

\section{معماری شناختی}

\subsection{چرا معماری شناختی لازم است؟}
پیش از ورود به جزئیات معماری‌های شناختی، یک مثال ساده نشان می‌دهد چرا صرفاً داشتن حسگر، قانون و خروجی درست کافی نیست. اگر بخواهیم مدل هوش مصنوعی به مهارتی نزدیک شود که انسان در زندگی روزمره به‌سادگی انجام می‌دهد، باید بدانیم چه اجزایی در فرایند شناختی دخیل‌اند: ادراک، حافظه، دانش عامه، پیش‌بینی پیامد، تصمیم‌گیری، عمل و دریافت بازخورد.

معماری شناختی دقیقاً برای همین لازم می‌شود. معماری شناختی فقط فهرستی از قابلیت‌ها نیست؛ بلکه باید روشن کند اطلاعات از کجا وارد می‌شود، در کجا نگه داشته می‌شود، با کدام دانش قبلی ترکیب می‌شود، چگونه تصمیم ساخته می‌شود و نتیجهٔ عمل چگونه به سیستم برمی‌گردد.

\subsection{مثال جارو رباتیک}
مثال مناسب، جاروبرقی رباتیک است. یک ربات ساده ممکن است چند حسگر و چند قانون داشته باشد:
\begin{itemize}
  \item اگر به گردوغبار رسید، آن را جمع کند.
  \item اگر مانعی دید، مسیرش را تغییر دهد.
  \item اگر به یک شیء کوچک رسید، آن را کنار بزند.
\end{itemize}

این قواعد در بسیاری از موقعیت‌های ساده کافی‌اند. اما اگر ربات به فنجانی پر از مایع برسد، مسئله تغییر می‌کند. انسان معمولاً فنجان را بدون فکر هل نمی‌دهد، چون چند رابطهٔ علّی و پیامدی را می‌فهمد: فنجان ممکن است پر باشد؛ اگر هل داده شود، مایع می‌ریزد؛ اگر مایع بریزد، زمین خیس می‌شود؛ خیس شدن زمین نامطلوب است؛ پس نباید فنجان را مانند یک مانع معمولی کنار زد.

این مثال نشان می‌دهد مسئله فقط تشخیص شیء نیست. ربات باید چیزی شبیه دانش عامه، پیش‌بینی پیامد، ارزیابی ریسک و انتخاب عمل داشته باشد. بدون این مهارت‌ها، سیستم ممکن است در معیارهای سادهٔ مهندسی خوب به نظر برسد، اما در موقعیت‌های انسانی و پرجزئیات رفتار نامناسبی تولید کند.

\subsubsection{راه‌حل اول: دوری از اشیای پرریسک}
یک راه ساده این است که ربات اشیای پرریسک را نادیده بگیرد یا از آن‌ها دور شود. این راه‌حل خسارت را کم می‌کند، اما مسئله را کامل حل نمی‌کند. اگر ربات هر شیء مشکوکی را رها کند، ممکن است بخش‌هایی از محیط هرگز تمیز نشوند. بنابراین با یک \eng{trade-off} روبه‌رو هستیم: کاهش خطر در برابر ناقص‌ماندن وظیفه.

\subsubsection{راه‌حل دوم: تعریف تابع ریسک}
راه دیگر این است که برای حرکت دادن هر شیء، یک تابع ریسک تعریف کنیم. مدل می‌تواند بپرسد: شیء چیست؟ ظرف است یا نه؟ شکستنی است یا نه؟ داخل آن مایع هست یا نه؟ احتمال ریختن چقدر است؟ پیامد ریختن چیست؟

این مسیر از نظر مهندسی ممکن است مفید باشد، اما اگر فقط به افزودن دائمی \eng{if-else}های تازه منجر شود، سیستم شکننده و پرهزینه می‌شود. هر موقعیت تازه ممکن است قانون تازه‌ای بخواهد.

\subsubsection{راه‌حل سوم: یادگیری از تجربه و بازخورد}
راه سوم این است که ربات از تجربه و \eng{feedback} یاد بگیرد. کودک نیز از ابتدا همهٔ قوانین فیزیکی و اجتماعی را به‌صورت صریح ندارد؛ تجربه می‌کند، پیامدها را می‌بیند و از بازخورد مثبت یا منفی یاد می‌گیرد. اگر فنجان بیفتد و آب بریزد، این تجربه می‌تواند در رفتارهای بعدی اثر بگذارد.

از نگاه علوم شناختی محاسباتی، پرسش مهم این است که این یادگیری فقط حفظ نمونه‌ها نباشد، بلکه به فهم روابط علّی و تعمیم به موقعیت‌های تازه نزدیک شود.

\subsubsection{دو مسیر برای بهبود مدل}
از این مثال دو مسیر کلی برای بهبود مدل‌های هوش مصنوعی به دست می‌آید:
\begin{enumerate}
  \item مدل را بزرگ‌تر کنیم و قوانین یا داده‌های بیشتری به آن بدهیم.
  \item یک مهارت یا ماژول شناختی به مدل اضافه کنیم.
\end{enumerate}

مسیر اول می‌تواند در عمل مفید باشد، اما همیشه خطر دارد که فقط به انباشتن قوانین سطحی تبدیل شود. مسیر دوم می‌پرسد کدام توانایی شناختی کم است: فهم علّیت، حافظه، توجه، پیش‌بینی پیامد، ارزیابی ریسک یا یادگیری از بازخورد؟ این همان نقطه‌ای است که علوم شناختی به طراحی مدل هوش مصنوعی وصل می‌شود.

\subsection{ماژول‌های اصلی در یک معماری شناختی}
از مثال جارو رباتیک می‌توان اجزای اولیهٔ یک معماری شناختی ساده را استخراج کرد. ربات فقط یک ورودی و خروجی ندارد؛ میان ورودی و خروجی چند سطح پردازش لازم است. ابتدا باید محیط دیده یا حس شود، سپس وضعیت فعلی در حافظهٔ موقت قرار بگیرد، بعد دانش مرتبط از حافظهٔ بلندمدت بازیابی شود، سپس یک قاعده یا مهارت مناسب انتخاب شود و در نهایت دستور حرکتی تولید گردد.

معماری شناختی یعنی همین سازمان‌دهی یکپارچهٔ فرایندها. اگر چند نظریهٔ شناختی جداگانه داشته باشیم، معماری مشخص می‌کند این نظریه‌ها چگونه در یک سیستم واحد با هم تعامل می‌کنند: اطلاعات از کجا وارد می‌شود، کدام ماژول آن را پردازش می‌کند، چه چیزی در حافظه می‌ماند، چه چیزی بازیابی می‌شود و عمل چگونه تولید می‌شود.

\subsubsection{ادراک و ماژول بینایی}
ماژول بینایی یا ادراکی، ورودی محیط را به بازنمایی قابل استفاده برای سیستم تبدیل می‌کند. در مثال ربات، این ماژول باید تشخیص دهد که یک فنجان وجود دارد، داخل آن احتمالاً آب است و فنجان در مسیر حرکت قرار گرفته است.

در نگاشت تقریبی به مغز، پردازش بینایی با نواحی بینایی، به‌ویژه لوب پس‌سری، ارتباط دارد. البته در مغز واقعی، بینایی فقط در یک نقطهٔ جداگانه انجام نمی‌شود و مسیرهای گسترده‌تری برای تشخیص شیء، مکان، حرکت و معنا درگیرند. اما برای سطح فعلی معماری، می‌توان ماژول بینایی را ورودی ادراکی سیستم دانست.

\subsubsection{حافظه و بافرهای کاری}
اطلاعاتی که از محیط می‌آید باید موقتاً نگه داشته شود. در معماری‌های شناختی معمولاً از مفهوم \term{بافر}{Buffer} استفاده می‌شود: فضای محدودی که وضعیت فعلی سیستم را در دسترس سایر ماژول‌ها قرار می‌دهد. بافر را می‌توان به حافظهٔ کوتاه‌مدت یا حافظهٔ کاری نزدیک دانست، هرچند در معماری‌هایی مانند \eng{ACT-R} معنای فنی دقیق‌تری دارد.

برای مثال، پس از دیدن فنجان، محتوای بافر می‌تواند چیزی شبیه این باشد: «فنجان در مسیر است» و «مایع داخل فنجان دیده می‌شود». این اطلاعات هنوز به‌تنهایی تصمیم ایجاد نمی‌کند؛ فقط وضعیت فعلی محیط را برای بازیابی دانش و انتخاب عمل آماده می‌کند.

\subsubsection{تصمیم‌گیری و انتخاب عمل}
پس از قرار گرفتن وضعیت فعلی در بافر، سیستم باید تصمیم بگیرد چه عملی مناسب است. این تصمیم از ترکیب چند منبع به دست می‌آید: وضعیت فعلی، دانش اعلانی دربارهٔ اشیا و پیامدها، و قواعد یا مهارت‌های رویه‌ای.

اگر سیستم بداند فنجان شکستنی است، آب ممکن است بریزد و خیس شدن زمین نامطلوب است، آنگاه یک قاعدهٔ رویه‌ای می‌تواند فعال شود: اگر فنجان پر از آب در مسیر است، با آن مثل مانع معمولی برخورد نکن. این مرحله همان جایی است که تفاوت یک مدل صرفاً واکنشی با یک مدل شناختی‌تر آشکار می‌شود.

\subsubsection{ماژول حرکتی}
ماژول حرکتی تصمیم مفهومی را به عمل فیزیکی تبدیل می‌کند. «تغییر مسیر بده» باید به دستورهایی مانند کاهش سرعت، چرخش، حرکت بازو، گرفتن آرام فنجان یا بازگشت به مسیر قبلی تبدیل شود.

در نگاشت مغزی، کنترل حرکت با بخش‌هایی از لوب پیشانی و مدارهای حرکتی مرتبط است. باید میان \eng{frontal lobe} به‌عنوان ناحیهٔ وسیع‌تر و \eng{prefrontal cortex} به‌عنوان بخشی که بیشتر با برنامه‌ریزی، استدلال، کنترل اجرایی و تصمیم‌گیری مرتبط است تمایز گذاشت.

\subsection{نگاشت معماری شناختی به مغز}
یکی از شروط مهم معماری شناختی این است که اجزای آن دست‌کم به‌صورت تقریبی با مغز قابل نگاشت باشند. اگر معماری یک ماژول بینایی، ماژول حرکتی، حافظهٔ اعلانی و حافظهٔ رویه‌ای دارد، باید بتوان توضیح داد این اجزا با چه شبکه‌ها یا نواحی مغزی ارتباط دارند.

این نگاشت نباید ساده‌انگارانه باشد. مغز شبکه‌ای و تعاملی عمل می‌کند و هیچ مرز کاملاً تیزی میان همهٔ کارکردها وجود ندارد. اما نگاشت تقریبی برای اعتبار شناختی مدل مهم است. برای نمونه:
\begin{longtable}{>{\raggedleft\arraybackslash}p{0.34\textwidth}p{0.50\textwidth}}
\toprule
\textbf{جزء معماری} & \textbf{نگاشت تقریبی مغزی} \\
\midrule
ماژول بینایی & نواحی بینایی، به‌ویژه لوب پس‌سری \\
ماژول حرکتی & نواحی حرکتی در لوب پیشانی و مدارهای مرتبط \\
تصمیم‌گیری و کنترل شناختی & نواحی پیش‌پیشانی \\
حافظهٔ اعلانی & هیپوکامپ و ساختارهای لوب گیجگاهی میانی \\
حافظهٔ رویه‌ای & عقده‌های قاعده‌ای و مدارهای قشری-زیرقشری مرتبط \\
\bottomrule
\end{longtable}

این جدول یک نقشهٔ آموزشی است، نه ادعای جداسازی کامل. هرکدام از این کارکردها با شبکه‌های گسترده‌تر و تعامل میان چند ناحیه انجام می‌شوند.

\subsection{\eng{ACT-R}}
\term{ACT-R}{Adaptive Control of Thought-Rational} یکی از معماری‌های شناختی اثرگذار است که عمدتاً با کارهای \eng{John R. Anderson} و همکارانش شناخته می‌شود. ریشه‌های تاریخی آن به نظریه‌های \eng{ACT} در دههٔ ۱۹۷۰ برمی‌گردد، اما خود \eng{ACT-R} در نسخه‌های بعدی توسعه یافته و همچنان در مدل‌سازی شناختی استفاده می‌شود.

ایدهٔ کلی \eng{ACT-R} این است که شناخت انسان را می‌توان با چند جزء اصلی مدل کرد: ماژول‌ها، بافرها، حافظهٔ اعلانی، حافظهٔ رویه‌ای و سازوکار تطبیق الگو. در این معماری، دانش اعلانی معمولاً به شکل \eng{chunk}ها و دانش رویه‌ای به شکل \eng{production rule}ها نمایش داده می‌شود. وضعیت لحظه‌ای سیستم در محتوای بافرها دیده می‌شود و \eng{pattern matcher} قاعده‌ای را پیدا می‌کند که با وضعیت فعلی سازگار است.

در مثال جارو رباتیک، مسیر ساده‌شده چنین است:
\[
\begin{aligned}
\text{\lr{Environment}}
&\rightarrow
\text{\lr{Visual Module}}
\rightarrow
\text{\lr{Buffer}}\\
&\rightarrow
\text{\lr{Memory Retrieval}}
\rightarrow
\text{\lr{Production}}
\rightarrow
\text{\lr{Motor Module}}.
\end{aligned}
\]

\coursefigure{0.52}{act-r-architecture.png}{نمای شماتیک معماری \eng{ACT-R}: محیط از مسیر ماژول‌های بینایی و حرکتی با بافرها، حافظهٔ اعلانی، حافظهٔ رویه‌ای، تطبیق الگو و اجرای \eng{production} در ارتباط است.}

\subsubsection{حافظهٔ اعلانی}
حافظهٔ اعلانی در \eng{ACT-R} محل دانش صریح و قابل بیان است. در مثال ربات، وقتی بافر شامل «فنجان» و «آب» است، حافظهٔ اعلانی می‌تواند دانشی مانند این را بازیابی کند:
\begin{itemize}
  \item فنجان می‌تواند شکستنی باشد.
  \item آب یک مایع است.
  \item برخورد با فنجان می‌تواند باعث ریختن آب شود.
  \item ریختن آب زمین را خیس می‌کند.
\end{itemize}

این نوع دانش برای توضیح‌پذیری نیز مهم است. اگر سیستم مسیرش را عوض کند، می‌توان گفت: «فنجان پر از آب در مسیر بود و برخورد با آن می‌توانست باعث ریختن آب شود.» این بخش به بحث \eng{interpretability} در هوش مصنوعی نزدیک می‌شود، چون تصمیم مدل به دانش قابل بیان متصل می‌شود.

\subsubsection{حافظهٔ رویه‌ای}
حافظهٔ رویه‌ای در \eng{ACT-R} شامل \eng{production}ها یا قواعدی از جنس اگر-آنگاه است. این قواعد شبیه مهارت‌هایی‌اند که می‌گویند در یک وضعیت خاص چه عملی مناسب است:
\[
\text{اگر فنجان پر از آب در مسیر بود، مسیر را تغییر بده.}
\]
یا:
\[
\text{اگر مانع شکستنی در مسیر بود، سرعت را کم کن و آرام جابه‌جایش کن.}
\]

در انسان، بسیاری از مهارت‌ها پس از تمرین زیاد سریع و کم‌هزینه اجرا می‌شوند. حافظهٔ رویه‌ای نیز دقیقاً به همین جنبهٔ «چگونه انجام دادن» مربوط است، نه فقط «چه چیزی دانستن».

\subsubsection{\eng{Pattern Matching}}
\eng{Pattern Matching} مرحله‌ای است که در آن وضعیت فعلی بافرها و اطلاعات بازیابی‌شده با قواعد حافظهٔ رویه‌ای مقایسه می‌شود. اگر وضعیت فعلی با شرط یک \eng{production} سازگار باشد، آن قاعده نامزد اجرا می‌شود.

برای مثال:
\begin{itemize}
  \item بافر: فنجان پر از آب در مسیر است.
  \item حافظهٔ اعلانی: فنجان شکستنی است و آب ممکن است بریزد.
  \item حافظهٔ رویه‌ای: اگر مانع پرخطر در مسیر بود، با آن مثل مانع معمولی برخورد نکن.
  \item نتیجه: قاعدهٔ تغییر مسیر یا جابه‌جایی آرام فعال می‌شود.
\end{itemize}

در \eng{ACT-R}، شناخت به‌صورت توالی‌ای از فعال‌شدن و اجرای \eng{production}ها پیش می‌رود. بنابراین فرایند فقط یک خروجی نهایی نیست؛ زنجیره‌ای از وضعیت‌ها و تغییر وضعیت‌هاست.

\subsubsection{\eng{Production Execution}}
پس از انتخاب قاعده، \eng{production} اجرا می‌شود. خروجی این مرحله باید بتواند وضعیت بافرها را تغییر دهد و در نهایت به دستور حرکتی تبدیل شود. برای ربات، «مسیر را تغییر بده» باید به دستورهایی مانند کاهش سرعت، چرخش چرخ‌ها یا فعال‌کردن بازو تبدیل شود.

نکتهٔ مهم این است که خروجی اجرای \eng{production} باید به حافظهٔ کوتاه‌مدت یا بافر برگردد. سیستم باید بداند چه تصمیمی گرفته و چه عملی را شروع کرده است. این بازگشت باعث می‌شود توالی عمل‌ها حفظ شود، عمل بعدی با عمل قبلی سازگار باشد و رفتار سیستم از هم گسسته نشود.

\subsection{اعتبارسنجی معماری شناختی}
اعتبارسنجی معماری شناختی فقط با مقایسهٔ خروجی انجام نمی‌شود. ممکن است انسان و مدل هر دو به پاسخ درست برسند، اما مسیر رسیدن آن‌ها کاملاً متفاوت باشد. در علوم شناختی، خود مسیر اهمیت دارد.

برای اعتبارسنجی، می‌توان یک \eng{task} مشخص طراحی کرد و آن را هم به انسان و هم به مدل داد. در انسان، ابزارهایی مانند \eng{fMRI}، \eng{EEG}، زمان واکنش یا \eng{eye-tracking} می‌توانند اطلاعاتی دربارهٔ توالی و زمان‌بندی فعالیت‌ها بدهند. در مدل نیز می‌توان دید کدام ماژول در هر مرحله فعال شده است.

اگر مدل ابتدا ماژول بینایی را فعال کند، سپس وضعیت را در بافر قرار دهد، بعد حافظهٔ اعلانی و رویه‌ای را درگیر کند و در نهایت عمل حرکتی تولید کند، باید بررسی شود آیا دادهٔ انسانی نیز جریان مشابهی را نشان می‌دهد یا نه. هرچه توالی فعال‌شدن اجزا، زمان‌بندی و الگوی خطا با انسان سازگارتر باشد، معماری از نظر شناختی معتبرتر است.

\section{حافظه در انسان و مدل}

\subsection{حافظهٔ حسی، کوتاه‌مدت و کاری}
حافظه در معماری شناختی فقط یک محل ذخیره‌سازی ساده نیست. اگر قرار است مدل مهارت شناختی داشته باشد، باید روشن شود چه نوع اطلاعاتی را موقت نگه می‌دارد، چه دانشی را در بلندمدت ذخیره می‌کند و چگونه میان این دو سطح رفت‌وبرگشت انجام می‌دهد.

حافظه را می‌توان ابتدا به‌صورت ساده به حافظهٔ کوتاه‌مدت و بلندمدت تقسیم کرد. حافظهٔ کوتاه‌مدت شبیه یک بافر یا \eng{cache} محدود است. اطلاعات تازه ابتدا وارد این فضا می‌شوند و ممکن است خیلی زود از بین بروند.

مثلاً اگر کسی آدرس اتاقی را برای لحظه‌ای بشنود، ممکن است همان روز از آن استفاده کند، اما اگر ماه بعد دوباره به آن نیاز داشته باشد، احتمالاً باید دوباره بپرسد. این مثال نشان می‌دهد حضور موقت یک اطلاعات در ذهن به معنی ورود پایدار آن به حافظهٔ بلندمدت نیست.

در ادبیات جدید، باید میان \term{حافظهٔ کوتاه‌مدت}{Short-Term Memory} و \term{حافظهٔ کاری}{Working Memory} تمایز گذاشت. حافظهٔ کوتاه‌مدت بیشتر بر نگهداری موقت تأکید دارد، در حالی که حافظهٔ کاری علاوه بر نگهداری، پردازش فعال اطلاعات فعلی را نیز شامل می‌شود. این دو مفهوم گاهی نزدیک به هم استفاده می‌شوند، اما از نظر نظری یکسان نیستند.

\subsection{حافظهٔ بلندمدت}
اطلاعاتی که با تمرین، تکرار، استفاده یا تجربهٔ مهم تثبیت می‌شوند، می‌توانند وارد حافظهٔ بلندمدت شوند. حافظهٔ بلندمدت یک مخزن یکنواخت نیست و انواع مختلفی دارد. دو دستهٔ اصلی در اینجا حافظهٔ اعلانی و حافظهٔ رویه‌ای هستند.

\subsubsection{حافظهٔ اعلانی}
\term{حافظهٔ اعلانی}{Declarative Memory} مربوط به دانشی است که می‌توان آن را به‌صورت صریح بیان کرد: واقعیت‌ها، مفاهیم، رویدادها و اطلاعاتی که می‌توان درباره‌شان گزارش داد. مثال‌ها:
\begin{itemize}
  \item دانستن اینکه آب زمین را خیس می‌کند.
  \item دانستن اینکه فنجان ممکن است شکستنی باشد.
  \item دانستن آدرس یک مکان.
  \item دانستن یک تعریف یا مفهوم علمی.
\end{itemize}

از نظر مغزی، حافظهٔ اعلانی به‌ویژه با هیپوکامپ و ساختارهای لوب گیجگاهی میانی ارتباط دارد، هرچند ذخیره و بازیابی دانش در شبکه‌های گسترده‌تری انجام می‌شود. بنابراین بهتر است نگوییم همهٔ حافظهٔ اعلانی «در هیپوکامپ ذخیره می‌شود»؛ دقیق‌تر این است که هیپوکامپ و نواحی پیرامونی برای شکل‌گیری و بازیابی بسیاری از خاطرات اعلانی نقش کلیدی دارند.

در مقایسهٔ آموزشی با مدل‌های زبانی، می‌توان گفت مدل‌هایی مانند \eng{ChatGPT} در بازیابی و ترکیب دانش زبانی و واقعیت‌های آموخته‌شده قوی به نظر می‌رسند؛ بنابراین از این جهت به حافظهٔ اعلانی تشبیه می‌شوند. این تشبیه کامل نیست، چون مدل زبانی حافظهٔ انسانی ندارد، اما برای فهم تفاوت «دانستن» و «انجام دادن» مفید است.

\coursefigure{0.50}{hippocampus-and-declarative-memory.png}{هیپوکمپ و ساختارهای پیرامونی آن در تثبیت و بازیابی حافظهٔ اعلانی نقش کلیدی دارند؛ این تصویر برای پیوند بحث حافظهٔ صریح با نگاشت مغزی آمده است.}

\subsubsection{حافظهٔ رویه‌ای}
\term{حافظهٔ رویه‌ای}{Procedural Memory} مربوط به مهارت‌ها، عادت‌ها و روش‌های انجام کار است. مثال‌ها:
\begin{itemize}
  \item راه رفتن
  \item رانندگی
  \item آشپزی
  \item ترمز گرفتن
  \item حرکت دادن یک ابزار یا بازوی ربات
\end{itemize}

حافظهٔ رویه‌ای بیشتر با «چگونه انجام دادن» سروکار دارد. ممکن است فردی نتواند به یاد بیاورد دیشب چه غذایی خورده، اما همان لحظه بتواند آشپزی کند. این جدایی نسبی نشان می‌دهد دانستن یک واقعیت و اجرای یک مهارت یکسان نیستند.

از نظر مغزی، حافظهٔ رویه‌ای با عقده‌های قاعده‌ای، به‌ویژه مدارهای مرتبط با \eng{striatum}، و نیز مدارهای حرکتی و یادگیری عادت ارتباط دارد. در مثال هوش مصنوعی، خودروی خودران می‌تواند تا حدی به حافظهٔ رویه‌ای تشبیه شود، چون بخشی از مهارت رانندگی را اجرا می‌کند. البته خودروی خودران هم دانش اعلانی، ادراک و تصمیم‌گیری دارد؛ پس این فقط یک تشبیه آموزشی است.

\coursefigure{0.50}{basal-ganglia-and-procedural-memory.png}{عقده‌های قاعده‌ای در یادگیری مهارت، عادت و حافظهٔ رویه‌ای نقش دارند؛ این نگاشت کمک می‌کند تفاوت «دانستن» و «انجام دادن» از نظر مغزی روشن‌تر شود.}

\begin{correctionbox}{دقت دربارهٔ \eng{Chain of Thought}}
\eng{Chain of Thought} را نباید مستقیماً معادل حافظهٔ رویه‌ای دانست. این روش بیشتر به آشکارسازی یا هدایت گام‌های استدلال کمک می‌کند و به بازیابی دانش و استدلال صریح نزدیک‌تر است. ممکن است در اجرای یک روش مرحله‌به‌مرحله اثر بگذارد، اما خودش مدل کامل حافظهٔ رویه‌ای انسان نیست.
\end{correctionbox}

\subsection{انتقال اطلاعات از حافظهٔ کوتاه‌مدت به بلندمدت}
اطلاعات معمولاً با تمرین، تکرار و استفاده وارد حافظهٔ بلندمدت می‌شوند. اما برخی شرایط می‌توانند این انتقال را بسیار سریع‌تر یا قوی‌تر کنند.

\subsubsection{احساس شدید}
اگر فرد هنگام دریافت یک اطلاعات یا تجربه، احساس شدیدی داشته باشد، احتمال ماندگاری آن بیشتر می‌شود. تصادف، تولد فرزند، یک خبر بسیار خوشحال‌کننده یا یک تجربهٔ ترسناک می‌تواند همراه با جزئیات زیادی در حافظه بماند: زمان، مکان، فصل، وضعیت هوا یا حتی صحنه‌های فرعی.

\subsubsection{تعجب و غیرمنتظره بودن}
چیزی که فرد را غافلگیر کند، راحت‌تر در حافظه می‌ماند. به همین دلیل در آموزش، تجربهٔ غیرمنتظره، آزمایش عملی یا مثال متفاوت می‌تواند یادگیری را تقویت کند. البته غافلگیری باید به فهم موضوع وصل شود؛ صرفاً عجیب بودن کافی نیست.

\subsubsection{تمایز و برجستگی}
اطلاعاتی که از زمینهٔ معمول جدا می‌شوند، شانس بیشتری برای ماندن دارند. اگر فردی ظاهر یا رفتاری بسیار متمایز داشته باشد، احتمال دارد نام یا تصویر او بهتر در ذهن بماند. در آموزش نیز برجسته‌سازی نکتهٔ مهم می‌تواند به ماندگاری کمک کند.

\subsubsection{علامت‌گذاری و ارتباط‌سازی}
روش‌هایی مانند علامت‌گذاری، ساخت مثال، تصویرسازی ذهنی، تجسم‌سازی و پیدا کردن ارتباط میان مفاهیم کمک می‌کنند اطلاعات جدید با ساختارهای قبلی ذهن پیوند بخورد. این پیوندها مسیر بازیابی را تقویت می‌کنند و احتمال انتقال به حافظهٔ بلندمدت را بالا می‌برند.

از منظر آموزشی، این پرسش خودش یک خط پژوهشی مهم است: چگونه می‌توان یادگیری را طوری طراحی کرد که اطلاعات مهم سریع‌تر، عمیق‌تر و پایدارتر وارد حافظهٔ بلندمدت شوند؟

\subsection{اختلالات و شواهد شناختی}
تقسیم‌بندی حافظه فقط زمانی ارزش نظری دارد که بتواند داده‌های انسانی، از جمله اختلالات، را توضیح دهد. اگر می‌گوییم حافظهٔ کوتاه‌مدت، بلندمدت، اعلانی و رویه‌ای از هم متمایزند، باید بتوانیم نشان دهیم در برخی شرایط یکی از این بخش‌ها بیشتر آسیب می‌بیند و دیگری بهتر باقی می‌ماند.

\subsubsection{استرس و بازیابی حافظه}
استرس می‌تواند بازیابی اطلاعات اعلانی را مختل کند. فرد ممکن است چیزی را بداند، اما در لحظهٔ فشار نتواند آن را توضیح دهد یا به یاد آورد. این حالت را می‌توان به اختلال در بازیابی \eng{fact}ها یا توضیح‌های صریح ربط داد.

اگر استرس روی خود عمل صحبت‌کردن، لکنت، کنترل تنفس یا اجرای مهارت اثر بگذارد، بخشی از مسئله می‌تواند با جنبه‌های رویه‌ای نیز مرتبط شود. بنابراین بهتر است استرس را فقط به یک نوع حافظه محدود نکنیم؛ اثر آن به نوع وظیفه و سطح درگیری سیستم بستگی دارد.

\subsubsection{پارکینسون و حافظهٔ رویه‌ای}
در بیماری پارکینسون، به‌ویژه در مراحل اولیه، اختلالات حرکتی و مهارت‌های رویه‌ای برجسته‌اند. لرزش، کندی حرکت، دشواری در شروع حرکت و تغییر در راه رفتن می‌تواند باعث شود فرد ابتدا تصور کند مشکل فقط از پا، کمر یا عضله است؛ در حالی که بخش مهمی از مسئله به مدارهای عصبی کنترل حرکت و یادگیری رویه‌ای مربوط است.

این مثال با نقش عقده‌های قاعده‌ای در حرکت، عادت و یادگیری رویه‌ای سازگار است. البته پارکینسون فقط یک اختلال سادهٔ حافظهٔ رویه‌ای نیست و در طول زمان می‌تواند جنبه‌های شناختی دیگری را نیز درگیر کند. بنابراین باید آن را به‌عنوان شاهدی برای تعامل سیستم‌های حرکتی، رویه‌ای و شناختی فهمید، نه به‌عنوان جداسازی کامل یک ماژول از بقیه.

\subsection{حافظه در مدل‌های زبانی}
در این چارچوب، مفاهیم حافظهٔ شناختی می‌توانند برای فکر کردن دربارهٔ مدل‌های زبانی مفید باشند. این نگاشت استعاری است و نباید آن را یکی‌گرفتن مستقیم مدل زبانی با مغز انسان دانست.

\subsubsection{پارامترها به‌عنوان حافظهٔ بلندمدت}
وزن‌ها و پارامترهای آموزش‌دیدهٔ مدل را می‌توان به حافظهٔ بلندمدت تشبیه کرد. دانشی که مدل در طول آموزش از داده‌ها جذب کرده، در این پارامترها توزیع شده است. وقتی مدل به پرسشی پاسخ می‌دهد، بخشی از این دانش آموخته‌شده را فعال و ترکیب می‌کند.

\subsubsection{\eng{Prompt} و \eng{Context Window} به‌عنوان حافظهٔ کاری}
\eng{Prompt} و \eng{context window} را می‌توان به حافظهٔ کاری یا بافر فعلی مدل تشبیه کرد. اطلاعاتی که در متن ورودی قرار دارد، در همان تعامل جاری در دسترس مدل است. اگر متن طولانی، پرنویز یا متناقض شود، مدل ممکن است در نگهداری و استفادهٔ درست از اطلاعات مهم دچار مشکل شود؛ مسئله‌ای که در فصل‌های بعد با عنوان بار شناختی و محدودیت حافظهٔ کاری بررسی می‌شود.

\subsubsection{\eng{RAG} به‌عنوان حافظهٔ بیرونی}
\eng{Retrieval-Augmented Generation} یا \eng{RAG} را می‌توان شبیه حافظهٔ بیرونی دانست. مدل به‌جای تکیهٔ کامل بر پارامترهای خودش، اطلاعات مرتبط را از منبعی بیرونی بازیابی می‌کند و سپس بر اساس آن پاسخ می‌دهد. این سازوکار به‌ویژه وقتی مهم است که دانش لازم بسیار خاص، تازه، حجیم یا نیازمند ارجاع دقیق باشد.

\section{نظریهٔ بار شناختی و حافظهٔ کاری}
تا اینجا حافظه به‌عنوان یکی از ماژول‌های اصلی معماری شناختی معرفی شد. در این فصل، تمرکز از «انواع حافظه» به محدودیت‌های عملی حافظهٔ کاری منتقل می‌شود. پرسش اصلی این است: وقتی اطلاعات زیاد، پراکنده یا مزاحم باشند، انسان و مدل چگونه افت می‌کنند؟

\subsection{\eng{Cognitive Load Theory}}
\term{نظریهٔ بار شناختی}{Cognitive Load Theory} بر این ایده تکیه دارد که حافظهٔ کاری انسان محدود است. یادگیری، حل مسئله و استدلال وقتی دشوار می‌شوند که مقدار اطلاعاتی که باید هم‌زمان نگه داشته و پردازش شود از ظرفیت مؤثر حافظهٔ کاری فراتر برود.

این نظریه در آموزش اهمیت ویژه‌ای دارد، چون طراحی بدِ محتوا می‌تواند بار اضافه ایجاد کند. اگر اسلاید، تمرین، منو، رابط کاربری یا توضیح آموزشی هم‌زمان اطلاعات زیادی را بدون ساختار مناسب ارائه کند، بخشی از توان شناختی خواننده صرف مدیریت آشفتگی می‌شود، نه فهم مفهوم اصلی.

در ادبیات کلاسیک آموزش، معمولاً میان چند نوع بار شناختی تمایز گذاشته می‌شود: دشواری ذاتی خود مطلب، بار اضافی ناشی از شیوهٔ ارائه، و باری که صرف ساختن طرحوارهٔ مفهومی مفید می‌شود. در این منبع، تمرکز اصلی روی ایدهٔ عمومی‌تر است: حافظهٔ کاری محدود است و طراحی محاسباتی یا آموزشی باید این محدودیت را جدی بگیرد.

\subsection{ظرفیت محدود حافظهٔ کاری}
حافظهٔ کاری فقط محل نگهداری موقت نیست؛ محل پردازش فعال اطلاعات فعلی نیز هست. وقتی شماره تلفنی را برای چند ثانیه در ذهن نگه می‌داریم تا آن را وارد گوشی کنیم، از حافظهٔ کاری استفاده کرده‌ایم. اگر هم‌زمان کسی با ما حرف بزند، عددها طولانی باشند یا الگوی آشنایی نداشته باشند، خطا بیشتر می‌شود.

برای بررسی ظرفیت حافظهٔ کاری می‌توان آزمایش‌های گوناگون طراحی کرد: نمایش سریع چند تصویر و پرسیدن تعداد یا نوع آیتم‌ها، پخش چند صدا و سنجش تشخیص آن‌ها، ارائهٔ رشته‌های عددی با طول‌های مختلف، یا دادن چند مزه و درخواست دسته‌بندی. وجه مشترک این آزمایش‌ها این است که با افزایش تعداد یا پیچیدگی آیتم‌ها، عملکرد افت می‌کند.

همین ایده برای مدل‌های زبانی نیز قابل ترجمه است. به جای نمایش تصویر یا رشتهٔ عدد به انسان، می‌توان \eng{prompt}هایی طراحی کرد که اطلاعات طولانی، پراکنده، چندمرحله‌ای یا همراه با نویز دارند و سپس دید آیا مدل همچنان می‌تواند اطلاعات مرتبط را نگه دارد و از آن‌ها استفاده کند یا نه.

\subsection{\texorpdfstring{عدد \eng{7 \(\pm\) 2}}{عدد 7 ± 2}}
مقالهٔ مشهور جورج میلر با عنوان \eng{The Magical Number Seven, Plus or Minus Two} در سال ۱۹۵۶ یکی از نقاط کلاسیک بحث دربارهٔ ظرفیت حافظهٔ کوتاه‌مدت است. برداشت رایج از این مقاله این است که انسان می‌تواند حدود \(7 \pm 2\) واحد را در حافظهٔ کوتاه‌مدت نگه دارد.

اما این عدد نباید به‌عنوان قانون ثابت همهٔ انسان‌ها و همهٔ موقعیت‌ها فهمیده شود. پژوهش‌های بعدی نشان داده‌اند ظرفیت مؤثر به نوع ماده، طول کلمات، آشنایی فرد با محتوا، روش اندازه‌گیری، سن، تمرین و راهبردهای سازمان‌دهی بستگی دارد. پیام اصلی برای این بحث، خود عدد هفت نیست؛ پیام اصلی این است که حافظهٔ کاری محدود است و واحد مهم در بسیاری از موقعیت‌ها \eng{chunk} است.

شکل زیر صفحهٔ آغازین مقالهٔ میلر را نشان می‌دهد تا روشن باشد این بحث از یک مشاهدهٔ تجربی دربارهٔ ظرفیت پردازش اطلاعات شروع شده و در طول زمان به نظریه‌های دقیق‌تر حافظهٔ کاری و بار شناختی رسیده است.

\coursefigure{0.74}{miller-1956-cognitive-capacity-paper.png}{صفحهٔ آغازین مقالهٔ کلاسیک جورج میلر دربارهٔ محدودیت ظرفیت پردازش اطلاعات؛ این مقاله نقطهٔ تاریخی مهمی برای بحث حافظهٔ کوتاه‌مدت و مفهوم چانک است.}

\begin{correctionbox}{دقت دربارهٔ عدد میلر}
عدد \(7 \pm 2\) یک بیان کلاسیک و آموزشی مفید است، اما ظرفیت حافظهٔ کاری را نمی‌توان همیشه با یک عدد ثابت توضیح داد. در بسیاری از بحث‌های جدید، ظرفیت مؤثر کمتر و وابسته‌تر به نوع اطلاعات و راهبردهای فرد در نظر گرفته می‌شود.
\end{correctionbox}

\subsection{\eng{Chunking}}
\eng{Chunking} راهی برای کاهش بار حافظهٔ کاری است. انسان به جای نگه داشتن تعداد زیادی جزء جداگانه، آن‌ها را به واحدهای معنادار بزرگ‌تر تبدیل می‌کند. به این ترتیب، حافظهٔ کاری به جای چندین جزء پراکنده، چند واحد سازمان‌یافته را نگه می‌دارد.

\subsubsection{تعریف \eng{Chunk}}
\term{چانک}{Chunk} یعنی گروهی از اطلاعات کوچک‌تر که برای فرد به‌صورت یک واحد معنادار در می‌آیند. برای مثال:
\[
\text{\lr{1 3 6 9}}
\]
اگر این چهار رقم جداگانه دیده شوند، چهار آیتم‌اند. اما اگر فرد آن‌ها را سال \eng{1369} یا بخشی از یک الگوی آشنا بداند، کل رشته می‌تواند یک \eng{chunk} شود.

بنابراین اندازهٔ واقعی حافظهٔ کاری فقط به تعداد خام نشانه‌ها بستگی ندارد؛ به این بستگی دارد که ذهن چگونه آن نشانه‌ها را سازمان‌دهی می‌کند.

\subsubsection{نقش دانش قبلی در \eng{Chunking}}
\eng{Chunking} به حافظهٔ بلندمدت وابسته است. ما برای فشرده‌سازی اطلاعات جدید از دانش قبلی، الگوهای آشنا، معنا، تصویر ذهنی و تداعی‌های قبلی استفاده می‌کنیم. کسی که یک رشتهٔ عددی را به تاریخ تولد، شمارهٔ ماشین، وزن ورزشی یا الگوی ریاضی ربط می‌دهد، آن را بهتر از کسی نگه می‌دارد که فقط رقم‌ها را جداگانه حفظ می‌کند.

به همین دلیل تخصص می‌تواند ظرفیت مؤثر حافظهٔ کاری را تغییر دهد. متخصص شطرنج ممکن است چیدمان صفحه را به چند الگوی معنادار تبدیل کند، در حالی که فرد مبتدی همان صفحه را مجموعه‌ای از مهره‌های پراکنده می‌بیند.

\subsubsection{تفاوت‌های فردی}
تفاوت افراد در حافظهٔ کاری فقط به ظرفیت خام محدود نمی‌شود. افراد در نوع و کیفیت \eng{chunking} نیز تفاوت دارند. یک نفر چانک تصویری می‌سازد، دیگری چانک زبانی، دیگری از داستان یا تداعی استفاده می‌کند و فردی دیگر الگوی عددی یا فضایی پیدا می‌کند.

بنابراین نظریهٔ خوب باید بتواند توضیح دهد چرا دو نفر با دادهٔ ظاهراً یکسان عملکرد متفاوت دارند. تفاوت می‌تواند از اندازهٔ چانک‌ها، نوع چانک‌ها، توانایی اتصال اطلاعات جدید به حافظهٔ بلندمدت، یا راهبردهای سازمان‌دهی ذهنی ناشی شود.

\subsection{کاربردهای \eng{Chunking}}
\eng{Chunking} فقط یک مفهوم آزمایشگاهی نیست. در طراحی آموزشی، طراحی اسلاید، طراحی رابط کاربری، طراحی منوها، بازی‌ها، شماره تلفن‌ها و سازمان‌دهی اطلاعات پیچیده از همین اصل استفاده می‌شود. اگر کاربر یا خواننده با تعداد زیادی آیتم بی‌ساختار روبه‌رو شود، حافظهٔ کاری او سریع‌تر اشباع می‌شود. اما اگر همان اطلاعات در گروه‌های معنادار، مرحله‌های کوچک‌تر یا ساختارهای قابل پیش‌بینی ارائه شود، پردازش ساده‌تر می‌شود.

\subsection{\eng{Multi-Hop Reasoning} به‌عنوان آزمون حافظهٔ کاری}
\term{استدلال چندمرحله‌ای}{Multi-Hop Reasoning} وظیفه‌ای است که پاسخ نهایی در آن از یک گام ساده به دست نمی‌آید. مدل یا انسان باید چند قطعه اطلاعات را نگه دارد، رابطهٔ میان آن‌ها را بفهمد، نتیجهٔ میانی بسازد و از نتیجهٔ میانی برای گام بعدی استفاده کند.

در یک مثال ساده:
\begin{itemize}
  \item انسان‌ها فانی‌اند.
  \item سقراط انسان است.
  \item پس سقراط فانی است.
  \item اگر فانی‌ها معمولاً تا سن خاصی عمر کنند، نتیجهٔ بعدی نیز به این زنجیره وابسته می‌شود.
\end{itemize}

یا در پرسش خانوادگی «پدرِ همسرِ برادرِ مادر من چه کسی است؟» باید چند رابطهٔ پشت سر هم پردازش شود. هر گام باید تا گام بعدی در حافظهٔ کاری فعال بماند. همین ویژگی، \eng{multi-hop reasoning} را برای سنجش بار شناختی در انسان و مدل مناسب می‌کند.

\subsection{محدودیت حافظهٔ کاری در مدل‌های زبانی}
در مدل‌های زبانی، معادل آزمایش حافظهٔ کاری را می‌توان با طراحی \eng{prompt}های دشوار انجام داد. هدف این است که نشان دهیم:
\begin{enumerate}
  \item مدل نیز در استفاده از زمینهٔ فعلی محدودیت دارد.
  \item با افزایش بار اطلاعاتی، عملکرد افت می‌کند.
  \item نویز، پراکندگی شواهد و تغییر موضوع می‌توانند بازیابی اطلاعات مرتبط را مختل کنند.
  \item سپس می‌توان با الهام از انسان، راهکارهایی مانند \eng{chunking}، خلاصه‌سازی، تقسیم مسئله یا مدیریت حافظه را بررسی کرد.
\end{enumerate}

در سال‌های اخیر، برخی کارها همین ایده را با عنوان بار شناختی محاسباتی در مدل‌های زبانی بررسی کرده‌اند. یکی از مسیرهای مطرح، سنجش مدل‌ها در \eng{multi-hop reasoning} همراه با دست‌کاری مقدار نویز یا ترتیب وظایف است.

\subsubsection{\eng{Context Saturation}}
\eng{Context Saturation} زمانی رخ می‌دهد که اطلاعات مهم در میان حجم زیادی از اطلاعات نامرتبط یا کم‌ارتباط پخش شود. در این حالت، مسئله فقط طول \eng{context window} نیست؛ مسئله این است که مدل باید از میان مقدار زیادی نویز، شواهد اصلی را پیدا کند.

ساختار آزمایش می‌تواند چنین باشد: ابتدا یک مسئلهٔ چندمرحله‌ای داده می‌شود، سپس مقدار زیادی متن نامرتبط یا اطلاعات مزاحم میان شواهد مرتبط اضافه می‌شود، و در نهایت عملکرد مدل اندازه‌گیری می‌شود. اگر با افزایش نویز، نرخ موفقیت کاهش یابد، می‌توان گفت زمینهٔ مدل از نظر شناختی اشباع شده است.

\subsubsection{\eng{Attentional Residue}}
\eng{Attentional Residue} به باقی‌ماندن اثر وظیفه یا موضوع قبلی روی وظیفهٔ فعلی اشاره دارد. در تجربهٔ عملی با مدل‌های زبانی، گاهی مدل سؤال جدید را با زمینهٔ قبلی قاطی می‌کند، یا هنوز در چارچوب مسئلهٔ قبلی پاسخ می‌دهد. در چنین حالتی، پاک‌کردن مکالمه، بازنویسی \eng{prompt} یا جداکردن وظایف می‌تواند کمک کند.

در طراحی آزمایش، می‌توان ابتدا چند بخش نامرتبط یا یک وظیفهٔ مزاحم به مدل داد، سپس مسئلهٔ اصلی را مطرح کرد و دید آیا مدل می‌تواند توجه خود را از اثر باقی‌ماندهٔ قبلی جدا کند یا نه.

\subsubsection{تفاوت دو نوع فشار}
هر دو پدیده با نویز و مزاحمت سروکار دارند، اما چیدمان آن‌ها متفاوت است:
\begin{itemize}
  \item در \eng{Context Saturation}، اطلاعات مرتبط در میان نویز پخش می‌شود.
  \item در \eng{Attentional Residue}، وظیفه یا نویز قبلی روی وظیفهٔ بعدی سایه می‌اندازد.
\end{itemize}
در هر دو حالت، افت عملکرد مدل نشانه‌ای از محدودیت در مدیریت زمینه، توجه و حافظهٔ کاری است.

\subsection{طراحی آزمایش برای سنجش بار شناختی مدل}
برای نشان دادن محدودیت حافظهٔ کاری در مدل زبانی، فقط طولانی کردن \eng{prompt} کافی نیست. باید طراحی آزمایش طوری باشد که عامل‌های مزاحم قابل کنترل و قابل تفسیر باشند. چند روش مطرح‌شده در این بحث:
\begin{itemize}
  \item افزایش طول \eng{prompt}
  \item اضافه کردن اطلاعات زیاد اما نامرتبط
  \item پخش کردن شواهد مرتبط در بخش‌های دور از هم
  \item اضافه کردن نویز میان اطلاعات مهم
  \item ارجاع دادن به بخش‌های قبلی مکالمه
  \item ایجاد تداخل میان موضوع قبلی و موضوع فعلی
  \item دادن یک کد یا دستور و سپس تغییر خواسته در ادامهٔ \eng{prompt}
\end{itemize}

در هر طراحی، باید روشن باشد چه چیزی اندازه‌گیری می‌شود: بازیابی اطلاعات مرتبط، حفظ نتیجهٔ میانی، مقاومت در برابر نویز، یا توانایی تغییر توجه از وظیفهٔ قبلی به وظیفهٔ جدید.

\subsection{تکامل آزمون‌های هوش مصنوعی}
این بحث در ادامهٔ تکامل ارزیابی هوش مصنوعی قرار می‌گیرد. در ابتدا آزمون‌هایی مانند \eng{Turing Test} برجسته بودند. بعدتر در یادگیری ماشین، معیارهایی مانند \eng{accuracy}، \eng{precision} و \eng{recall} روی \eng{test set} اهمیت پیدا کردند. با پیشرفت مدل‌ها، مسئلهٔ \term{تعمیم خارج از توزیع}{Out-of-Distribution Generalization} مطرح شد.

اکنون یکی از مسیرهای مهم، آزمون‌های شناختی است. به جای اینکه فقط بپرسیم مدل روی یک بنچمارک چند درصد امتیاز می‌گیرد، می‌پرسیم: آیا مدل مانند انسان محدودیت حافظهٔ کاری دارد؟ آیا خطاهای مدل با خطاهای انسان همبستگی دارد؟ آیا در مهارت‌هایی مانند استدلال، حافظه، برنامه‌ریزی و انتزاع ضعف نظام‌مند دارد؟ آیا می‌توان با الهام از نظریه‌های شناختی این ضعف را کاهش داد؟

\subsection{راهکارهای شناختی}
وقتی محدودیت حافظهٔ کاری شناسایی شد، مرحلهٔ بعد طراحی راهکار است. در انسان، راهکارهایی مانند \eng{chunking}، تقسیم مسئله، ساختن نشانه‌های بازیابی، خلاصه‌سازی و استفاده از دانش قبلی به کاهش بار شناختی کمک می‌کنند. در مدل‌های زبانی نیز می‌توان از ایده‌های مشابه استفاده کرد: تقسیم مسئله به زیرمسئله، خلاصه کردن شواهد، ساختن حافظهٔ بیرونی، یا طراحی \eng{prompt}هایی که اطلاعات را مرحله‌به‌مرحله سازمان‌دهی می‌کنند.

\subsubsection{\eng{Divide and Conquer}}
\eng{Divide and Conquer} یعنی مسئلهٔ بزرگ را به چند مسئلهٔ کوچک‌تر تقسیم کنیم. این ایده وقتی مفید است که نتیجهٔ هر بخش بتواند به‌صورت خلاصه و قابل استفاده برای گام بعدی نگه داشته شود. اگر مدل مجبور نباشد همهٔ جزئیات خام را هم‌زمان در \eng{context} نگه دارد، فشار حافظهٔ کاری کمتر می‌شود.

از نظر الگوریتمی، این ایده را می‌توان با رابطهٔ کلی زیر نمایش داد:
\[
T(n)=aT(n/b)+f(n)
\]
که در آن \(a\) تعداد زیرمسئله‌ها، \(n/b\) اندازهٔ هر زیرمسئله، و \(f(n)\) هزینهٔ شکستن مسئله و ترکیب پاسخ‌هاست. پیام آموزشی این فرمول این است که تقسیم درست مسئله می‌تواند هزینه را به‌شدت کم کند؛ برای مثال، در بعضی مسائل به جای هزینه‌های مرتبهٔ \(O(n^2)\) می‌توان به ساختارهایی نزدیک‌تر به \(O(n\log n)\) رسید. البته این مقایسه قانون عمومی همهٔ مسائل نیست و به ساختار دقیق مسئله و الگوریتم بستگی دارد.

یک مثال ساده، خانه‌تکانی است. اگر کل خانه یک‌باره موضوع کار باشد، همهٔ جزئیات هم‌زمان روی حافظهٔ کاری فشار می‌آورند. اما اگر خانه به اتاق‌ها، اتاق به کمد و کشو، و هر کشو به چند دستهٔ کوچک‌تر تقسیم شود، در هر لحظه فقط یک بخش فعال است. پس از تمام شدن هر بخش، لازم نیست همهٔ حرکت‌های کوچک در حافظه بماند؛ کافی است نتیجهٔ آن بخش به‌عنوان یک وضعیت خلاصه نگه داشته شود.

\subsubsection{\eng{Sub-goal} و مدیریت سلسله‌مراتبی حافظه}
در وظایف طولانی، هر \eng{sub-goal} می‌تواند نقش یک چانک سطح بالاتر را داشته باشد. مدل ابتدا یک زیرهدف را حل می‌کند، خروجی مهم آن را نگه می‌دارد، سپس به زیرهدف بعدی می‌رود. این ایده در \eng{agent}های بلندمدت مهم است، چون مدل باید بداند چه چیزی از مرحلهٔ قبلی را نگه دارد و چه جزئیاتی را می‌تواند کنار بگذارد.

در مقالهٔ \eng{HiAgent} این ایده به شکل صریح به مدیریت حافظهٔ کاری \eng{agent}های زبانی وصل می‌شود. در راهبرد ساده، همهٔ جفت‌های «عمل--مشاهده» در طول انجام کار وارد زمینهٔ مدل می‌شوند. این کار اطلاعات را کامل نگه می‌دارد، اما در وظایف بلندمدت زمینه را پر از جزئیات تکراری و کم‌اهمیت می‌کند. راه‌حل پیشنهادی این است که مدل ابتدا زیرهدف بسازد، عمل‌های مربوط به همان زیرهدف را دنبال کند، و وقتی زیرهدف کامل شد، جزئیات اجرایی آن را با یک مشاهدهٔ خلاصه جایگزین کند.

به این ترتیب، حافظهٔ کاری به‌جای یک فهرست بلند از همهٔ حرکت‌ها، چیزی شبیه این دارد:
\[
\left(g_1,s_1\right),\left(g_2,s_2\right),\ldots,\left(g_k,\text{جزئیات زیرهدف فعلی}\right)
\]
که در آن \(g_i\) زیرهدف و \(s_i\) خلاصهٔ وضعیت پس از انجام آن زیرهدف است. اگر بعداً جزئیات یک زیرهدف گذشته لازم شود، می‌توان آن بخش را دوباره بازیابی کرد؛ اما حالت پیش‌فرض این نیست که همهٔ جزئیات همیشه در حافظهٔ فعال بمانند.

\subsubsection{مثال \eng{Blocks World}}
در \eng{Blocks World} چند بلوک روی میز یا روی یکدیگر قرار دارند و هدف این است که به وضعیت نهایی مشخصی برسیم؛ مثلاً بلوک آبی روی بلوک نارنجی قرار گیرد. برای رسیدن به این هدف، مسئله را می‌توان به چند زیرهدف شکست:
\begin{enumerate}
  \item بلوک آبی \eng{clear} شود؛ یعنی چیزی روی آن نباشد.
  \item بلوک نارنجی \eng{clear} شود.
  \item بلوک آبی روی بلوک نارنجی گذاشته شود.
\end{enumerate}

نکتهٔ شناختی این نیست که زیرهدف در برنامه‌ریزی رباتیک چیز تازه‌ای است؛ زیرهدف‌ها از قبل در \eng{planning} وجود داشته‌اند. نکته این است که در اینجا زیرهدف به‌عنوان واحد حافظه نیز تفسیر می‌شود. پس از کامل شدن زیرهدف «بلوک نارنجی clear است»، لازم نیست همهٔ حرکت‌هایی که به این وضعیت منجر شده‌اند در حافظهٔ فعال بمانند. نتیجهٔ فشردهٔ زیرهدف برای ادامهٔ حل مسئله کافی است.

\subsubsection{مثال \eng{Gripper}}
در \eng{Gripper}، ربات باید اشیایی را بین اتاق‌ها جابه‌جا کند. در نتایج گزارش‌شدهٔ \eng{HiAgent}، چانک‌سازی سلسله‌مراتبی در بعضی وظایف مانند \eng{Blocksworld} سود آشکارتری داشت، اما در \eng{Gripper} بهبود موفقیت ایجاد نکرد و حتی شاخص پیشرفت کمی افت کرد، هرچند مصرف زمینه کمتر شد. تفسیر آموزشی این است که \eng{chunking} وقتی بیشترین فایده را دارد که مسئله واقعاً ساختار سلسله‌مراتبی و وابستگی میان زیرهدف‌ها داشته باشد. اگر کارها بسیار مستقل و خرد باشند، تقسیم اضافه ممکن است فقط هزینهٔ مدیریتی بسازد.

\subsubsection{محدودیت‌های \eng{Chunking}}
\eng{Chunking} همیشه راه‌حل کامل نیست. اگر مسئله ذاتاً ساختار تقسیم‌پذیر نداشته باشد، یا اگر تقسیم‌بندی غلط انجام شود، چانک‌ها می‌توانند اطلاعات مهم را پنهان کنند. همچنین در مدل‌های زبانی، ساختن چانک معمولاً به طراحی بیرونی، \eng{prompting} یا ماژول حافظه نیاز دارد؛ مدل همیشه خودش نمی‌داند چه چیزی را باید خلاصه کند و چه چیزی را باید نگه دارد.

تفاوت مهم انسان و مدل نیز همین‌جاست. انسان‌ها در بسیاری از موقعیت‌ها خودشان راهبرد می‌سازند: موضوع را دسته‌بندی می‌کنند، چیزهای کم‌اهمیت را کنار می‌گذارند، و تشخیص می‌دهند چه چیزی باید در حافظه بماند. در بسیاری از \eng{agent}های فعلی، ما هنوز باید صریحاً بگوییم زیرهدف چیست، چه چیزی چانک شود، چه چیزی فراموش شود و چه چیزی بازیابی شود. بنابراین راهبرد شناختی به مدل الهام می‌دهد، اما پیاده‌سازی آن هنوز به طراحی محاسباتی دقیق نیاز دارد.

\section{ترنسفورمر، توجه و حافظهٔ کاری}

\subsection{حافظه در \eng{Transformer}}
برای پیوند دادن بحث حافظهٔ کاری با معماری‌های امروزی، باید ببینیم در \eng{Transformer} چه چیزی را می‌توان به‌صورت تقریبی معادل حافظهٔ بلندمدت و کوتاه‌مدت دانست. این تشبیه کامل و زیستی نیست، اما برای تحلیل مدل مفید است.

حافظهٔ بلندمدت مدل در وزن‌ها و پارامترهای آن قرار دارد. دانشی که مدل در زمان آموزش از متن، گفتار، تصویر یا داده‌های دیگر آموخته، در همین پارامترها کدگذاری شده است. برای مثال، اینکه مدل واژهٔ \eng{data} را به چه حوزه‌هایی نزدیک می‌داند، یا چه رابطه‌هایی میان کلمات، ساختارهای نحوی و مفاهیم یاد گرفته است، محصول آموزش و پارامترهاست.

در مقابل، حافظهٔ کوتاه‌مدت یا حافظهٔ کاری به پردازش ورودی فعلی مربوط است: \eng{prompt}، جمله، متن، مکالمه یا دنباله‌ای که در زمان \eng{inference} پیش روی مدل قرار دارد. بخشی از معماری که بیشترین ارتباط را با این پردازش دارد، لایهٔ \eng{attention} است؛ زیرا مدل در آنجا تصمیم می‌گیرد هر توکن، برای ساختن بازنمایی فعلی خود، باید به کدام توکن‌های دیگر توجه کند.

\begin{correctionbox}{تشبیه، نه همسانی کامل}
وقتی می‌گوییم پارامترها شبیه حافظهٔ بلندمدت و \eng{attention} شبیه حافظهٔ کاری است، منظور یک نگاشت دقیق عصبی نیست. مدل زبانی مغز ندارد و حافظهٔ انسانی نیز صرفاً ماتریس وزن نیست. این تشبیه فقط کمک می‌کند نقش «دانش آموخته‌شده» و «پردازش زمینهٔ فعلی» را جدا کنیم.
\end{correctionbox}

\subsection{\eng{Contextual Embedding}}
مدل در ابتدا برای هر \eng{token} یک بازنمایی اولیه دارد. این بازنمایی از دانش آموخته‌شدهٔ مدل و جدول‌های \eng{embedding} یا لایه‌های ورودی می‌آید. اما هدف \eng{Transformer} این نیست که یک معنای ثابت برای هر کلمه نگه دارد؛ هدف این است که بازنمایی کلمه با توجه به زمینهٔ فعلی تغییر کند. به این بازنمایی وابسته به زمینه \term{بردار زمینه‌مند}{Contextual Embedding} می‌گوییم.

برای مثال، واژهٔ «شیر» در فارسی می‌تواند به حیوان، نوشیدنی یا وسیلهٔ باز و بسته کردن آب اشاره کند. اگر مدل فقط یک بردار ثابت برای «شیر» داشته باشد، نمی‌تواند این تفاوت‌ها را به‌خوبی نشان دهد. اما پس از لایه‌های \eng{self-attention}، بازنمایی «شیر» در «شیر در جنگل زندگی می‌کند» باید با بازنمایی آن در «شیر را در لیوان ریخت» متفاوت شود.

همین نکته برای واژهٔ \eng{data} نیز برقرار است. بسته به زمینه، \eng{data} ممکن است به پایگاه داده، یادگیری ماشین، مصورسازی، آمار، حریم خصوصی یا سیگنال آزمایشگاهی نزدیک شود. \eng{Contextual embedding} یعنی مدل معنای عملی توکن را در همان متن خاص بسازد، نه اینکه فقط به معنای میانگین و ثابت آن تکیه کند.

\subsubsection{حل ارجاع}
\term{حل ارجاع}{Coreference Resolution} یعنی مدل بفهمد چند عبارت مختلف به یک موجودیت واحد اشاره می‌کنند. در جملهٔ «علی رفت مدرسه. او کتابش را برداشت»، مدل باید تشخیص دهد «او» و «ـش» به علی برمی‌گردند. در سطح بازنمایی، انتظار داریم بردارهای مربوط به این ارجاع‌ها به اطلاعاتی دربارهٔ «علی» وصل شوند، نه اینکه ضمیرها جدا و بی‌ارتباط پردازش شوند.

این کار نمونهٔ روشنی از حافظهٔ کاری در پردازش زبان است. مدل باید موجودیتی را که در جملهٔ قبلی آمده نگه دارد، ضمیر را در جملهٔ بعدی ببیند، و میان آن‌ها پیوند بسازد.

\subsubsection{رفع ابهام معنای واژه}
\term{رفع ابهام معنای واژه}{Word Sense Disambiguation} یعنی تشخیص اینکه یک واژهٔ چندمعنا در متن فعلی کدام معنا را دارد. برای «شیر»، نشانه‌هایی مانند «جنگل»، «یال»، «لیوان»، «کلسیم»، «آب» یا «باز کردن» معنا را تعیین می‌کنند. مدل باید این نشانه‌ها را در زمینه پیدا کند و بازنمایی واژه را به معنای درست نزدیک کند.

پس \eng{contextual embedding} فقط یک اصطلاح فنی در \eng{NLP} نیست؛ همان نقطه‌ای است که مدل از حافظهٔ بلندمدت خود، یعنی دانش آماری و معنایی آموخته‌شده، برای پردازش وضعیت فعلی استفاده می‌کند.

\subsection{\eng{Query}، \eng{Key} و \eng{Value}}
در \eng{self-attention}، بازنمایی هر توکن با چند ماتریس وزنی به سه بردار تبدیل می‌شود:
\begin{itemize}
  \item \term{پرس‌وجو}{Query}: این توکن دنبال چه اطلاعاتی است؟
  \item \term{کلید}{Key}: این توکن چه نشانه‌ای برای یافتن خود ارائه می‌کند؟
  \item \term{مقدار}{Value}: اگر به این توکن توجه شود، چه اطلاعاتی منتقل می‌شود؟
\end{itemize}

اگر بردار اولیهٔ یک توکن مثلاً ابعاد \(1 \times 768\) داشته باشد، مدل با ضرب آن در ماتریس‌های آموخته‌شده، بردارهای \(Q\)، \(K\) و \(V\) می‌سازد. سپس \eng{query} یک توکن با \eng{key} توکن‌های دیگر مقایسه می‌شود. هرچه شباهت یا سازگاری آن‌ها بیشتر باشد، وزن توجه بیشتر می‌شود. در شکل ساده، وزن‌ها با ضرب داخلی و سپس \eng{softmax} ساخته می‌شوند:
\[
\operatorname{Attention}(Q,K,V)=\operatorname{softmax}\left(\frac{QK^{T}}{\sqrt{d_k}}\right)V
\]
که \(d_k\) بعد بردارهای کلید است. ضرب نهایی در \(V\) باعث می‌شود بازنمایی جدید هر توکن از ترکیب اطلاعات توکن‌هایی ساخته شود که برای آن مهم‌تر بوده‌اند.

در شکل زیر، همین تبدیل خطی به‌صورت تصویری دیده می‌شود: embeddingهای ورودی با وزن‌های آموخته‌شده ترکیب می‌شوند تا نمایش‌های \(Q\)، \(K\) و \(V\) ساخته شوند. این تصویر کمک می‌کند توجه را نه به‌عنوان یک مفهوم مبهم، بلکه به‌عنوان چند ضرب ماتریسی روی بازنمایی‌های فعلی ببینیم.

\coursefigure{0.88}{query-key-value-projections.png}{ساخت بردارهای \eng{Query}، \eng{Key} و \eng{Value} از embeddingهای ورودی با ضرب در وزن‌های آموخته‌شده و اضافه شدن بایاس.}

\subsection{\eng{Self-Attention} از دید حافظهٔ کاری}
\eng{Self-attention} را می‌توان به‌صورت آموزشی بخشی از سازوکار حافظهٔ کاری مدل دانست، چون روی ورودی فعلی کار می‌کند. پیش از \eng{attention}، مدل بازنمایی‌هایی دارد که عمدتاً از حافظهٔ بلندمدت و جایگاه توکن‌ها آمده‌اند. پس از \eng{attention}، هر توکن بازنمایی تازه‌ای دارد که به متن فعلی وابسته است.

برای مثال، در جمله‌ای که چند بار از «شیر» با معناهای متفاوت استفاده شده، \eng{attention} باید نشانه‌های اطراف هر نمونه را جمع کند تا بازنمایی مناسب بسازد. در مسئلهٔ چندمرحله‌ای، \eng{attention} باید میان قطعه‌های دور از هم پیوند برقرار کند. بنابراین اگر مدل در بازیابی شواهد مرتبط یا نادیده گرفتن اطلاعات نامرتبط ضعف داشته باشد، این ضعف را می‌توان تا حدی در نحوهٔ مدیریت توجه و زمینه جست‌وجو کرد.

البته در \eng{decoder-only transformer}های خودرگرسیو، هر توکن فقط به توکن‌های قبلی توجه می‌کند، چون تولید متن نباید به آینده دسترسی داشته باشد. در \eng{encoder-only}ها، مانند مدل‌های خانوادهٔ \eng{BERT}، توجه معمولاً دوسویه است و هر توکن می‌تواند از توکن‌های قبل و بعد استفاده کند. پس نوع ترنسفورمر بر شکل حافظهٔ کاری محاسباتی اثر می‌گذارد.

\coursefigure{0.76}{attention-weighted-values.png}{ضرب ماتریس توجه در بردارهای \eng{Value}: وزن‌های توجه مشخص می‌کنند هر توکن چه مقدار از اطلاعات توکن‌های دیگر را در بازنمایی خروجی خود دریافت کند.}

\subsection{آیا افزایش طول \eng{Context} کافی است؟}
یک پاسخ ساده به محدودیت حافظهٔ کاری مدل این است که پنجرهٔ زمینه را بزرگ‌تر کنیم. این کار در بسیاری از کاربردها مفید است، اما به‌تنهایی مسئله را حل نمی‌کند. انسان نیز اگر حافظهٔ خام بیشتری داشته باشد، بدون راهبرد سازمان‌دهی الزاماً بهتر فکر نمی‌کند. مسئلهٔ اصلی فقط «جا داشتن اطلاعات» نیست؛ مسئله این است که مدل بتواند اطلاعات مرتبط را بازیابی کند، اطلاعات نامرتبط را کنار بگذارد و مسیر استدلال را حفظ کند.

اگر \eng{prompt} طولانی باشد اما پر از نویز، تکرار یا موضوع‌های نامرتبط شود، افزایش طول زمینه حتی می‌تواند بار پردازشی بیشتری ایجاد کند. بنابراین هدف بهتر، افزایش بی‌قید حافظه نیست؛ هدف \term{مدیریت توجه و بازیابی}{Attention and Retrieval Management} است.

\subsection{کاهش یا حذف توجه به اطلاعات کم‌اهمیت}
یک ایدهٔ معماری این است که ارتباط‌هایی با \eng{attention score} پایین حذف یا صفر شوند تا مدل روی موارد مهم‌تر تمرکز کند. این ایده از نظر شهودی جذاب است، اما ساده نیست.

اول، وزن توجه نسبی است. ممکن است یک توکن نسبت به توکن دوم وزن کمی داشته باشد، اما نسبت به توکن سوم وزن بالایی داشته باشد. بنابراین «کم بودن» یک وزن در یک رابطه به‌تنهایی نشان نمی‌دهد آن توکن یا مسیر اطلاعاتی بی‌اهمیت است.

دوم، وابستگی‌ها می‌توانند زنجیره‌ای باشند. اگر \(A\) به \(B\) مربوط باشد و \(B\) به \(C\)، ممکن است \(A\) از راه \(B\) به \(C\) وابسته شود، حتی اگر توجه مستقیم \(A\) به \(C\) کوچک باشد. حذف خام وزن‌های کوچک می‌تواند این مسیرهای غیرمستقیم را خراب کند.

پس کاهش توجه باید با طراحی دقیق انجام شود: با پنجره‌های محلی، الگوهای پراکنده، بازیابی هدفمند، خلاصه‌سازی، یا لایه‌هایی که اجازه دهند اطلاعات مهم در سطح بالاتر دوباره ترکیب شود.

\subsection{\eng{Local Attention}}
\eng{Local Attention} یعنی هر توکن به جای مقایسه با همهٔ توکن‌های دنباله، فقط به یک پنجرهٔ محلی از توکن‌های نزدیک خود توجه کند. این کار از نظر محاسباتی مهم است، چون \eng{self-attention} کامل نسبت به طول دنباله هزینهٔ درجهٔ دوم دارد؛ اگر طول دنباله \(n\) باشد، مقایسهٔ همه با همه تقریباً از مرتبهٔ \(O(n^2)\) است.

با محلی کردن توجه، هزینه و مصرف حافظه کاهش می‌یابد و مدل مجبور نیست همهٔ ارتباط‌های ممکن را هم‌زمان بررسی کند. از نظر مفهومی، این ایده به \eng{divide and conquer} شباهت دارد: ورودی به بخش‌های کوچک‌تر تقسیم می‌شود، در هر بخش پردازش محلی انجام می‌گیرد، و سپس لایه‌های بعدی یا سازوکارهای مکمل باید اطلاعات بخش‌ها را با هم ترکیب کنند.

اما \eng{local attention} خطر مهمی دارد: وابستگی‌های دوربرد ممکن است از دست بروند. اگر پاسخ یک سؤال به جمله‌ای در ابتدای متن و جمله‌ای در انتهای متن وابسته باشد، پنجرهٔ محلی ساده ممکن است این دو را به هم وصل نکند. به همین دلیل، مدل‌های عملی معمولاً فقط به پنجرهٔ محلی خام تکیه نمی‌کنند؛ ممکن است از پنجره‌های لغزان، توکن‌های جهانی، الگوهای پراکنده یا لایه‌های سلسله‌مراتبی برای عبور اطلاعات بین بخش‌ها استفاده کنند.

یک مثال مناسب این است که به جای نگاه هم‌زمان به کل محیط، ابتدا با یک بخش تعامل شود و سپس توجه به بخش‌های دیگر منتقل گردد. کل محیط حذف نشده است؛ فقط توجه به‌صورت مرحله‌ای و محلی اعمال می‌شود. در مدل نیز \eng{local attention} به معنی حذف متن نیست، بلکه به معنی محدود کردن میدان توجه در هر گام و نیاز به ترکیب بعدی است.

\subsubsection{ترکیب پنجره‌های محلی}
اگر در هر پنجرهٔ محلی عناصر مهم پیدا شوند، هنوز باید مشخص شود چگونه این عناصر با عناصر مهم پنجره‌های دیگر ترکیب می‌شوند. یک راه ساده این است که از بازنمایی‌های خلاصه، میانگین‌ها یا توکن‌های تجمیعی استفاده شود. راه دقیق‌تر این است که لایهٔ بالاتری ساخته شود که بازنمایی‌های محلی را دوباره با هم مقایسه کند. بدون چنین مرحله‌ای، مدل ممکن است فقط چند نگاه محلی جداگانه داشته باشد و بازنمایی منسجم کل متن را نسازد.

\subsection{\eng{Multi-Head Attention} و \eng{Mixture of Experts}}
\eng{Multi-head attention} و \eng{Mixture of Experts} را نباید مستقیماً معادل راهکار حافظهٔ کاری دانست، اما هر دو می‌توانند به استفادهٔ بهتر از ظرفیت مدل کمک کنند.

در \eng{multi-head attention}، مدل چند مجموعهٔ جدا از \(Q\)، \(K\) و \(V\) دارد. هر \eng{head} می‌تواند جنبه‌ای متفاوت از ورودی را برجسته کند: یک سر ممکن است رابطه‌های نحوی را دنبال کند، سر دیگر رابطه‌های معنایی، سر دیگر ارجاع‌های موجودیت‌ها، و سر دیگر نشانه‌های موقعیت یا ساختار جمله را. سپس خروجی سرها با هم ترکیب می‌شود. این کار شبیه نگاه کردن از چند زاویه است، نه الزاماً بزرگ‌تر کردن حافظهٔ کاری.

\eng{Mixture of Experts} نیز معمولاً به این معناست که بخش‌هایی از مدل یا «متخصص‌ها» برای نمونه‌ها یا توکن‌های مختلف فعال می‌شوند. هدف اصلی آن افزایش ظرفیت مؤثر یا تخصصی کردن پردازش است، بدون اینکه همیشه همهٔ پارامترها فعال شوند. این ایده می‌تواند غیرمستقیم به مدیریت اطلاعات کمک کند، اما مسئلهٔ اصلی حافظهٔ کاری، یعنی نگه‌داری و بازیابی هدفمند اطلاعات زمینهٔ فعلی، همچنان به طراحی توجه، حافظه، خلاصه‌سازی و راهبردهای استدلال وابسته است.

\subsection{جمع‌بندی}
در این فصل دو مسیر برای کاهش محدودیت حافظهٔ کاری مدل‌ها از هم جدا شد. مسیر اول شناختی و الگوریتمی بود: \eng{chunking}، زیرهدف، خلاصه‌سازی و تقسیم مسئله. مسیر دوم معماری بود: بررسی اینکه در \eng{Transformer}، توجه چگونه بازنمایی زمینه‌مند می‌سازد و چه تغییرهایی مانند \eng{local attention} می‌توانند فشار محاسباتی و حافظه‌ای را کم کنند.

نکتهٔ مشترک هر دو مسیر این است که حافظهٔ مؤثر فقط به مقدار فضای خام بستگی ندارد. آنچه اهمیت دارد، سازمان‌دهی، انتخاب، بازیابی و فشرده‌سازی اطلاعات است.

\include{chapters/03-prediction-bias-attention}
\include{chapters/04-language-speech-auditory-models}
\chapter{پاداش، یادگیری تقویتی و تصمیم‌گیری}

\section{سیستم پاداش و یادگیری تقویتی}
پس از پایان بحث مدل‌های شنیداری، بحث وارد موضوع سیستم‌های پاداش در مغز و ارتباط آن با \eng{Reinforcement Learning} شد. نقطهٔ شروع این بحث این است که یک عامل در محیط قرار دارد، در هر لحظه در یک وضعیت یا \eng{state} است، عملی یا \eng{action} انجام می‌دهد و از یک وضعیت به وضعیت دیگر منتقل می‌شود. به این تغییر وضعیت \eng{transition} گفته می‌شود.

هر \eng{transition} می‌تواند با پاداش یا تنبیه همراه باشد. پاداش ممکن است مثبت، منفی، کم، زیاد یا خنثی باشد. این پاداش‌ها نقش \eng{reinforcer} دارند: باعث می‌شوند عامل یاد بگیرد کدام رفتارها را بیشتر تکرار کند و کدام رفتارها را کم‌تر انجام دهد.

هدف این بخش ورود سنگین به معادلات فیزیکی یا ریاضی نبود. ابتدا قرار بود سازوکار کلی مغز در یادگیری از طریق بازخورد فهمیده شود و سپس این بحث به فضای \eng{Reinforcement Learning} وصل شود.

در معماری شناختی‌ای که پیش‌تر دربارهٔ آن صحبت شده بود، عامل از محیط محرک‌های حسی دریافت می‌کند. این محرک‌ها، برای مثال از مسیر ماژول بینایی یا شنیداری، وارد سیستم می‌شوند. سپس با فرایندهایی مانند توجه انتخابی، بخشی از اطلاعات محیط انتخاب می‌شود. اطلاعات فعلی وارد حافظهٔ کوتاه‌مدت می‌شود؛ سیستم از حافظهٔ بلندمدت اطلاعات مرتبط را بازیابی می‌کند؛ با ترکیب اطلاعات فعلی و دانش قبلی، یک الگو استخراج می‌شود؛ و بر اساس آن، عمل انتخاب و از مسیر ماژول حرکتی اجرا می‌شود.

پس از انجام عمل، عامل از زمان \(t\) به زمان \(t+1\) می‌رود و به وضعیت جدیدی منتقل می‌شود. در همین فاصله، پاداش یا تنبیه دریافت می‌شود. در این نقطه باید چیزی شبیه یک \eng{utility module} یا ماژول ارزش‌گذاری به معماری اضافه شود؛ ماژولی که بررسی کند عمل انجام‌شده خوب بود، بد بود، یا تفاوت خاصی با انتظار قبلی نداشت. هدف این ماژول این است که پس از هر تجربه، مدل قبلی یا \eng{prior model} به‌روزرسانی شود.

\subsection{پاداش در مغز}
در انسان، پاداش معمولاً به‌صورت حس رضایت، لذت یا موفقیت تجربه می‌شود. غذا خوردن، موفق شدن در یک کار، دریافت پیام مهم یا دیدن چیزی که قبلاً تجربهٔ خوبی با آن داشته‌ایم، نمونه‌هایی از تجربه‌های پاداش‌دهنده‌اند.

این تجربه‌ها می‌توانند با آزاد شدن دوپامین همراه شوند. نکتهٔ مطرح‌شده در این بحث این بود که دوپامین در واکنش به یک تغییر حالت معنادار آزاد می‌شود؛ تغییری که برای سیستم مهم یا آموخته‌شده است. برای مثال خاموش شدن چراغ در روز ممکن است تجربهٔ خاصی نباشد، اما دریافت یک پیامک به‌دلیل انتظار یا تجربهٔ آموخته‌شده می‌تواند برای فرد پاداش‌دهنده باشد.

\subsection{دوپامین و \eng{VTA}}
وقتی دوپامین آزاد می‌شود، در مغز از مسیرهایی عبور می‌کند که به آن‌ها مسیرهای دوپامینرژیک یا \eng{dopaminergic pathways} گفته می‌شود. در بحث پاداش، یکی از نقاط مهم تولید دوپامین ناحیهٔ \eng{Ventral Tegmental Area} یا \eng{VTA} است.

در چارچوب این بخش، کارکرد مهم \eng{VTA} محاسبهٔ خطای پیش‌بینی پاداش است: سیستم چه پاداشی را پیش‌بینی کرده بود، در واقعیت چه پاداشی دریافت شد، و اختلاف این دو چقدر است. وقتی پاداش واقعی بهتر از انتظار باشد، \eng{VTA} با ترشح دوپامین مرتبط می‌شود. بنابراین دوپامین صرفاً نشانهٔ «لذت» نیست؛ بلکه در این بحث بیشتر به غافلگیری مثبت و خطای پیش‌بینی پاداش مربوط است.

به همین دلیل، هر چیز خوب یا لذت‌بخشی الزاماً همیشه دوپامین زیادی ایجاد نمی‌کند. اگر تجربه‌ای کاملاً قابل پیش‌بینی شده باشد، ممکن است غافلگیری کمی ایجاد کند و در نتیجه پاسخ دوپامینی آن نیز کم‌تر شود.

\subsection{مسیرهای دوپامینرژیک}
چهار مسیر اصلی دوپامینرژیک در مغز مطرح می‌شود، اما در بحث پاداش این بخش تمرکز بر دو مسیر بود:
\[
\text{\lr{VTA}}\rightarrow
\begin{cases}
\text{\lr{Mesolimbic Pathway}}\\
\text{\lr{Mesocortical Pathway}}
\end{cases}
\]

\coursefigure{0.62}{dopaminergic-vta-pathways.png}{مسیرهای دوپامینرژیک مرتبط با \eng{VTA}: ارتباط ناحیهٔ تگمنتوم شکمی با هستهٔ اکومبنس، استریاتوم، قشر پیشانی، هیپوکمپ و تودهٔ سیاه.}

\subsubsection{\eng{Mesolimbic Pathway}}
مسیر \eng{Mesolimbic} از \eng{VTA} به سمت سیستم لیمبیک می‌رود. این مسیر با احساس رضایت، شرطی‌شدن، تقویت رفتار، اعتیاد و میل به تکرار رفتار ارتباط دارد.

برای مثال اگر فرد بعد از کار چای بنوشد و این کار برای او تجربهٔ خوشایندی باشد، سیستم \eng{Mesolimbic} می‌تواند به این الگو شرطی شود. در نتیجه، بعد از کار میل به نوشیدن چای افزایش پیدا می‌کند. به زبان ساده، این مسیر به مغز می‌گوید: این رفتار قبلاً خوب بوده است؛ دوباره انجامش بده.

در توضیح این بخش، مسیر \eng{Mesolimbic} از \eng{VTA} به نواحی لیمبیک، به‌ویژه \eng{Nucleus Accumbens}، می‌رسد. این مسیر بیشتر با تقویت یا تضعیف ارزش محرک‌ها در حافظهٔ بلندمدت ارتباط دارد. وقتی خطای پیش‌بینی پاداش مثبت است، این مسیر کمک می‌کند محرک یا عمل مربوطه در حافظه با ارزش بیشتری ثبت شود.

\subsubsection{\eng{Mesocortical Pathway}}
مسیر \eng{Mesocortical} نیز از \eng{VTA} شروع می‌شود، اما به سمت \eng{Prefrontal Cortex} یا قشر پیش‌پیشانی می‌رود. قشر پیش‌پیشانی با تصمیم‌گیری، تحلیل، ارزیابی پیامدها، کنترل رفتار، بررسی ارزش تلاش و فکر کردن به نتایج بلندمدت ارتباط دارد.

مثلاً وقتی فرد کیک می‌بیند، مسیر \eng{Mesolimbic} ممکن است او را به سمت خوردن کیک سوق دهد، چون کیک قبلاً تجربه‌ای لذت‌بخش بوده است. اما مسیر \eng{Mesocortical} پیامدها را ارزیابی می‌کند: آیا خوردن کیک منطقی است؟ آیا فرد گرسنه است؟ آیا رژیم دارد؟ آیا نتیجهٔ بلندمدت این رفتار خوب است؟ در این بیان استعاری، \eng{Mesolimbic} بیشتر نقش گاز و \eng{Mesocortical} بیشتر نقش تحلیل‌گر یا ترمز را دارد.

\subsection{اختلال‌های مرتبط با سیستم پاداش}
در ادامه، چند پدیده با همین چارچوب توضیح داده شد. هدف این بخش تشخیص پزشکی نیست؛ بلکه نشان دادن این است که نظریهٔ پاداش، دوپامین و خطای پیش‌بینی چگونه می‌تواند برای فهم بعضی تغییرات رفتاری به‌کار رود.

\subsubsection{پیری}
در پیری، سیستم پاداش معمولاً کم‌تر دچار غافلگیری می‌شود. فرد تجربه‌های زیادی کسب کرده و بسیاری از محرک‌ها برای او قابل پیش‌بینی شده‌اند. بنابراین خطای پیش‌بینی پاداش اغلب به صفر نزدیک‌تر است، دوپامین کم‌تر ترشح می‌شود، به‌روزرسانی مدل ذهنی کم‌تر رخ می‌دهد و میل به تجربهٔ چیزهای تازه یا \eng{exploration} کاهش می‌یابد.

از این دید، یکی از راه‌های ایجاد حال بهتر برای فرد سالمند می‌تواند فراهم کردن تجربه‌های تازه و کمتر تجربه‌شده باشد؛ مانند سفرهای جدید، فعالیت‌های جدید یا موقعیت‌هایی که پیش‌تر در زندگی فرد وجود نداشته‌اند.

\subsubsection{افسردگی}
افسردگی نیز در این بحث با همین زبان توضیح داده شد. در این حالت، سیستم پاداش به‌درستی فعال نمی‌شود و فرد ممکن است تجربه کند که چیزی خوشحالش نمی‌کند، کارهایی که قبلاً لذت‌بخش بودند دیگر جذاب نیستند، یا انجام دادن هیچ کاری ارزش ندارد.

در این چارچوب، فرد ممکن است در حالت قطعیت بماند و عمل تازه‌ای انجام ندهد. وقتی تجربهٔ تازه‌ای رخ ندهد، خطای پیش‌بینی پاداش نیز صفر می‌ماند. اما این با یادگیری سالم فرق دارد. در یادگیری سالم، فرد محیط را تجربه می‌کند تا مدل ذهنی‌اش دقیق‌تر شود و خطا کاهش یابد؛ اما در این توصیف از افسردگی، فرد با انجام ندادن عمل و تجربه نکردن، به‌صورت ناسالم به خطای نزدیک به صفر می‌رسد.

\subsubsection{اعتیاد}
اعتیاد یکی از مثال‌های اصلی سوءاستفاده از سیستم پاداش است. در ابتدای مصرف، محرک اعتیادآور می‌تواند خطای پیش‌بینی پاداش مثبت ایجاد کند:
\[
\text{\lr{RPE}} > 0
\]
چون تجربه قوی‌تر یا متفاوت‌تر از انتظار قبلی است. اما پس از تکرار، سیستم یاد می‌گیرد و آن تجربه برای مغز عادی‌تر می‌شود؛ بنابراین خطای پیش‌بینی پاداش به سمت صفر می‌رود:
\[
\text{\lr{RPE}} \rightarrow 0
\]

در اینجا چرخهٔ معیوب شروع می‌شود. فرد برای اینکه دوباره خطای مثبت و پاسخ دوپامینی ایجاد کند، دوز را بالا می‌برد. با افزایش دوز دوباره غافلگیری ایجاد می‌شود، اما پس از مدتی مغز دوباره سازگار می‌شود و خطا کاهش می‌یابد. چرخهٔ کلی چنین است: تجربهٔ اولیه خطای مثبت ایجاد می‌کند؛ سیستم یاد می‌گیرد؛ خطا به صفر نزدیک می‌شود؛ دوز افزایش می‌یابد؛ دوباره خطای مثبت ایجاد می‌شود؛ و مغز دوباره سازگار می‌شود.

مغز در برابر این چرخه منفعل نمی‌ماند. یکی از واکنش‌های دفاعی آن کاهش حساسیت به دوپامین است. در نتیجه، فرد نه‌تنها از محرک اعتیادآور کم‌تر لذت می‌برد، بلکه ممکن است از محرک‌های عادی زندگی نیز لذت کم‌تری تجربه کند.

در این بحث برای این پرسش که چرا همهٔ افراد وارد چرخهٔ اعتیاد نمی‌شوند، چند عامل مطرح شد. برخی افراد اصلاً در معرض محرکی قرار نمی‌گیرند که خطای مثبت شدید ایجاد کند. آگاهی و محیط نیز مهم‌اند، چون شناخت پیامدهای اعتیاد می‌تواند از ورود به چرخه جلوگیری کند. عامل دیگر تفاوت‌های فردی، ژنتیک و قدرت مسیر \eng{Mesocortical} است. اگر این مسیر در کنترل، راهبرد و ارزیابی پیامدهای بلندمدت قوی‌تر عمل کند، فرد بهتر می‌تواند در برابر پاداش فوری مقاومت کند.

مثال مطرح‌شده \eng{Marshmallow Test} بود: آزمایشی که در آن به کودک گفته می‌شود اگر اکنون شیرینی یا شکلات را نخورد، بعداً پاداش بیشتری می‌گیرد. برخی کودکان می‌توانند صبر کنند و برخی نه. این مثال برای توضیح خودکنترلی و توانایی دیدن پیامدهای بلندمدت استفاده شد.

\subsection{ارتباط سیستم پاداش با \eng{Predictive Coding}}
در این بخش پرسیده شد کدام نظریهٔ شناختی می‌تواند این فرایند را توضیح دهد. پاسخ اصلی \eng{Predictive Coding} بود. در کدگذاری پیش‌بینانه، سیستم همواره میان دانش قبلی یا \eng{prior knowledge} و واقعیت مشاهده‌شده مقایسه انجام می‌دهد. آنچه برای یادگیری مهم است، صرفاً خود واقعیت خام نیست؛ بلکه خطای میان پیش‌بینی و واقعیت است.

در سیستم پاداش نیز همین اتفاق می‌افتد. عامل پیش از عمل، انتظاری از پاداش دارد. سپس عمل انجام می‌دهد و پاداش واقعی را دریافت می‌کند. اختلاف میان پاداش واقعی و پاداش پیش‌بینی‌شده همان چیزی است که باید یاد گرفته شود.

\subsection{\eng{Reward Prediction Error}}
مفهوم مرکزی این بخش \eng{Reward Prediction Error} یا خطای پیش‌بینی پاداش است:
\[
\text{\lr{Reward Prediction Error}} =
\text{\lr{Actual Reward}} - \text{\lr{Predicted Reward}}
\]
یا به‌شکل خلاصه:
\[
\text{\lr{RPE}} = R_{\text{\lr{actual}}} - R_{\text{\lr{predicted}}}
\]

اگر پاداش واقعی بیشتر از پاداش پیش‌بینی‌شده باشد، خطا مثبت است:
\[
\text{\lr{RPE}} > 0
\]
یعنی نتیجه بهتر از انتظار بوده است. در این حالت سیستم غافلگیر می‌شود و عمل یا محرک مربوطه تقویت می‌شود. مثال رستورانی بود که فرد انتظار متوسطی از آن دارد، اما تجربهٔ واقعی بسیار بهتر از انتظار اوست. در این حالت، ارزش آن رستوران در ذهن افزایش پیدا می‌کند.

اگر پاداش واقعی تقریباً برابر با پیش‌بینی قبلی باشد، خطا نزدیک به صفر است:
\[
\text{\lr{RPE}} = 0
\]
در این حالت اتفاق تازه‌ای برای یادگیری رخ نمی‌دهد و مدل قبلی تقریباً ثابت می‌ماند.

اگر پاداش واقعی کم‌تر از پاداش پیش‌بینی‌شده باشد، خطا منفی است:
\[
\text{\lr{RPE}} < 0
\]
یعنی نتیجه بدتر از انتظار بوده است. در این حالت سیستم باید یاد بگیرد که ارزش آن عمل یا محرک را کاهش دهد.

\subsection{\eng{Free Energy Principle} و کاهش \eng{Surprise}}
در ادامه، \eng{Free Energy Principle} مطرح شد. بر اساس این نگاه، انسان فقط به‌دنبال دریافت پاداش بیشتر نیست؛ بلکه تلاش می‌کند عدم قطعیت و خطای پیش‌بینی را کاهش دهد. سیستم شناختی محیط را تجربه می‌کند تا مدل ذهنی‌اش دقیق‌تر شود و اختلاف میان پیش‌بینی و واقعیت کم‌تر شود.

هدف نهایی این است که خطای پیش‌بینی پاداش به صفر نزدیک شود:
\[
\text{\lr{RPE}} \rightarrow 0
\]
یعنی سیستم به حالتی برسد که پاداش واقعی با پاداش پیش‌بینی‌شده همخوان شود. در این حالت مدل ذهنی به‌روزرسانی شده و غافلگیری زیادی رخ نمی‌دهد.

\section{مدل‌های محاسباتی یادگیری تقویتی}
از دید محاسباتی، سیستم پاداش انسان شباهت زیادی به \eng{Reinforcement Learning} دارد. در یادگیری تقویتی، عامل در محیط قرار دارد، عملی انجام می‌دهد، از یک وضعیت به وضعیت دیگر می‌رود، پاداش دریافت می‌کند و بر اساس آن پاداش یاد می‌گیرد در آینده چه رفتاری انجام دهد.

\subsection{مسئلهٔ \eng{Reinforcement Learning}}
مسئلهٔ اصلی \eng{Reinforcement Learning} یادگیری از طریق تعامل با محیط است. عامل برخلاف مدل نظارت‌شده، از ابتدا برچسب درست هر تصمیم را ندارد؛ باید عملی انجام دهد و پیامد آن را ببیند. اگر پاداش دریافت‌شده مطلوب باشد، احتمال تکرار آن رفتار بیشتر می‌شود و اگر پیامد نامطلوب باشد، احتمال آن رفتار کاهش می‌یابد.

یادگیری تقویتی وقتی اهمیت ویژه پیدا می‌کند که محیط آنلاین باشد؛ یعنی عامل مدل آماده و دقیقی از محیط ندارد و باید تجربه را نمونه‌به‌نمونه به‌دست آورد. در هر گام، عامل یک \eng{state} را مشاهده می‌کند، یک \eng{action} انجام می‌دهد، به وضعیت جدیدی منتقل می‌شود و از محیط \eng{reward} یا \eng{punishment} دریافت می‌کند. سپس مدل خود را بر اساس همین تجربه به‌روزرسانی می‌کند.

حالت مهم دیگر، محیط تصادفی یا \eng{stochastic} است. در چنین محیطی، اگر در یک وضعیت مشخص یک عمل مشخص انجام شود، لزوماً همیشه به یک وضعیت مشخص نمی‌رسیم. گذار کلی چنین است:
\[
S \xrightarrow{A} S'
\]
اما \(S'\) می‌تواند چند حالت مختلف داشته باشد؛ برای مثال ممکن است با احتمال ۰٫۵ به \(S'_1\) و با احتمال ۰٫۲ به \(S'_2\) برسیم. چون مدل گذار از ابتدا معلوم نیست، عامل باید با تجربه کردن و گرفتن بازخورد آن را یاد بگیرد.

\subsection{تفاوت با یادگیری نظارت‌شده و بدون‌ناظر}
در \eng{Supervised Learning}، داده‌ها برچسب دارند و مدل از رابطهٔ میان \eng{input} و \eng{label} یاد می‌گیرد. اما در سیستم پاداش و یادگیری تقویتی، برچسب آماده وجود ندارد. عامل باید \eng{action} انجام دهد، از محیط \eng{feedback} بگیرد و بر اساس نتیجهٔ تجربه یاد بگیرد.

در \eng{Unsupervised Learning} نیز برچسب مستقیم وجود ندارد، اما مدل معمولاً از ساختار داده‌ها الگو استخراج می‌کند؛ مثلاً خوشه‌بندی می‌کند یا بازنمایی می‌آموزد. تفاوت \eng{Reinforcement Learning} این است که مسئله فقط نبودن برچسب نیست. نکتهٔ اصلی تعامل فعال با محیط و دریافت \eng{reward} پس از عمل است.

\subsection{\eng{Temporal Difference Learning}}
در پایان این بخش، بحث به فرمول یادگیری رسید. فرض کنیم عامل در وضعیت فعلی \(S\) قرار دارد و با انجام عمل \(A\) به وضعیت بعدی \(S'\) می‌رود. عامل برای وضعیت فعلی یک ارزش قبلی دارد که آن را \(V_t\) می‌نامیم. پس از تجربهٔ جدید، یک نمونهٔ تازه از ارزش آن تجربه به‌دست می‌آید. اکنون سیستم باید ارزش قبلی را با توجه به تجربهٔ جدید به‌روزرسانی کند.

نه دانش قبلی به‌تنهایی کافی است و نه نمونهٔ جدید به‌تنهایی. بنابراین سیستم میان این دو درون‌یابی می‌کند. برای این کار از ضریب \(\alpha\) استفاده می‌شود:
\[
0 \leq \alpha \leq 1
\]
\(\alpha\) همان \eng{learning rate} است. اگر \(\alpha = 1\) باشد، فقط تجربهٔ جدید مهم است و دانش قبلی نادیده گرفته می‌شود. اگر \(\alpha = 0\) باشد، تجربهٔ جدید اثری ندارد و سیستم چیزی یاد نمی‌گیرد.

فرمول کلی به‌روزرسانی چنین است:
\[
V_{t+1} = V_t + \alpha(\text{\lr{Sample}} - V_t)
\]
ترم داخل پرانتز همان خطای پیش‌بینی پاداش است:
\[
\text{\lr{Sample}} - V_t = \text{\lr{RPE}}
\]
پس:
\[
V_{t+1} = V_t + \alpha \times \text{\lr{RPE}}
\]

این منطق \eng{Temporal Difference Learning} نامیده می‌شود، چون اختلاف میان دو زمان را محاسبه می‌کند: پیش از عمل در زمان \(t\) و پس از عمل در زمان \(t+1\). سیستم ارزش قبلی را به اندازهٔ نرخ یادگیری و بر اساس خطای پیش‌بینی به‌روزرسانی می‌کند.

هدف نهایی این فرایند آن است که نمونهٔ تجربه‌شده با ارزش پیش‌بینی‌شده همگرا شود:
\[
\text{\lr{Sample}} = V_t
\]
در نتیجه:
\[
\text{\lr{RPE}} = 0
\]
این همان ایده‌ای است که در بحث \eng{Free Energy Principle} نیز دیده شد: کاهش غافلگیری و کاهش خطای پیش‌بینی.

\subsection{\eng{Discount Factor}}
در این بخش، هدف اصلی عمیق‌تر کردن مفاهیم بخش‌های قبل بود؛ به‌ویژه \eng{Q-Learning}، \eng{Active RL} و ارتباط آن‌ها با سیستم پاداش. نقطهٔ شروع این بود که عامل در یک فضای جست‌وجو حرکت می‌کند. ممکن است دربارهٔ بعضی حالت‌ها تجربهٔ قبلی یا \eng{prior knowledge} داشته باشد، و دربارهٔ بعضی حالت‌ها هیچ تجربه‌ای نداشته باشد؛ در این حالت، مقدار اولیه معمولاً صفر در نظر گرفته می‌شود.

فرض کنید عامل در یک حالت شروع قرار دارد و دو انتخاب دارد: رفتن به چپ یا رفتن به راست. اگر عامل فقط پاداش فوری را ببیند، ممکن است بگوید چپ پاداش ۱ دارد و راست فعلاً پاداش ۰ دارد؛ پس در حالت کاملاً \eng{greedy} به چپ می‌رود. اما ممکن است رفتن به راست در آینده به پاداش بزرگ‌تری، مثلاً ۱۰، برسد. بنابراین عامل فقط نباید به پاداش فوری نگاه کند؛ باید تا حدی آینده را نیز در نظر بگیرد.

برای مدل کردن میزان اهمیت آینده از \eng{Discount Factor} یا ضریب تخفیف استفاده می‌کنیم. این ضریب معمولاً با \(\gamma\) نشان داده می‌شود:
\[
0 \leq \gamma \leq 1
\]
اگر \(\gamma = 0\) باشد، عامل فقط پاداش فوری را می‌بیند و کاملاً کوتاه‌نگر است. اگر \(\gamma = 1\) باشد، عامل آینده را بدون تخفیف کامل لحاظ می‌کند. اگر \(0 < \gamma < 1\) باشد، پاداش‌های آینده در محاسبه وارد می‌شوند، اما هرچه دورتر باشند اثر کم‌تری دارند:
\[
r_t + \gamma r_{t+1} + \gamma^2 r_{t+2} + \gamma^3 r_{t+3} + \cdots
\]

نکتهٔ مهم دیگر این است که مفهوم ضریب تخفیف را می‌توان تا حدی با سازوکارهای زیستی نیز مرتبط دانست. برای مثال دربارهٔ سروتونین گفته شد که اگر سطح آن بالاتر باشد، فرد ممکن است کم‌تر عجولانه و کم‌تر کاملاً \eng{greedy} تصمیم بگیرد و توان بیشتری برای در نظر گرفتن آینده داشته باشد. البته این رفتار فقط به یک مادهٔ زیستی وابسته نیست و ترکیبی از عوامل مختلف می‌تواند روی چیزی شبیه ضریب تخفیف اثر بگذارد.

\subsection{\eng{Value Function}}
در \eng{TD Learning} هدف به‌روزرسانی مقدار یک حالت است:
\[
V(s)
\]
فرمول کامل‌تر به‌روزرسانی با ضریب تخفیف چنین است:
\[
V_{t+1}(s) = (1-\alpha)V_t(s) + \alpha \left[r + \gamma V_t(s')\right]
\]
در این فرمول، \(V_t(s)\) مقدار قبلی حالت \(s\)، \(r\) پاداش دریافت‌شده، \(s'\) حالت بعدی، \(V_t(s')\) مقدار حالت بعدی، \(\alpha\) نرخ یادگیری و \(\gamma\) ضریب تخفیف است.

این فرمول نوعی درون‌یابی میان دانش قبلی و تجربهٔ جدید است. بخش \((1-\alpha)V_t(s)\) دانش قبلی را نگه می‌دارد و بخش \(\alpha[r+\gamma V_t(s')]\) تجربهٔ جدید را وارد می‌کند.

مثال سه تجربه داشت:
\[
A \xrightarrow{a_1,\, r=-3} B
\]
\[
B \xrightarrow{a_1,\, r=4} A
\]
\[
A \xrightarrow{a_2,\, r=-4} A
\]
فرض شد \(\alpha=0.1\)، \(\gamma=1\)، و در ابتدا \(V(A)=0\) و \(V(B)=0\).

در تجربهٔ اول، از \(A\) به \(B\) با پاداش \(-3\) می‌رویم:
\[
V(A)=0.9V(A)+0.1[-3+V(B)]
\]
چون مقدارهای اولیه صفرند:
\[
V(A)=0.9(0)+0.1[-3+0]=-0.3
\]
پس مقدار \(A\) کاهش پیدا می‌کند. در تجربهٔ دوم:
\[
V(B)=0.9V(B)+0.1[4+V(A)]
\]
و چون \(V(A)=-0.3\):
\[
V(B)=0.9(0)+0.1[4-0.3]=0.37
\]
در تجربهٔ سوم دوباره \(V(A)\) به‌روزرسانی می‌شود:
\[
V(A)=0.9(-0.3)+0.1[-4-0.3]
\]
نکتهٔ اصلی این است که در \eng{TD Learning}، در هر مرحله مقدار همان حالتی به‌روزرسانی می‌شود که عامل از آن حرکت کرده است.

\subsection{\eng{Q Function}}
تا اینجا مقدار حالت‌ها را داشتیم، اما هدف نهایی عامل فقط دانستن ارزش حالت‌ها نیست. عامل باید در هر لحظه تصمیم بگیرد کدام عمل را انتخاب کند. بنابراین لازم است بدانیم اگر عامل در حالت \(s\) باشد و عمل \(a\) را انجام دهد، این انتخاب چقدر ارزش دارد.

این مقدار با \(Q(s,a)\) نشان داده می‌شود. \(Q(s,a)\) ارزش انجام عمل \(a\) در حالت \(s\) است. رابطهٔ کلی میان \(V\) و \(Q\) چنین است:
\[
V(s)=\max_a Q(s,a)
\]
یعنی ارزش یک حالت برابر است با بیشترین مقدار \(Q\) میان عمل‌های ممکن در آن حالت.

در دنیای واقعی، اگر یک عمل را انتخاب کنیم، همیشه دقیقاً به همان حالتی که انتظار داریم نمی‌رسیم. برای مثال ممکن است عامل تصمیم بگیرد به چپ برود؛ با احتمال ۰٫۹ واقعاً به خانهٔ چپ برسد، اما با احتمال ۰٫۱ به حالت دیگری منتقل شود. این احتمال گذار با تابع \(T(s,a,s')\) نشان داده می‌شود:
\[
T(s,a,s')
\]
یعنی احتمال اینکه عامل از حالت \(s\)، با انجام عمل \(a\)، به حالت \(s'\) برود. در این حالت می‌توان مقدار \(Q\) را با جمع روی حالت‌های بعدی ممکن نوشت:
\[
Q(s,a)=\sum_{s'} T(s,a,s')\left[R(s,a,s')+\gamma V(s')\right]
\]
در این رابطه، عامل هم پاداش فوری را در نظر می‌گیرد و هم ارزش حالت بعدی را، اما ارزش آینده با ضریب تخفیف وارد می‌شود.

\subsection{\eng{Q-Learning}}
در \eng{Q-Learning} به‌جای اینکه فقط \(V(s)\) به‌روزرسانی شود، مستقیماً \(Q(s,a)\) به‌روزرسانی می‌شود:
\[
Q_{t+1}(s,a) = (1-\alpha)Q_t(s,a)
+ \alpha \left[r+\gamma \max_{a'}Q_t(s',a')\right]
\]
در این فرمول، \(s\) حالت فعلی، \(a\) عمل انجام‌شده، \(r\) پاداش دریافت‌شده، \(s'\) حالت بعدی، و \(a'\) عمل‌های ممکن در حالت بعدی هستند. عبارت \(\max_{a'}Q_t(s',a')\) بهترین عمل ممکن در حالت بعدی را نشان می‌دهد.

تفاوت مهم \eng{Q-Learning} با \eng{TD Learning} ساده این است که در \eng{Q-Learning} مستقیماً ارزش عمل‌ها در حالت‌ها یاد گرفته می‌شود، نه فقط ارزش خود حالت‌ها. به همین دلیل، \eng{Q-Learning} یک گام به تصمیم‌گیری واقعی نزدیک‌تر است.

برای مثال، همان اپیزود قبل را در نظر می‌گیریم:
\[
A \xrightarrow{a_1,\, r=-3} B,\qquad
B \xrightarrow{a_1,\, r=4} A,\qquad
A \xrightarrow{a_2,\, r=-4} A
\]
فرض می‌کنیم \(\alpha=0.1\)، \(\gamma=1\)، و همهٔ مقدارهای \(Q\) در ابتدا صفرند.

در تجربهٔ اول، عامل از \(A\) با عمل ۱ به \(B\) می‌رود و پاداش \(-3\) می‌گیرد:
\[
Q(A,1)=0.9Q(A,1)+0.1[-3+\max_{a'}Q(B,a')]
\]
چون هنوز دربارهٔ عمل‌های حالت \(B\) چیزی نمی‌دانیم، بیشینه صفر است:
\[
Q(A,1)=0.1[-3]=-0.3
\]

در تجربهٔ دوم، عامل از \(B\) با عمل ۱ به \(A\) می‌رود و پاداش ۴ می‌گیرد:
\[
Q(B,1)=0.9Q(B,1)+0.1[4+\max_{a'}Q(A,a')]
\]
در حالت \(A\)، مقدار \(Q(A,1)\) برابر \(-0.3\) است و عمل‌های دیگر هنوز مقدار صفر دارند؛ پس بیشینه صفر می‌شود:
\[
Q(B,1)=0.1[4+0]=0.4
\]

در تجربهٔ سوم، عامل از \(A\) با عمل ۲ به \(A\) می‌رود و پاداش \(-4\) می‌گیرد:
\[
Q(A,2)=0.9Q(A,2)+0.1[-4+\max_{a'}Q(A,a')]
\]
در این مرحله، بیشینهٔ مقدارهای \(Q\) در \(A\) همچنان صفر است، پس:
\[
Q(A,2)=0.1[-4+0]=-0.4
\]
روند به همین شکل ادامه پیدا می‌کند. در هر تجربه، مقدار \(Q(s,a)\) مربوط به همان حالت و همان عمل به‌روزرسانی می‌شود.

\subsection{محیط‌های تصادفی و عدم قطعیت}
در محیط قطعی، انجام یک عمل مشخص در یک وضعیت مشخص همیشه به نتیجه‌ای مشخص می‌رسد. اما در محیط تصادفی یا غیرقطعی، یک عمل می‌تواند نتایج مختلفی داشته باشد. مثال رانندگی بود: در یک محیط ایده‌آل، اگر راننده فرمان را به چپ بچرخاند، انتظار داریم ماشین به چپ برود؛ اما در محیطی غیرقطعی ممکن است ماشین کمی منحرف شود، متوقف شود یا رفتار دیگری نشان دهد.

بنابراین عامل باید با عدم قطعیت در گذارها کنار بیاید. یادگیری در این حالت یعنی تجربه کردن، دریافت بازخورد و اصلاح تدریجی مدل از اینکه هر عمل در هر وضعیت معمولاً به چه پیامدهایی منجر می‌شود.

\subsection{\eng{Exploration} و \eng{Exploitation}}
اگر عامل فقط از رابطهٔ
\[
V(s)=\max_a Q(s,a)
\]
برای انتخاب عمل استفاده کند، همیشه عملی را انتخاب می‌کند که تا آن لحظه بهترین تجربهٔ شناخته‌شده را داشته است. این حالت کاملاً \eng{greedy} یا \eng{exploitative} است؛ یعنی عامل فقط از تجربه‌های قبلی خود استفاده می‌کند.

اما در تصمیم‌گیری واقعی، همیشه فقط بهترین تجربهٔ قبلی را تکرار نمی‌کنیم. مثال خرید این نکته را روشن می‌کند: معمولاً همیشه فقط از جایی خرید نمی‌کنیم که قبلاً بهترین تجربه را از آن داشته‌ایم؛ گاهی مکان‌های جدید را هم امتحان می‌کنیم، به‌ویژه اگر فکر کنیم شاید گزینهٔ بهتری وجود داشته باشد.

\eng{Exploitation} یعنی عامل از بهترین تجربه‌های قبلی خود استفاده کند: «تا الان این عمل بهترین بوده، پس دوباره همان را انتخاب می‌کنم.» در مقابل، \eng{Exploration} یعنی عامل به عمل‌هایی هم فرصت بدهد که کم‌تر امتحان شده‌اند: «این عمل هنوز بهترین شناخته نشده، اما چون کم‌تر امتحان شده، شاید ارزش کشف کردن داشته باشد.» تصمیم‌گیری در دنیای واقعی معمولاً ترکیبی از این دو است.

\subsection{\eng{Passive RL} و \eng{Active RL}}
برای اینکه عامل فقط \eng{greedy} نباشد، باید عدم قطعیت را هم وارد تصمیم‌گیری کنیم. اگر امتیاز یک عمل فقط برابر با مقدار \(Q\) باشد:
\[
\text{\lr{Score}}(s,a)=Q(s,a)
\]
عامل بیشتر به سمت بهره‌برداری از تجربه‌های قبلی می‌رود. اما می‌توان جمله‌ای مربوط به عدم قطعیت نیز اضافه کرد:
\[
\text{\lr{Score}}(s,a)=Q(s,a)+\beta U(s,a)
\]
در این رابطه، \(Q(s,a)\) تجربهٔ قبلی و پاداش یادگرفته‌شده را نشان می‌دهد، \(U(s,a)\) میزان عدم قطعیت دربارهٔ آن عمل است، و \(\beta\) مشخص می‌کند عامل چقدر به عدم قطعیت اهمیت می‌دهد.

اگر عملی کم‌تر امتحان شده باشد، عدم قطعیت آن بیشتر است. بنابراین حتی اگر مقدار \(Q\) آن فعلاً خیلی بالا نباشد، ممکن است به خاطر عدم قطعیت بالا انتخاب شود تا عامل آن را کشف کند. این ایده به روش‌هایی مانند \eng{Upper Confidence Bound} یا \eng{UCB} شبیه است؛ یعنی عامل فقط میانگین تجربه‌شده را نگاه نمی‌کند، بلکه برای گزینه‌های کم‌تر امتحان‌شده یک حد بالای خوش‌بینانه هم در نظر می‌گیرد.

در حالت \eng{passive}، عامل بیشتر بر اساس تجربه‌های قبلی تصمیم می‌گیرد؛ بهترین گزینه‌های شناخته‌شده را انتخاب می‌کند، تمایل کمی به امتحان گزینه‌های جدید دارد و بیشتر \eng{exploit} می‌کند. در حالت \eng{active}، عامل علاوه بر تجربه‌های قبلی، عدم قطعیت را هم در نظر می‌گیرد؛ گزینه‌های کم‌تر امتحان‌شده را بررسی می‌کند و میان \eng{exploration} و \eng{exploitation} تعادل برقرار می‌کند.

برای توضیح شهودی گفت می‌توان حالت \eng{active} را شبیه جوانی دانست که بیشتر دوست دارد تجربه و کشف کند، و حالت \eng{passive} را شبیه پیری دانست که بیشتر به تجربه‌های قبلی اتکا دارد. این تشبیه برای فهم شهودی است، نه یک ادعای زیستی دقیق.

\chapter{شبکه‌های اسپایکی و محاسبات نورومورفیک}

\section{مبانی شبکه‌های عصبی اسپایکی}
پس از بحث سیستم پاداش و خطای پیش‌بینی پاداش، بحث وارد شبکه‌های عصبی اسپایکی یا \eng{Spiking Neural Networks} شد. هدف این بود که ابتدا معماری کلی این شبکه‌ها فهمیده شود و سپس ببینیم چگونه می‌توان سیستم پاداش را از دریچهٔ شبکه‌های اسپایکی مدل‌سازی کرد.

\subsection{چرا شبکه‌های اسپایکی؟}
در شبکه‌های عصبی مصنوعی معمول یا \eng{ANN}، ورودی‌ها اغلب اعداد پیوسته‌اند؛ برای مثال \eng{embedding}ها یا بردارهای عددی. این ورودی‌ها در وزن‌ها ضرب می‌شوند، جمع وزن‌دار ساخته می‌شود، \eng{bias} اضافه می‌شود، مقدار از \eng{activation function} عبور می‌کند و در نهایت خروجی یا پیش‌بینی تولید می‌شود. سپس پیش‌بینی با پاسخ درست یا \eng{ground truth} مقایسه می‌شود و اگر خطا وجود داشته باشد، وزن‌ها با \eng{backpropagation} اصلاح می‌شوند.

اما این تصویر از نظر زیستی کاملاً واقع‌گرایانه نیست. دو تفاوت مهم در این بحث برجسته شد: مغز با اعداد پیوسته‌ای شبیه ورودی‌های معمول \eng{ANN} کار نمی‌کند، و مغز \eng{backpropagation} کلاسیک و سراسری ندارد. در مغز، یادگیری بیشتر محلی است؛ هر نورون تا حد زیادی بر اساس ورودی‌ها و وضعیت خودش تصمیم می‌گیرد، نه اینکه کل شبکه منتظر خروجی نهایی بماند و سپس گرادیان سراسری به عقب برگردد.

شبکه‌های اسپایکی تلاش می‌کنند به این تصویر زیستی نزدیک‌تر شوند. در این شبکه‌ها، ارتباط نورون‌ها با رخدادهای گسستهٔ زمانی انجام می‌شود: اسپایک رخ می‌دهد یا رخ نمی‌دهد. به همین دلیل، این شبکه‌ها برای مدل‌سازی‌ای اهمیت دارند که می‌خواهد زمان، رویداد و مصرف انرژی را جدی بگیرد.

\subsection{تفاوت \eng{ANN} و \eng{SNN}}
در \eng{ANN} معمولاً با مقدارهای پیوسته سروکار داریم. نورون‌ها در هر گام مقدار عددی دریافت می‌کنند و خروجی عددی تولید می‌کنند. اما در \eng{SNN}، ورودی و خروجی نورون‌ها به‌شکل اسپایک است. اگر نورون اسپایک بزند، مقدار آن را می‌توان ۱ در نظر گرفت؛ اگر اسپایک نزند، مقدار ۰ است.

بنابراین در \eng{SNN}، ورودی یک دنبالهٔ زمانی از صفر و یک‌هاست:
\[
S = [0, 1, 0, 1, 1, 0, \ldots]
\]
وقتی مقدار ۱ باشد، یعنی اسپایک رخ داده است؛ و وقتی مقدار ۰ باشد، یعنی اسپایکی وجود ندارد. هر اسپایک در وزن سیناپسی ضرب می‌شود:
\[
S_i \times W_i
\]
اگر مجموع ورودی‌های وزن‌دار باعث شود ولتاژ نورون از آستانه عبور کند، نورون پس‌سیناپسی نیز اسپایک می‌زند.

\coursefigure{0.70}{artificial-neuron.png}{نورون مصنوعی کلاسیک در \eng{ANN}: ورودی‌های پیوستهٔ وزن‌دار جمع می‌شوند، از تابع فعال‌ساز عبور می‌کنند و خروجی عددی تولید می‌شود.}

\coursefigure{0.70}{spiking-neuron.png}{نورون اسپایکی در \eng{SNN}: ورودی‌های وزن‌دار به پتانسیل غشا تبدیل می‌شوند و خروجی به‌صورت رخداد شلیک اسپایک ظاهر می‌شود.}

\subsection{نورون زیستی}
برای فهم \eng{SNN}، ابتدا ساختار سادهٔ نورون زیستی مرور شد. یک نورون زیستی را می‌توان به‌صورت ساده شامل دندریت، بدنهٔ سلولی، آکسون و سیناپس دانست.

\subsubsection{دندریت}
دندریت محل دریافت ورودی است. سیگنال‌های نورون‌های دیگر از طریق سیناپس‌ها به دندریت می‌رسند.

\subsubsection{\eng{Soma}}
\eng{Soma} یا بدنهٔ سلولی بخش پردازشی نورون است. در توضیح ساده، اثر ورودی‌ها در این بخش جمع می‌شود و وضعیت نورون برای اسپایک زدن یا نزدن تعیین می‌شود.

\subsubsection{آکسون}
آکسون خروجی نورون را به نورون‌های دیگر منتقل می‌کند. اگر نورون اسپایک بزند، این رخداد از مسیر آکسون به سمت سیناپس‌ها منتقل می‌شود.

\subsubsection{سیناپس}
سیناپس محل اتصال آکسون یک نورون به دندریت نورون دیگر است. اگر نورونی به نورون دیگر اطلاعات بفرستد، نورون اول \eng{presynaptic neuron} و نورون دوم \eng{postsynaptic neuron} نامیده می‌شود. نورون پیش‌سیناپسی اطلاعات را می‌فرستد و نورون پس‌سیناپسی آن را دریافت می‌کند.

\subsection{اسپایک و زمان‌محور بودن محاسبه}
در \eng{SNN} زمان بخشی اصلی از محاسبه است. در \eng{ANN} معمولاً با یک بردار عددی برای یک شیء یا ورودی کار می‌کنیم، اما در \eng{SNN} ورودی یک توالی زمانی است. مهم است کدام اسپایک زودتر آمده، کدام بعداً آمده و فاصلهٔ زمانی میان آن‌ها چقدر بوده است.

به همین دلیل، خروجی \eng{SNN} نیز فقط یک مقدار عددی ساده نیست؛ بلکه معمولاً باید الگوی اسپایک‌ها در طول زمان تفسیر شود. هر اسپایک می‌تواند در همان لحظه بخشی از تصمیم‌گیری را فعال کند، بی‌آنکه لازم باشد همهٔ نورون‌های شبکه درگیر محاسبه شوند.

\subsection{مدل \eng{LIF}}
برای مدل‌سازی نورون اسپایکی، مدل‌های مختلفی وجود دارد؛ از جمله \eng{Leaky Integrate-and-Fire} یا \eng{LIF}، مدل \eng{Izhikevich} و مدل \eng{Hodgkin-Huxley}. در این بخش تمرکز روی مدل \eng{LIF} بود، چون هم ساده است و هم برخی رفتارهای مهم نورون زیستی را نشان می‌دهد.

فرمول سادهٔ مدل چنین است:
\[
V(t+1) = \lambda V(t) + I(t)
\]
در این رابطه، \(V(t)\) ولتاژ قبلی، \(\lambda\) ضریب نشت یا \eng{leak factor} و \(I(t)\) جریان ورودی است. جریان ورودی معمولاً از اسپایک‌های ورودی و وزن‌های سیناپسی ساخته می‌شود:
\[
I(t) = \sum_i W_i S_i(t)
\]
اگر ولتاژ از آستانه عبور کند:
\[
V(t+1) \geq V_{\text{\lr{threshold}}}
\]
نورون اسپایک می‌زند.

\coursefigure{0.78}{leaky-integrate-and-fire-neuron.png}{مدل \eng{LIF}: اسپایک‌های ورودی وزن‌دار در دندریت و سوما جمع می‌شوند، ولتاژ غشا با نشت زمانی تغییر می‌کند و پس از عبور از آستانه اسپایک تولید می‌شود.}

\subsubsection{\eng{Integration}}
\eng{Integration} یعنی اثر ورودی‌ها با هم جمع می‌شود. اگر نورون از چند نورون پیش‌سیناپسی ورودی بگیرد، اسپایک‌های ورودی در وزن‌های سیناپسی ضرب می‌شوند. سپس اثر آن‌ها در ولتاژ نورون تجمع پیدا می‌کند.

\subsubsection{\eng{Fire}}
\eng{Fire} یعنی اگر ولتاژ تجمع‌یافته از آستانه عبور کند، نورون اسپایک می‌زند. در این لحظه خروجی نورون برابر با رخداد اسپایک است.

\subsubsection{\eng{Leak}}
\eng{Leak} یعنی اگر ورودی کافی وجود نداشته باشد، ولتاژ نورون به‌تدریج کاهش پیدا می‌کند. این کاهش تدریجی همان نشت است و باعث می‌شود اثر ورودی‌ها برای همیشه در نورون باقی نماند.

\subsubsection{\eng{Refractory Period}}
پس از اسپایک، نورون برای مدت کوتاهی وارد دوره‌ای می‌شود که در آن فعالیتی ندارد یا به‌سادگی دوباره اسپایک نمی‌زند. این دوره \eng{refractory period} نامیده می‌شود و از رفتار نورون‌های واقعی الهام گرفته شده است.

\subsection{مصرف انرژی و پردازش رویدادمحور}
یکی از ویژگی‌های مهم \eng{SNN} این است که شبکه می‌تواند همیشه در حالت آماده‌باش یا \eng{on-call} باشد، اما مصرف انرژی پایینی داشته باشد. دلیل آن این است که وقتی ورودی صفر است، محاسبهٔ سنگین انجام نمی‌شود. نورون فقط در صورت وجود اسپایک و عبور از آستانه فعال می‌شود.

در \eng{ANN} معمولاً در هر گام محاسباتی، لایه‌ها درگیر ضرب و جمع می‌شوند، حتی اگر تغییر مهمی در ورودی رخ نداده باشد. اما در \eng{SNN} رخدادهای مهم، یعنی اسپایک‌ها، محاسبه را فعال می‌کنند. به همین دلیل، این شبکه‌ها برای سیستم‌هایی مناسب‌اند که باید همیشه فعال باشند اما مصرف انرژی کمی داشته باشند.

مثال‌های مناسب شامل دستیار صوتی همیشه‌روشن برای تشخیص \eng{wake-up word}، سمعک‌ها، سیستم‌های شنود کم‌مصرف، پهپادها، ربات‌هایی که باید نسبت به مانع یا حمله واکنش سریع نشان دهند، و سیستم‌هایی بود که به واکنش سریع و محلی نیاز دارند.

\section{یادگیری در شبکه‌های اسپایکی}
پس از معرفی معماری \eng{SNN}، بحث به این پرسش رسید که شبکهٔ اسپایکی چگونه یاد می‌گیرد. در مغز، وقتی برخی مسیرهای نورونی بارها فعال می‌شوند، اتصال میان آن‌ها تقویت می‌شود. بنابراین یادگیری در این سطح را می‌توان به تغییر وزن‌های سیناپسی ربط داد:
\[
\Delta W
\]
وقتی یک مسیر بارها استفاده شود، وزن سیناپسی آن افزایش پیدا می‌کند و آن مسیر در شبکه تقویت می‌شود.

برای توضیح این ایده، می‌توان از مثال مسیر روی چمن استفاده کرد. فرض کنید در یک دانشگاه مسیر رسمی برای عبور افراد طراحی شده است، اما پس از مدتی مسیر دیگری روی چمن بیشتر پا می‌خورد و سبزه‌هایش از بین می‌رود. این یعنی افراد بدون دستور مستقیم، مسیر دیگری را بارها انتخاب کرده‌اند. در شبکهٔ نورونی نیز مسیری که بارها فعال شود، وزن بیشتری می‌گیرد و تقویت می‌شود.

\subsection{قانون \eng{Hebbian Learning}}
ایدهٔ پایه برای یادگیری بدون‌ناظر در شبکه‌های اسپایکی، قانون \eng{Hebb} است. عبارت مشهور آن چنین است:
\[
\text{\lr{Fire together, wire together}}
\]
یعنی اگر دو نورون مرتب با هم فعال شوند، اتصال بین آن‌ها تقویت می‌شود.

در اینجا منظور از «با هم فعال شدن» الزاماً هم‌زمانی دقیق نیست. منظور این است که نورون پیش‌سیناپسی اسپایک بزند و نورون پس‌سیناپسی در ادامه فعال شود. اگر نورون پیش‌سیناپسی اسپایک بزند و بعد نورون پس‌سیناپسی نیز اسپایک بزند، وزن سیناپسی میان آن‌ها افزایش پیدا می‌کند.

این اتفاق مهم است، چون نشان می‌دهد میان آن دو نورون نوعی ارتباط اطلاعاتی وجود دارد. ممکن است هر دو نسبت به موضوعی مشابه حساس باشند، یا میان آن‌ها رابطه‌ای علّی وجود داشته باشد؛ یعنی فعال شدن نورون پیش‌سیناپسی در فعال شدن نورون پس‌سیناپسی نقش داشته باشد.

\subsection{\eng{STDP}}
\eng{STDP} مخفف \eng{Spike-Timing-Dependent Plasticity} است. یعنی تغییرپذیری وزن‌ها به زمان‌بندی اسپایک‌ها وابسته است. قانون \eng{Hebb} می‌گوید اگر دو نورون با هم فعال شوند، اتصالشان تقویت می‌شود. اما \eng{STDP} نکتهٔ دقیق‌تری اضافه می‌کند: ترتیب و فاصلهٔ زمانی اسپایک‌ها مهم است.

اگر نورون پیش‌سیناپسی پیش از نورون پس‌سیناپسی اسپایک بزند و به فعال شدن آن کمک کند، وزن سیناپسی باید تقویت شود. اما اگر نورون پس‌سیناپسی پیش از نورون پیش‌سیناپسی فعال شده باشد، احتمالاً فعال شدن پس‌سیناپسی به علت آن نورون پیش‌سیناپسی خاص نبوده است؛ بنابراین وزن میان آن دو باید تضعیف شود.

\subsubsection{\eng{LTP}}
\eng{LTP} یا \eng{Long-Term Potentiation} زمانی رخ می‌دهد که نورون پیش‌سیناپسی قبل از نورون پس‌سیناپسی اسپایک بزند:
\[
t_{\text{\lr{pre}}} < t_{\text{\lr{post}}}
\]
در این حالت رابطهٔ علّی-زمانی معقولی وجود دارد و وزن سیناپسی تقویت می‌شود:
\[
\Delta W > 0
\]

\subsubsection{\eng{LTD}}
\eng{LTD} یا \eng{Long-Term Depression} زمانی رخ می‌دهد که نورون پیش‌سیناپسی پس از نورون پس‌سیناپسی اسپایک بزند:
\[
t_{\text{\lr{pre}}} > t_{\text{\lr{post}}}
\]
در این حالت نورون پیش‌سیناپسی احتمالاً علت فعال شدن نورون پس‌سیناپسی نبوده است، پس وزن سیناپسی تضعیف می‌شود:
\[
\Delta W < 0
\]

\subsubsection{نقش اختلاف زمانی}
در \eng{STDP} فقط مثبت یا منفی بودن اختلاف زمانی مهم نیست؛ اندازهٔ اختلاف زمانی نیز اهمیت دارد. اختلاف زمانی را می‌توان چنین تعریف کرد:
\[
\Delta t = t_{\text{\lr{post}}} - t_{\text{\lr{pre}}}
\]
اگر \(\Delta t > 0\)، یعنی نورون پیش‌سیناپسی زودتر از نورون پس‌سیناپسی فعال شده و این حالت به \eng{LTP} مربوط است. اگر \(\Delta t < 0\)، یعنی نورون پس‌سیناپسی زودتر فعال شده و این حالت به \eng{LTD} مربوط است.

هرچه \(|\Delta t|\) کوچک‌تر باشد، دو اسپایک از نظر زمانی به هم نزدیک‌تر بوده‌اند و ارتباط میان آن‌ها قوی‌تر فرض می‌شود؛ پس تغییر وزن بیشتر خواهد بود. هرچه \(|\Delta t|\) بزرگ‌تر باشد، ارتباط ضعیف‌تر است و تغییر وزن کم‌تر می‌شود.

فرمول شهودی برای تقویت چنین بیان شد:
\[
\Delta W = A_+ e^{-\Delta t / \tau_+}
\]
و برای تضعیف:
\[
\Delta W = -A_- e^{\Delta t / \tau_-}
\]
در این فرمول‌ها، \(A_+\) دامنهٔ تقویت، \(A_-\) دامنهٔ تضعیف، \(\tau_+\) و \(\tau_-\) ثابت‌های زمانی، و \(\Delta t\) اختلاف زمانی اسپایک‌ها هستند.

اگر نورون پیش‌سیناپسی قبل از پس‌سیناپسی و با فاصلهٔ کم اسپایک بزند، تقویت زیاد است. اگر همین ترتیب با فاصلهٔ زیاد رخ دهد، تقویت کم‌تر است. اگر پیش‌سیناپسی بعد از پس‌سیناپسی و با فاصلهٔ کم اسپایک بزند، تضعیف زیاد است؛ و اگر فاصله زیاد باشد، تضعیف کم‌تر است.

اختلاف زمانی زمانی محاسبه می‌شود که هر دو نورون اسپایک زده باشند. اگر فقط نورون پیش‌سیناپسی اسپایک بزند و نورون پس‌سیناپسی فعال نشود، هنوز اختلاف زمانی معناداری برای \eng{STDP} وجود ندارد.

\subsection{\eng{Reward-Modulated STDP}}
تا اینجا دو جزء جداگانه داشتیم: از سیستم پاداش و \eng{TD Learning}، مفهوم خطای پیش‌بینی پاداش را داشتیم:
\[
\text{\lr{Reward Prediction Error}} =
R_{\text{\lr{actual}}} - R_{\text{\lr{predicted}}}
\]
و از شبکه‌های اسپایکی، تغییر وزن بر اساس زمان‌بندی اسپایک‌ها را داشتیم:
\[
\Delta W_{\text{\lr{STDP}}}
\]
هدف \eng{Reward-Modulated STDP} این است که سیستم پاداش روی \eng{STDP} سوار شود. یعنی هم اثر زمان‌بندی اسپایک‌ها لحاظ شود و هم اثر خطای پیش‌بینی پاداش.

فرم کلی تغییر وزن در این بحث چنین بیان شد:
\[
\Delta W = \eta \times \delta \times E(t)
\]
در این رابطه، \(\eta\) نرخ یادگیری، \(\delta\) خطای پیش‌بینی پاداش، \(E(t)\) اثر \eng{STDP} یا \eng{eligibility trace}، و \(\Delta W\) تغییر وزن سیناپسی است.

در این مدل، \eng{STDP} مشخص می‌کند کدام سیناپس‌ها از نظر زمانی در رفتار نقش داشته‌اند، و خطای پیش‌بینی پاداش مشخص می‌کند که نتیجهٔ رفتار باید آن سیناپس‌ها را تقویت کند یا تضعیف. اگر سیناپسی از نظر \eng{STDP} فعال و مرتبط بوده باشد و سپس سیستم پاداش بگوید نتیجه خوب بوده است، آن سیناپس بیشتر تقویت می‌شود. اما اگر نتیجه بدتر از انتظار باشد، همان سیناپس می‌تواند تضعیف شود.

این ایده همچنین به مسئلهٔ نسبت دادن علت‌ها کمک می‌کند. وقتی نورون پس‌سیناپسی اسپایک می‌زند، ممکن است از چند نورون پیش‌سیناپسی ورودی گرفته باشد و دقیقاً ندانیم کدام ورودی علت اصلی عبور از آستانه بوده است. \eng{STDP} ردپای نورون‌هایی را که کمی قبل‌تر اسپایک زده‌اند نگه می‌دارد، و پاداش بعدی تعیین می‌کند این ردپا در جهت تقویت استفاده شود یا تضعیف.

\subsection{از \eng{SNN} پایه به \eng{Deep SNN}}
در بخش‌های قبل، نورون اسپایکی با مدل \eng{LIF} و یادگیری سیناپسی با \eng{STDP} و \eng{Reward-Modulated STDP} بررسی شد. در این بخش، بحث از شبکه‌های کم‌عمق و پایه‌ای به سمت \eng{Deep SNN} رفت؛ یعنی شبکه‌هایی که چندین لایهٔ اسپایکی پشت سر هم دارند.

در یادآوری زیستی، نورون از دندریت، جسم سلولی، آکسون، پایانهٔ آکسونی، شکاف سیناپسی و وزن سیناپسی میان نورون پیش‌سیناپسی و پس‌سیناپسی تشکیل می‌شود. وقتی نورون پیش‌سیناپسی اسپایک می‌زند، این رویداد از راه وزن سیناپسی روی نورون بعدی اثر می‌گذارد و به شکل جریان ورودی در نورون پس‌سیناپسی ظاهر می‌شود.

در مدل \eng{LIF}، نورون جریان‌های ورودی را جمع می‌کند و ولتاژ غشایی آن افزایش می‌یابد. اگر ولتاژ از آستانه عبور کند، نورون اسپایک می‌زند:
\[
V > \theta
\]
پس از اسپایک، ولتاژ ریست می‌شود و نورون وارد دورهٔ \eng{refractory} می‌شود؛ یعنی برای مدتی دوباره فعال نمی‌شود یا فعال‌شدن آن دشوارتر است. خروجی سادهٔ نورون را می‌توان چنین نوشت:
\[
S =
\begin{cases}
1 & V > \theta \\
0 & V \leq \theta
\end{cases}
\]
بنابراین خروجی نورون اسپایکی معمولاً دودویی است: یا اسپایک داریم یا نداریم.

در ساده‌ترین حالت، اگر ورودی نورون \(x\) و وزن سیناپسی \(w\) باشد، جریان یا اثر سیناپسی چنین نوشته می‌شود:
\[
I = xw
\]
و در حالت چند ورودی:
\[
V = x_1w_1 + x_2w_2 + x_3w_3 + \cdots
\]
اگر ولتاژ از آستانه عبور کند، خروجی ۱ می‌شود؛ در غیر این صورت خروجی صفر باقی می‌ماند.

در \eng{Deep SNN} چندین لایهٔ نورون اسپایکی پشت سر هم قرار می‌گیرند. خروجی هر لایه، یعنی دنباله‌ای از اسپایک‌ها، به لایهٔ بعدی منتقل می‌شود و در آنجا دوباره وزن سیناپسی، جریان، ولتاژ و تصمیم اسپایک‌زدن یا نزدن محاسبه می‌شود. از این نظر، شبکه شبیه شبکه‌های عمیق معمولی چندلایه است؛ اما ورودی و خروجی هر لایه به شکل اسپایکی و دودویی است.

\subsection{الهام زیستی از ساختار چندلایهٔ مغز}
ایدهٔ چندلایه‌بودن شبکه‌ها از ساختار مغز نیز الهام می‌گیرد. در سیستم بینایی انسان، ناحیه‌هایی مانند \eng{V1}، \eng{V2}، \eng{V3}، \eng{V4}، \eng{MT/V5} و \eng{IT} در پردازش اطلاعات بینایی نقش دارند. در لایه‌های اولیه، ویژگی‌های ساده‌تری مانند لبه، خط، رنگ و جهت استخراج می‌شوند. هرچه پردازش جلوتر می‌رود، بازنمایی‌ها پیچیده‌تر و انتزاعی‌تر می‌شوند: از لبه و خط به شکل‌های پیچیده‌تر، از شکل‌ها به چهره، و در سطح‌های بالاتر به هویت فرد یا مفهوم شیء.

پس عمیق‌تر شدن شبکه با توانایی ساخت \eng{representation}های پیچیده‌تر ارتباط دارد. البته در مغز، پردازش همیشه یک مسیر خطی واحد نیست. خروجی ناحیه‌های اولیه می‌تواند وارد مسیرهای مختلف شود؛ برخی لایه‌ها میان چند وظیفه مشترک‌اند و سپس مسیرها تخصصی‌تر می‌شوند. این تصویر به ایدهٔ \eng{multi-task learning} شباهت دارد.

در سیستم بینایی، مسیر \eng{ventral} بیشتر با تشخیص «چیستی» شیء و مسیر \eng{dorsal} بیشتر با «کجایی»، حرکت و تعامل با شیء مرتبط دانسته می‌شود. اگر یک مسیر آسیب ببیند، ممکن است فرد شیء را ببیند اما درست با آن تعامل نکند؛ یا برعکس، بتواند با اشیا تعامل کند اما نتواند دقیق تشخیص دهد چه هستند.

این ساختار چندلایه فقط به بینایی محدود نیست. در سیستم لامسه نیز ناحیه‌هایی مانند \eng{S1} یا \eng{Primary Somatosensory Cortex} و \eng{S2} یا \eng{Secondary Somatosensory Cortex} وجود دارند. در \eng{S1} اطلاعات اولیه‌تری مانند لمس، درد و دما پردازش می‌شود و در \eng{S2} اطلاعات حسی ترکیب‌شده‌تر و انتزاعی‌تر شکل می‌گیرند.

\coursefigure{0.72}{visual-processing-hierarchy.png}{سلسله‌مراتب پردازش بینایی از نواحی اولیه مانند \eng{V1} تا مسیرهای بالاتر مانند \eng{V4}، \eng{MT/V5} و قشر گیجگاهی تحتانی.}

\coursefigure{0.78}{dorsal-and-ventral-visual-streams.png}{تقسیم پردازش بینایی به مسیر پشتی \eng{Where/How} و مسیر شکمی \eng{What} که برای فهم الهام زیستی شبکه‌های چندلایه مفید است.}

\subsection{مصرف انرژی و انگیزهٔ \eng{Deep SNN}}
یکی از انگیزه‌های مهم برای مطالعهٔ \eng{SNN}، مصرف انرژی پایین مغز است. مغز انسان با وجود پیچیدگی زیاد، حدود ۲۰ وات انرژی مصرف می‌کند؛ در حالی که مدل‌های عمیق معمولی روی \eng{GPU}ها و سامانه‌های بزرگ می‌توانند انرژی بسیار بیشتری مصرف کنند.

در این بحث دو دلیل برای این بهینگی برجسته شد. نخست، اسپایکی‌بودن است: در مغز، ورودی و خروجی نورون‌ها پیوسته نیست و بیشتر به شکل رویدادهای صفر و یک دیده می‌شود:
\[
0 \quad \text{یا} \quad 1
\]
این دودویی و رویدادمحور بودن می‌تواند مصرف انرژی را کاهش دهد و سطحی از انتزاع ایجاد کند.

دلیل دوم، محلی بودن محاسبات است. در مغز، هر نورون تا حد زیادی بر اساس اطلاعات محلی خود تصمیم می‌گیرد. این وضعیت با شبکه‌های عمیق معمولی متفاوت است که در آن‌ها خطا به‌صورت سراسری محاسبه می‌شود و با \eng{backpropagation} از انتها به ابتدا برمی‌گردد. محلی بودن نیز یکی از عواملی است که به کاهش مصرف انرژی کمک می‌کند.

سؤال اصلی این است که آیا می‌توان \eng{Deep SNN}هایی ساخت که هم توانایی شبکه‌های عمیق را داشته باشند و هم در فاز اجرا به مزیت انرژی و رویدادمحوری مغز نزدیک‌تر شوند. برای پاسخ به این سؤال، باید روش یادگیری در شبکه‌های اسپایکی عمیق را بررسی کرد.

\subsection{محدودیت‌های \eng{STDP} در شبکه‌های عمیق}
\eng{STDP} برای شبکه‌های کوچک یا کم‌عمق می‌تواند مفید باشد، به‌ویژه وقتی یادگیری بدون ناظر است، برچسب و پاداش نداریم، و شبکه قرار است از الگوهای تکرارشوندهٔ ورودی یاد بگیرد. در \eng{Reward-Modulated STDP} نیز همان ایدهٔ زمان‌بندی اسپایک با سیگنال پاداش ترکیب می‌شود تا پاداش تعیین کند ردپای سیناپسی باید تقویت شود یا تضعیف.

اما وقتی تعداد لایه‌ها زیاد می‌شود، \eng{STDP} به‌تنهایی کافی نیست. همگرایی سخت‌تر می‌شود، دقت پایین می‌آید، تنظیم وزن‌ها فقط با قوانین محلی دشوار می‌شود و شبکه برای \eng{fine-tune} کردن وزن‌ها به روش‌های دقیق‌تری نیاز پیدا می‌کند. بنابراین در \eng{Deep SNN}، مخصوصاً در وظایف نظارت‌شده، معمولاً به روش‌هایی شبیه \eng{backpropagation} نیاز داریم.

\subsection{\eng{Surrogate Gradient}}
در یادگیری نظارت‌شده، داده‌ها و برچسب‌ها مشخص‌اند. شبکه ورودی را می‌گیرد، خروجی تولید می‌کند و آن را با \eng{label} مقایسه می‌کند. سپس وزن‌ها اصلاح می‌شوند. در شبکه‌های عمیق معمولی، این اصلاح با \eng{backpropagation} انجام می‌شود.

مشکل اصلی در \eng{SNN} این است که خروجی نورون یک تابع پله‌ای است:
\[
S =
\begin{cases}
1 & V > \theta \\
0 & V \leq \theta
\end{cases}
\]
تابع پله‌ای تقریباً همه‌جا مشتق صفر دارد. بنابراین اگر بخواهیم \eng{backpropagation} انجام دهیم، مشتق لازم برای به‌روزرسانی وزن‌ها از مسیر خروجی اسپایک به دست نمی‌آید. پرسش اصلی این است: چگونه می‌توان در شبکهٔ اسپایکی \eng{backpropagation} انجام داد، وقتی تابع اسپایک مشتق‌پذیر نیست؟

\subsubsection{مسئلهٔ مشتق‌پذیری اسپایک}
راه‌حل رایج، استفاده از \eng{Surrogate Gradient} است. ایده این است که در فاز \eng{forward}، شبکه همچنان اسپایکی باقی بماند و خروجی نورون واقعاً صفر یا یک باشد؛ اما در فاز \eng{backward}، برای محاسبهٔ گرادیان، تابع پله‌ای با تابعی نرم‌تر و مشتق‌پذیر جایگزین شود.

\subsubsection{تابع جایگزین}
تابع جایگزین یا \eng{surrogate function} می‌تواند شکلی شبیه \eng{sigmoid} یا تابع نرم دیگر داشته باشد. نقش آن این نیست که رفتار واقعی نورون را در فاز پیش‌رو تغییر دهد؛ نقش آن فقط فراهم کردن مشتقی است که در \eng{backpropagation} بتوان از آن استفاده کرد.

\coursefigure{0.65}{surrogate-gradient-functions.png}{نمونه‌هایی از مشتق‌های جایگزین برای تابع شلیک؛ تقریب‌هایی مانند \eng{SuperSpike} و \eng{SLAYER} در اطراف آستانه گرادیان قابل استفاده فراهم می‌کنند.}

پس در این چارچوب:
\begin{itemize}
  \item در \eng{forward}، تابع واقعی اسپایک استفاده می‌شود؛
  \item در \eng{backward}، از تابع جایگزین مشتق‌پذیر برای محاسبهٔ گرادیان استفاده می‌شود.
\end{itemize}

\subsubsection{تفاوت با \eng{Activation Function} معمولی}
در شبکه‌های عصبی معمولی، \eng{activation function} در خود \eng{forward pass} استفاده می‌شود. یعنی خروجی هر لایه واقعاً از تابعی مانند \eng{sigmoid}، \eng{ReLU} یا \eng{tanh} عبور می‌کند.

اما در \eng{SNN} با \eng{surrogate gradient}، تابع جایگزین فقط برای \eng{backpropagation} استفاده می‌شود. در فاز \eng{forward} همچنان اسپایک واقعی داریم، ورودی و خروجی همچنان صفر و یک هستند، و تابع \eng{surrogate} فقط برای محاسبهٔ مشتق و اصلاح وزن‌ها وارد می‌شود. بنابراین \eng{surrogate function} همان \eng{activation function} معمولی نیست، چون نقش اصلی آن در فاز \eng{backward} است، نه در فاز \eng{forward}.

\subsection{\eng{Feedforward} در \eng{Deep SNN}}
در فاز \eng{feedforward}، اسپایک‌های لایهٔ قبلی به لایهٔ بعدی منتقل می‌شوند. اگر خروجی لایهٔ قبلی را \(S\) و وزن سیناپسی میان دو لایه را \(W\) بگیریم، جریان ورودی به نورون بعدی در ساده‌ترین حالت چنین است:
\[
I = SW
\]
در حالت زمانی، جریان می‌تواند اثر مقدار قبلی خود را نیز حفظ کند:
\[
I[t+1] = \alpha I[t] + S[t]W
\]
در این رابطه، \(I[t]\) جریان قبلی، \(\alpha\) ضریب کاهش یا \eng{decay} جریان، \(S[t]\) اسپایک ورودی و \(W\) وزن سیناپسی است.

ولتاژ نورون نیز با توجه به ولتاژ قبلی، ضریب نشت و جریان ورودی به‌روزرسانی می‌شود:
\[
V[t+1] = \beta V[t] + I[t+1]
\]
که در آن، \(V[t]\) ولتاژ قبلی، \(\beta\) ضریب نشت ولتاژ و \(I[t+1]\) جریان ورودی جدید است.

اگر ولتاژ از آستانه بیشتر شود، اسپایک تولید می‌شود:
\[
S[t+1] =
\begin{cases}
1 & V[t+1] > \theta \\
0 & V[t+1] \leq \theta
\end{cases}
\]
پس از اسپایک، نورون باید ریست شود؛ یعنی ولتاژ و گاهی جریان آن کاهش یابد یا صفر شود.

\subsubsection{نقش \eng{reset}}
وقتی نورون اسپایک می‌زند، فقط به لایهٔ بعدی سیگنال نمی‌فرستد؛ بلکه خودش نیز ریست می‌شود. بنابراین اسپایک دو اثر دارد: به لایهٔ بعدی منتقل می‌شود و در وزن سیناپسی ضرب می‌شود، و هم‌زمان روی نورونی که اسپایک زده اثر می‌گذارد و ولتاژ آن را ریست می‌کند. به همین دلیل، در \eng{SNN} زمان و وضعیت داخلی نورون اهمیت اساسی دارند.

\subsection{\eng{Backpropagation} با \eng{Surrogate Gradient}}
در انتهای شبکه، خروجی با هدف مقایسه می‌شود و خطا محاسبه می‌شود. برای نمونه، اگر از \eng{Mean Squared Error} استفاده کنیم:
\[
E = \frac{1}{2}(T - O)^2
\]
که در آن \(T\) هدف و \(O\) خروجی شبکه است.

برای به‌روزرسانی وزن‌ها، مانند شبکه‌های معمولی از گرادیان استفاده می‌کنیم:
\[
W_{\text{\lr{new}}} =
W_{\text{\lr{old}}} - \eta \frac{\partial E}{\partial W}
\]
در این رابطه، \(\eta\) نرخ یادگیری و \(\frac{\partial E}{\partial W}\) مشتق خطا نسبت به وزن است.

با استفاده از قاعدهٔ زنجیره‌ای:
\[
\frac{\partial E}{\partial W}
=
\frac{\partial E}{\partial O}
\cdot
\frac{\partial O}{\partial Net}
\cdot
\frac{\partial Net}{\partial W}
\]
اگر:
\[
Net = XW
\]
آنگاه:
\[
\frac{\partial Net}{\partial W} = X
\]
و برای خطای مربعی:
\[
\frac{\partial E}{\partial O} = O - T
\]
بخش اصلی \eng{surrogate gradient} در جملهٔ میانی است:
\[
\frac{\partial O}{\partial Net}
\]
زیرا اگر \(O\) را خروجی تابع پله‌ای واقعی بگیریم، مشتق آن صفر می‌شود. بنابراین این مشتق را از تابع \eng{surrogate} می‌گیریم.

\subsubsection{مثال عددی}
فرض کنید:
\[
X = 2,\qquad W = 0.4,\qquad \theta = 1,\qquad T = 1,\qquad V_0 = 0
\]
ابتدا مقدار \(Net\) را حساب می‌کنیم:
\[
Net = XW = 2 \times 0.4 = 0.8
\]
چون:
\[
0.8 < 1
\]
نورون اسپایک نمی‌زند و بنابراین:
\[
O = 0
\]
اما هدف برابر ۱ بوده است:
\[
T = 1
\]
پس خروجی شبکه اشتباه است و وزن باید افزایش یابد.

مشتق خطا نسبت به خروجی:
\[
\frac{\partial E}{\partial O} = O - T = 0 - 1 = -1
\]
فرض شد مقدار مشتق تابع \eng{surrogate} در نقطهٔ \(Net=0.8\) تقریباً \(0.25\) است. همچنین:
\[
\frac{\partial Net}{\partial W} = X = 2
\]
پس:
\[
\frac{\partial E}{\partial W}
=
(-1)(0.25)(2)
=
-0.5
\]
اگر نرخ یادگیری \(\eta = 0.1\) باشد:
\[
W_{\text{\lr{new}}}
=
0.4 - 0.1(-0.5)
=
0.45
\]
بنابراین وزن از \(0.4\) به \(0.45\) افزایش پیدا می‌کند. این نتیجه منطقی است، چون خروجی باید ۱ می‌شد اما نورون به آستانه نرسیده بود؛ در نتیجه وزن افزایش می‌یابد تا دفعهٔ بعد ولتاژ راحت‌تر از آستانه عبور کند.

\subsection{محلی بودن، آموزش سراسری و ارزش عملی \eng{SNN}}
استفاده از \eng{surrogate gradient} در مرحلهٔ آموزش تا حدی ویژگی محلی بودن \eng{SNN} را کاهش می‌دهد، چون خطا به‌صورت سراسری محاسبه می‌شود و \eng{backpropagation} انجام می‌گیرد. اما نکتهٔ مهم این است که \eng{surrogate gradient} فقط در مرحلهٔ آموزش استفاده می‌شود.

پس از آموزش، شبکه همچنان اسپایکی است: ورودی و خروجی صفر و یک باقی می‌مانند، مدل می‌تواند روی سخت‌افزار \eng{neuromorphic} اجرا شود، و دیگر نیازی به \eng{surrogate gradient} یا \eng{backpropagation} در فاز اجرا نیست. وزن‌ها پیش‌تر \eng{fine-tune} شده‌اند.

بنابراین چارچوب اصلی \eng{SNN} حفظ می‌شود. فقط برای اینکه بتوانیم وزن‌ها را در شبکهٔ عمیق و نظارت‌شده تنظیم کنیم، در مرحلهٔ آموزش از \eng{surrogate gradient} کمک می‌گیریم.

ارزش عملی \eng{SNN} بیشتر در فاز اجرا روشن می‌شود. در کاربردهایی مانند ربات، پهپاد، سامانه‌های \eng{embedded}، سخت‌افزار \eng{neuromorphic} و پردازش \eng{real-time} با مصرف انرژی کم، شبکهٔ آموزش‌دیده می‌تواند به‌صورت اسپایکی اجرا شود. در این حالت، مهم نیست وزن‌ها دقیقاً با چه روشی آموزش دیده‌اند؛ مهم این است که \eng{inference} با اسپایک‌های صفر و یک انجام شود.

\subsection{کاربرد \eng{STDP} در برابر \eng{Surrogate Gradient}}
\eng{STDP} مناسب‌تر است وقتی شبکه کوچک یا کم‌عمق است، یادگیری بدون ناظر است، یادگیری \eng{online} لازم داریم، یا هدف، پیاده‌سازی زیستی‌تر است. در مقابل، \eng{surrogate gradient} برای \eng{Deep SNN} و وظایف نظارت‌شده مناسب‌تر است، جایی که باید خطای خروجی با هدف مقایسه شود و وزن‌ها در چندین لایه تنظیم شوند.

\chapter{مدل‌های اسپایکی مدرن}

\section{\eng{Spiking Transformer} و \eng{SpikeFormer}}
در آغاز این بخش، مسیر قبلی بحث جمع‌بندی می‌شود. بحث سیستم پاداش تا حد لازم پیش رفته بود و تمرکز بحث به شبکه‌های اسپایکی عمیق منتقل شد. تا اینجا چند نکته روشن شده بود: \eng{SNN}ها از نظر ورودی رویدادمحور و مصرف انرژی به مغز نزدیک‌ترند، ورودی آن‌ها معمولاً به‌صورت اسپایک‌های صفر و یک است، داده در آن‌ها ماهیت زمانی دارد، و در مقایسه با شبکه‌های عمیق معمولی می‌توانند از نظر انرژی کم‌هزینه‌تر باشند.

در \eng{Deep SNN} نیز اگر یادگیری نظارت‌شده داشته باشیم، شبکه برچسب را می‌بیند، خروجی خود را با هدف مقایسه می‌کند و بر اساس خطا وزن‌ها را به‌روزرسانی می‌کند. مسئلهٔ اصلی همان گسسته بودن خروجی‌های اسپایکی است؛ به همین دلیل برای آموزش، از ایدهٔ \eng{surrogate gradient} استفاده می‌شود. اگر در یک \eng{Deep SNN} از \eng{linear layer} صحبت می‌کنیم، منظور این است که وزن‌ها در اسپایک‌ها ضرب می‌شوند و خروجی اسپایکی تازه‌ای برای لایهٔ بعدی ساخته می‌شود:
\[
W \times \text{\lr{spike}}
\]

\subsection{چرا \eng{Transformer} در \eng{SNN}؟}
وقتی بتوانیم معماری‌های عمیق اسپایکی بسازیم، پرسش طبیعی این است که آیا می‌توان ایدهٔ \eng{Transformer} را نیز وارد \eng{SNN} کرد. در شبکه‌های عمیق معمولی، \eng{Transformer}ها در بسیاری از حوزه‌ها معماری‌های بسیار موفقی هستند. بنابراین اگر بخواهیم یک \eng{SNN} قدرتمندتر داشته باشیم، یکی از مسیرها ترکیب \eng{Transformer} و \eng{Spiking Neural Network} است؛ مسیری که در مدل‌هایی مانند \eng{SpikeFormer} دنبال می‌شود.

برای ساختن \eng{Transformer} اسپایکی، سه بخش اصلی باید روشن شود:
\begin{enumerate}
  \item کدگذاری دادهٔ ورودی؛
  \item \eng{Attention}؛
  \item بلوک \eng{MLP}.
\end{enumerate}
اگر این سه بخش مشخص شوند، می‌توان گفت معماری‌ای شبیه \eng{Transformer} اما مبتنی بر \eng{SNN} داریم.

\subsection{مسئلهٔ کدگذاری ورودی}
در شبکه‌های معمولی، داده را می‌توان مستقیم به شبکه داد؛ اما در \eng{SNN} داده باید به دنباله‌ای از اسپایک‌ها تبدیل شود. اگر دادهٔ اولیه شامل عددهای پیوسته یا \eng{floating-point} باشد، باید این مقدارها به صفر و یک تبدیل شوند و این تبدیل در چند گام زمانی رخ دهد.

پس کدگذاری ورودی دو کار انجام می‌دهد: مقدارهای پیوسته را به اسپایک‌های صفر و یک تبدیل می‌کند، و آن‌ها را در چند گام زمانی پخش می‌کند. برای مثال، ممکن است داده را در چهار گام زمانی نمایش دهیم:
\[
T = 4
\]
در این حالت، هر مقدار عددی باید به یک دنبالهٔ اسپایکی در طول زمان تبدیل شود.

\subsubsection{کدگذاری تصویر}
فرض کنید ورودی یک تصویر است و مقدار پیکسل‌ها در بازهٔ ۰ تا ۲۵۵ قرار دارند. نخستین گام معمولاً نرمال‌سازی است:
\[
x_{\text{\lr{norm}}} = \frac{x}{255}
\]
با این کار مقدار هر پیکسل به بازهٔ صفر تا یک می‌رود.

برای نمونه، اگر چند مقدار نرمال‌شده داشته باشیم:
\[
0.2,\quad 0.8,\quad 0.5,\quad 0.5
\]
و بخواهیم آن‌ها را در چهار گام زمانی پخش کنیم، هر مقدار را در ۴ ضرب می‌کنیم:
\[
0.2 \times 4 = 0.8,\qquad
0.8 \times 4 = 3.2,\qquad
0.5 \times 4 = 2,\qquad
0.5 \times 4 = 2
\]
سپس مقدارها را گرد می‌کنیم:
\[
1,\quad 3,\quad 2,\quad 2
\]
یعنی مقدار اول در کل زمان یک اسپایک تولید می‌کند، مقدار دوم سه اسپایک، مقدار سوم دو اسپایک و مقدار چهارم نیز دو اسپایک.

به این ترتیب، دادهٔ ثابت اولیه در طول چند گام زمانی پخش می‌شود و به دنباله‌ای اسپایکی تبدیل می‌گردد. این روند شبیه نوعی «برعکس \eng{integration}» است: عدد پیوسته به جای اینکه مستقیم نگه داشته شود، در طول زمان باز می‌شود.

\subsubsection{نکتهٔ محاسباتی در کدگذاری}
پشت بسیاری از روش‌های کدگذاری، لزوماً معنای زیستی یا معنایی بسیار دقیقی وجود ندارد. این روش‌ها بیشتر مدل‌های محاسباتی‌اند. هدف اصلی این است که دادهٔ پیوسته به اسپایک تبدیل شود، نرخ اسپایک با مقدار عدد اولیه رابطهٔ معنادار داشته باشد، و مدل در نهایت از نظر دقت و مصرف انرژی عملکرد قابل قبولی پیدا کند.

در شبکه‌های اسپایکی معمولاً یک مبادله وجود دارد:
\[
\text{\lr{Energy / Cost}} \quad \leftrightarrow \quad \text{\lr{Accuracy}}
\]
اگر از \eng{Transformer} معمولی استفاده کنیم، معمولاً دقت بیشتر است، اما مصرف انرژی و هزینهٔ محاسباتی نیز بیشتر می‌شود. اگر از \eng{SNN} استفاده کنیم، مصرف انرژی کمتر است، اما ممکن است دقت کاهش یابد. به همین دلیل، در مقاله‌های مربوط به \eng{SpikeFormer} معمولاً مدل را با \eng{SNN}های دیگر مقایسه می‌کنند، نه لزوماً با \eng{Transformer}های معمولی؛ چون از همان مرحلهٔ کدگذاری بخشی از اطلاعات پیوستهٔ داده از دست می‌رود.

\subsubsection{کدگذاری متن}
در متن نیز پس از ساخت \eng{embedding}، با عددهای پیوسته سروکار داریم. برای نمونه، ممکن است یکی از ویژگی‌ها مقدار زیر را داشته باشد:
\[
x = 0.8
\]
این مقدار باید به دنباله‌ای اسپایکی تبدیل شود. یکی از ایده‌ها \eng{Rate Encoding} است؛ یعنی مقدار عددی به نرخ اسپایک ربط داده می‌شود.

اگر:
\[
x = 0.8,\qquad T = 4
\]
باشد، انتظار داریم تقریباً سه اسپایک داشته باشیم، چون:
\[
0.8 \times 4 \approx 3
\]
پس یک نمایش ممکن چنین است:
\[
[1,1,1,0]
\]
یا هر چینش دیگری که سه اسپایک در چهار گام زمانی داشته باشد.

\subsubsection{کدگذاری با ایدهٔ \eng{Integration}}
روش دیگر این است که عدد در هر گام زمانی به پتانسیل غشا وارد شود و بررسی کنیم چه زمانی از آستانه عبور می‌کند. فرض کنید:
\[
x = 0.8,\qquad \theta = 1
\]
در گام اول:
\[
0.8 < 1
\]
پس اسپایک نداریم. در گام دوم:
\[
0.8 + 0.8 = 1.6
\]
از آستانه عبور می‌کنیم، پس اسپایک داریم. پس از اسپایک، مقدار آستانه از پتانسیل کم می‌شود و باقی‌مانده \(0.6\) می‌ماند.

در گام بعد:
\[
0.6 + 0.8 = 1.4
\]
دوباره اسپایک داریم و پس از کم کردن ۱، باقی‌مانده \(0.4\) می‌شود. در گام بعد:
\[
0.4 + 0.8 = 1.2
\]
باز هم اسپایک داریم. پس خروجی می‌تواند چنین باشد:
\[
[0,1,1,1]
\]
در این روش نیز نرخ اسپایک معنای اصلی را حفظ می‌کند: عدد بزرگ‌تر، اسپایک بیشتری تولید می‌کند. حداقل معیار برای کدگذاری معنادار این است که نرخ اسپایک با مقدار عدد پیوسته رابطهٔ مستقیم داشته باشد.

\subsection{\eng{Attention} در شبکهٔ اسپایکی}
در \eng{Transformer} معمولی، هدف \eng{attention} این است که بخش‌های مرتبط با هم در یک \eng{context} به هم نزدیک‌تر شوند. برای مثال، اگر جمله‌ای دربارهٔ «شیر» و «جنگل» داشته باشیم، \eng{attention} باید کمک کند بازنمایی کلمهٔ «شیر» بیشتر به معنای حیوانی آن نزدیک شود، نه معنای نوشیدنی. در تصویر نیز \eng{attention} باعث می‌شود بخش‌های مهم تصویر، با توجه به زمینه، وزن بیشتری بگیرند.

در \eng{SNN} نیز همین هدف را داریم، اما ورودی‌ها دنباله‌های اسپایکی‌اند. بنابراین انتظار داریم در خروجی \eng{attention}، نرخ اسپایک برای بخش‌هایی که به هم مرتبط‌ترند بیشتر شود.

\subsubsection{\eng{Query}، \eng{Key} و \eng{Value} اسپایکی}
در \eng{Transformer} معمولی، \eng{embedding} ورودی به سه بردار تبدیل می‌شود:
\[
Q,\quad K,\quad V
\]
در \eng{SNN} نیز همین ایده وجود دارد، اما ورودی یک دنبالهٔ اسپایکی است. برای این کار سه ماتریس وزن داریم:
\[
W_Q,\quad W_K,\quad W_V
\]
دنبالهٔ اسپایکی ورودی در این وزن‌ها ضرب می‌شود و سه دنبالهٔ جدید \(Q\)، \(K\) و \(V\) ساخته می‌شود. این وزن‌ها در طول آموزش، با همان ایدهٔ \eng{surrogate gradient} به‌روزرسانی می‌شوند.

\subsubsection{حذف \eng{Softmax}}
در \eng{Transformer} معمولی، نمرهٔ توجه از ضرب زیر به دست می‌آید:
\[
QK^T
\]
سپس معمولاً نرمال‌سازی و \eng{softmax} انجام می‌شود:
\[
\operatorname{softmax}\left(\frac{QK^T}{\sqrt{d}}\right)
\]
و نتیجه در \(V\) ضرب می‌شود. هدف \eng{softmax} این است که نمره‌های خام به نوعی توزیع وزن تبدیل شوند؛ یعنی وزن‌ها نرمال شوند و مجموعشان ۱ شود.

در \eng{SNN}، داده‌ها پیوسته نیستند، بلکه اسپایک‌های صفر و یک‌اند. وقتی اسپایک‌ها در هم ضرب می‌شوند، خود این ضرب تا حدی نقش شباهت را بازی می‌کند: اگر هر دو اسپایک فعال باشند، خروجی ۱ می‌شود و اگر یکی یا هر دو فعال نباشند، خروجی ۰ می‌شود. بنابراین \eng{softmax} همان معنای قبلی را ندارد، چون داده از فضای پیوسته وارد فضای زمانی اسپایک‌ها شده است.

در \eng{Spiking Self-Attention} می‌توان ساختاری شبیه زیر داشت:
\[
QK^TV
\]
بدون اینکه الزاماً \eng{softmax} اعمال شود. البته مقیاس‌گذاری مشابه \eng{Transformer} معمولی، مانند تقسیم بر \(\sqrt{d}\)، می‌تواند همچنان در معماری وجود داشته باشد.

\subsubsection{تغییر ترتیب ضرب‌ها و کاهش پیچیدگی}
وقتی \eng{softmax} حذف شود، می‌توان ترتیب ضرب‌ها را تغییر داد. در \eng{Transformer} معمولی معمولاً ابتدا \(QK^T\) محاسبه می‌شود. این کار ماتریسی با ابعاد وابسته به تعداد توکن‌ها می‌سازد و پیچیدگی بالایی دارد.

اما اگر \eng{softmax} وجود نداشته باشد، می‌توان ابتدا \(K^TV\) را محاسبه کرد و سپس نتیجه را در \(Q\) ضرب کرد:
\[
Q(K^TV)
\]
این تغییر می‌تواند پیچیدگی محاسباتی را کاهش دهد. در مقاله‌های \eng{SpikeFormer} نیز یکی از ادعاهای مهم همین کاهش مصرف انرژی و هزینهٔ محاسباتی است. چون مدل اسپایکی معمولاً از نظر دقت به \eng{Transformer} معمولی نمی‌رسد، باید مزیت خود را در انرژی و پیچیدگی نشان دهد.

\subsubsection{تفسیر \eng{Attention} اسپایکی}
در \eng{attention} اسپایکی، وقتی توکنی مانند «شیر» به توکن دیگری مانند «جنگل» توجه می‌کند، معنایش این است که اسپایک‌های مربوط به «جنگل» اجازه پیدا می‌کنند روی پتانسیل یا بازنمایی مربوط به «شیر» اثر بگذارند. اگر نرخ اسپایک میان این دو بالا باشد، یعنی ارتباط بیشتری بین آن‌ها وجود دارد.

در نهایت، اگر وظیفه مثلاً \eng{next-token generation} باشد، انتظار داریم \eng{logit} مربوط به معنای درست افزایش پیدا کند. در زمینهٔ «جنگل»، \eng{logit} مربوط به معنای حیوانی «شیر» باید از معنای نوشیدنی بیشتر شود.

\subsection{معماری کلی \eng{SpikeFormer}}
فرآیند کلی \eng{SpikeFormer} چنین است: دادهٔ ورودی کدگذاری می‌شود، به دنبالهٔ اسپایکی تبدیل می‌شود، وارد \eng{Spiking Self-Attention} یا \eng{SSA} می‌شود، خروجی توجه وارد بلوک \eng{MLP} می‌شود، و در اطراف بلوک‌های توجه و \eng{MLP} اتصال‌های \eng{residual} قرار می‌گیرند.

\coursefigure{0.95}{spikeformer-architecture.png}{نمای کلی \eng{SpikeFormer}: تصویر ورودی به توالی‌های اسپایکی تبدیل می‌شود، از بلوک‌های \eng{SSA} و \eng{MLP} عبور می‌کند و در انتها با سر طبقه‌بندی تصمیم گرفته می‌شود.}

در بخش توجه، ورودی اسپایکی از \eng{linear layer} عبور می‌کند و معمولاً پس از آن یک نورون \eng{LIF} قرار می‌گیرد. خروجی \eng{LIF} دوباره یک دنبالهٔ اسپایکی است. برای ساختن \(Q\)، \(K\) و \(V\)، سه وزن جداگانه داریم:
\[
W_Q,\quad W_K,\quad W_V
\]
و سپس دنباله‌های \(Q\)، \(K\) و \(V\) با ضرب‌های ماتریسی یا ضرب داخلی ترکیب می‌شوند.

\subsection{نقش \eng{LIF}}
در \eng{SNN}، پس از بسیاری از \eng{linear layer}ها یک نورون \eng{LIF} یا \eng{Leaky Integrate-and-Fire} قرار می‌گیرد. نقش \eng{LIF} این است که پتانسیل غشا را در طول زمان جمع کند، نشتی داشته باشد، اگر پتانسیل از آستانه عبور کرد اسپایک تولید کند، پس از اسپایک ریست شود، مدتی در دورهٔ \eng{refractory} بماند، و سپس دوباره شروع به جمع‌کردن ورودی کند. بنابراین خروجی \eng{LIF} دوباره یک دنبالهٔ اسپایکی است.

\subsection{\eng{MLP} در \eng{Transformer} اسپایکی}
در \eng{Transformer} معمولی، \eng{MLP} معمولاً شامل دو \eng{linear layer} و یک تابع فعال‌سازی میان آن‌هاست:
\[
\text{\lr{Linear}} \rightarrow \text{\lr{Activation}} \rightarrow \text{\lr{Linear}}
\]
تابع فعال‌سازی می‌تواند \eng{ReLU}، \eng{GELU}، \eng{sigmoid} یا تابعی مشابه باشد. نقش آن ایجاد غیرخطی‌بودن میان دو لایهٔ خطی است تا مدل بتواند بازنمایی عمیق‌تر و انتزاعی‌تری از هر توکن بسازد.

در \eng{SNN} نمی‌توان همان تابع‌های فعال‌سازی را به همان شکل روی اسپایک‌های صفر و یک استفاده کرد. در معماری اسپایکی، خود \eng{LIF} نقش غیرخطی‌سازی را بازی می‌کند:
\[
\text{\lr{Linear}} \rightarrow \text{\lr{LIF}} \rightarrow
\text{\lr{Linear}} \rightarrow \text{\lr{LIF}}
\]
پس \eng{LIF} معادل دقیق \eng{activation function} معمولی نیست، اما نقش اصلی آن، یعنی ایجاد \eng{non-linearity} و تبدیل خروجی به دنبالهٔ اسپایکی، را انجام می‌دهد.

\subsection{\eng{Residual Connection}}
در \eng{SpikeFormer} نیز مانند \eng{Transformer} معمولی، اطراف بلوک‌های مهم اتصال \eng{residual} داریم؛ هم اطراف بلوک \eng{attention} و هم اطراف بلوک \eng{MLP}. دلیل استفاده از \eng{residual connection} کمک به کاهش مشکل \eng{vanishing gradient}، حفظ اطلاعات قبلی، و جلوگیری از نابودی بازنمایی پیشین است.

اگر یک بلوک چیز تازه‌ای یاد نگیرد، دست‌کم اطلاعات قبلی همچنان از مسیر \eng{residual} باقی می‌ماند. این اتصال مانند یک پشتوانه عمل می‌کند: اگر مدل در یک بلوک موفق به ساخت بازنمایی بهتر نشود، بازنمایی قبلی کاملاً از دست نمی‌رود.

\subsection{جمع‌بندی \eng{SpikeFormer}}
در این بخش، سه بخش اصلی \eng{Transformer} اسپایکی بررسی شد. نخست، کدگذاری: دادهٔ پیوسته باید به دنبالهٔ اسپایکی تبدیل شود و نرخ اسپایک با مقدار اولیهٔ داده رابطهٔ مستقیم داشته باشد. روش‌هایی مانند \eng{Rate Encoding} و روش‌های مبتنی بر \eng{integration} برای این هدف استفاده می‌شوند.

دوم، توجه: در \eng{attention} اسپایکی همچنان ایدهٔ \(Q\)، \(K\) و \(V\) وجود دارد، اما به دلیل صفر و یکی بودن اسپایک‌ها و زمانی بودن داده، \eng{softmax} ضرورت قبلی را ندارد. حذف \eng{softmax} معماری را ساده‌تر می‌کند، پیچیدگی محاسباتی را کاهش می‌دهد، اجازهٔ تغییر ترتیب ضرب‌ها را می‌دهد و می‌تواند مصرف انرژی را کم کند.

\coursefigure{0.95}{classic-vs-spiking-attention.png}{مقایسهٔ توجه کلاسیک و توجه اسپایکی: حذف \eng{Softmax} و استفاده از اسپایک‌ها هزینهٔ محاسباتی و انرژی \eng{Spiking Self-Attention} را کاهش می‌دهد.}

سوم، \eng{MLP}: در \eng{MLP} معمولی از تابع‌هایی مانند \eng{GELU} یا \eng{ReLU} استفاده می‌شود. در \eng{SNN}، \eng{LIF} این نقش را بر عهده می‌گیرد، چون رفتار غیرخطی ایجاد می‌کند و خروجی را دوباره به دنباله‌ای اسپایکی تبدیل می‌کند.

\section{\eng{SpikeBERT} و تقطیر دانش}

بخش قبل نشان داد که پس از \eng{Deep SNN} می‌توان ایدهٔ \eng{Transformer} را نیز به فضای اسپایکی برد. \eng{SpikeFormer} نمونه‌ای از همین مسیر بود و بیشتر در حوزهٔ بینایی ماشین مطرح شد. در این بخش پرسش اصلی این است که آیا می‌توان همین ایده را برای متن و پردازش زبان طبیعی نیز به کار برد یا نه.

این دسته از کارها، از جمله \eng{SpikeFormer} و \eng{SpikeBERT}، بیشتر نقش \eng{Proof of Concept} دارند: نشان می‌دهند تبدیل معماری‌های \eng{ANN} به \eng{SNN} ممکن است، اما هنوز چالش‌های جدی باقی است. در \eng{SpikeFormer} سه ایدهٔ اصلی دیده شد: در بخش \eng{Attention}، ضرب ماتریسی و \eng{softmax} تا حدی با عملیات ساده‌تر در فضای اسپایکی جایگزین می‌شود؛ در بخش \eng{MLP}، نورون‌هایی مانند \eng{LIF} نقش غیرخطی‌سازی را بر عهده می‌گیرند؛ و برای آموزش، از \eng{Surrogate Gradient} استفاده می‌شود.

\subsection{چرا متن برای مدل اسپایکی دشوار است؟}
وقتی ایدهٔ \eng{SpikeFormer} از تصویر به متن منتقل می‌شود، مدل به‌سادگی همگرا نمی‌شود. دلیل اصلی این مشکل تفاوت ماهوی متن و تصویر است. متن هم از نظر تراکم اطلاعات، هم از نظر وابستگی‌های بلندمدت، و هم از نظر نمادین‌بودن، برای یک مدل اسپایکی دشوارتر از تصویر است.

\subsubsection{تراکم اطلاعات در متن}
در متن، تراکم اطلاعات بسیار بالاست. گاهی یک تغییر کوچک می‌تواند معنای کل جمله را عوض کند. برای نمونه:

\begin{itemize}
  \item «من رفتم»
  \item «من نرفتم»
\end{itemize}

اضافه‌شدن یک نشانهٔ نفی، معنای جمله را تقریباً برعکس می‌کند. در تصویر معمولاً چنین نیست؛ چند پیکسل کوچک، در بسیاری از موقعیت‌ها، معنای کل تصویر را به‌طور کامل عوض نمی‌کنند، چون تصویر افزونگی بیشتری دارد.

در \eng{SNN}، نمایش‌های پیوسته به دنباله‌هایی از صفر و یک تبدیل می‌شوند. اگر این تبدیل ظرافت‌های معنایی متن را از بین ببرد، شبکه نمی‌تواند به‌خوبی یاد بگیرد. به همین دلیل، اسپایکی‌کردن متن نسبت به اسپایکی‌کردن تصویر حساس‌تر است.

\subsubsection{وابستگی‌های بلندمدت}
در زبان، معنای یک عبارت ممکن است به جمله‌ها یا کلمه‌هایی وابسته باشد که خیلی زودتر آمده‌اند. گاهی برای فهم یک جمله باید کل \eng{context} پاراگراف در نظر گرفته شود.

نورون \eng{LIF} ذاتاً نشتی دارد. نشتی یعنی پتانسیل غشا به‌تدریج کم می‌شود و این رفتار نوعی فراموشی ایجاد می‌کند. بنابراین یک شبکهٔ اسپایکی ممکن است در نگهداری وابستگی‌های بلندمدت متن دچار مشکل شود. در تصویر، وابستگی‌های فضایی معمولاً چنین شکل زبانی و بلندمدتی ندارند و همین مسئله انتقال مستقیم ایدهٔ \eng{SpikeFormer} به متن را دشوار می‌کند.

\subsubsection{ماهیت نمادین و انتزاعی زبان}
متن از نمادهای گسسته و انتزاعی ساخته شده است. همین که این نمادها به \eng{embedding} پیوسته تبدیل شوند، خود یک مسئلهٔ اصلی در \eng{NLP} است. اگر پس از آن بخواهیم همین \eng{embedding}های پیوسته را نیز اسپایکی کنیم، مسیر یادگیری پیچیده‌تر می‌شود.

هدف از \eng{SpikeBERT} این نیست که حتماً از همان ابتدا به دقت مدل‌های زبانی معمولی برسد. هدف این است که مدل اسپایکی بتواند از مزیت‌هایی مانند ماهیت زمانی، مصرف انرژی کمتر و سرعت بالقوهٔ بیشتر استفاده کند، در حالی که هنوز از پس وظایف اصلی \eng{NLP} بربیاید.

\subsection{\eng{Knowledge Distillation}}
ایدهٔ اصلی \eng{SpikeBERT} استفاده از \eng{Knowledge Distillation} است. اگر \eng{Student} فقط در انتهای شبکه و فقط بر اساس برچسب نهایی آموزش ببیند، روی متن به‌خوبی همگرا نمی‌شود. بنابراین باید در بخش‌های مختلف شبکه به آن کمک کرد.

در این روش، مدل معمولی \eng{BERT} نقش \eng{Teacher} را دارد و مدل اسپایکی \eng{SpikeBERT} نقش \eng{Student} را. \eng{Student} نه‌تنها با برچسب نهایی، بلکه با خروجی‌های میانی \eng{Teacher} نیز مقایسه می‌شود. به این ترتیب، یادگیری فقط به یک \eng{loss} نهایی محدود نمی‌ماند.

در \eng{SpikeBERT} چند نوع خطا در کنار هم استفاده می‌شود:

\begin{itemize}
  \item \eng{Embedding Loss}
  \item \eng{Hidden Feature Loss}
  \item \eng{Logit Loss}
  \item \eng{Cross Entropy Loss} روی برچسب واقعی
\end{itemize}

پس \eng{loss} نهایی ترکیبی از چند \eng{loss} است و برای هر بخش می‌توان ضریب وزنی جداگانه در نظر گرفت.

\subsection{\eng{Teacher} و \eng{Student}}
\eng{Teacher} همان \eng{BERT} معمولی است؛ مدلی که از قبل روی متن‌های خام آموزش دیده و نمایش‌های پیوستهٔ زبانی تولید می‌کند. این نمایش‌ها حامل رابطه‌های معنایی‌اند. برای مثال، در فضای برداری زبان، رابطهٔ «تهران» با «ایران» می‌تواند شبیه رابطهٔ «پاریس» با «فرانسه» باشد.

\eng{Student} همان \eng{SpikeBERT} است. ورودی اولیهٔ آن می‌تواند \eng{embedding}های \eng{BERT} باشد، اما این \eng{embedding}ها باید به فضای اسپایکی برده شوند.

اگر ورودی \eng{embedding} را به شکل زیر بنویسیم:

\[
X \in \mathbb{R}^{N \times D}
\]

در این رابطه، \(N\) تعداد \eng{token}ها و \(D\) تعداد ویژگی‌های هر \eng{token} است. برای نمونه، نمایش یک \eng{token} می‌تواند چنین باشد:

\[
[0.3,\ 0.7,\ 0]
\]

این نمایش پیوسته باید به دنباله‌ای از اسپایک‌ها تبدیل شود تا بتواند وارد مدل اسپایکی شود.

\subsection{تبدیل \eng{Embedding} پیوسته به اسپایک}
برای اسپایکی‌کردن \eng{embedding}، یک ابرپارامتر مهم وارد می‌شود:

\[
T = \text{\eng{Time Step}}
\]

اگر \(T=5\) باشد، هر مقدار پیوسته در طول پنج گام زمانی بازنمایی می‌شود. برای مثال، اگر مقدار ورودی \(0.3\) باشد و آستانهٔ نورون \eng{LIF} برابر \(1\) در نظر گرفته شود، پتانسیل می‌تواند به این شکل جمع شود:

\begin{itemize}
  \item گام اول: \(0.3\)
  \item گام دوم: \(0.6\)
  \item گام سوم: \(0.9\)
  \item گام چهارم: \(1.2\)، تولید اسپایک
\end{itemize}

پس از تولید اسپایک، پتانسیل reset یا کاهش داده می‌شود. به این ترتیب، مقدار پیوستهٔ \(0.3\) در طول زمان به یک الگوی اسپایکی تبدیل می‌شود.

در نتیجه، شکل نمایش از حالت زیر:

\[
N \times D
\]

به حالت زیر تبدیل می‌شود:

\[
T \times N \times D
\]

یعنی همان نمایش زبانی در طول زمان کشیده می‌شود.

\subsection{بازگرداندن اسپایک‌ها به نمایش پیوسته}
چون خروجی‌های \eng{Teacher} پیوسته‌اند، برای مقایسهٔ \eng{Student} با \eng{Teacher} باید خروجی‌های اسپایکی دوباره به یک نمایش پیوسته تبدیل شوند. یک راه ساده این است که تعداد اسپایک‌ها یا نرخ اسپایک در چند گام زمانی به یک مقدار عددی پیوسته نگاشت شود.

جریان کلی چنین است:

\begin{enumerate}
  \item گرفتن \eng{embedding} اولیه از \eng{BERT}
  \item تبدیل \eng{embedding} پیوسته به اسپایک
  \item پردازش در \eng{SpikeBERT}
  \item تبدیل خروجی اسپایکی به مقدار پیوسته
  \item مقایسه با خروجی \eng{Teacher}
\end{enumerate}

با این حال، پیوسته‌کردن خروجی کافی نیست. فضای هندسی \eng{Student} نیز باید با فضای \eng{Teacher} قابل مقایسه شود.

\subsection{\eng{Projection} و هم‌ترازسازی بازنمایی}
پس از تبدیل خروجی اسپایکی به مقدار پیوسته، باید آن را به فضایی منتقل کنیم که با خروجی \eng{Teacher} قابل مقایسه باشد. یک راه ساده استفاده از یک لایهٔ خطی است:

\[
H' = HW + b
\]

در \eng{SpikeBERT} به‌جای یک لایهٔ خطی ساده، از \eng{MLP} استفاده می‌شود تا \eng{projection} قوی‌تری انجام شود:

\[
H' = MLP(H)
\]

پس از این تبدیل، خروجی \eng{Student} و خروجی \eng{Teacher} با هم مقایسه می‌شوند و خطا محاسبه می‌شود.

\subsection{\eng{Feature Loss}}
برای محاسبهٔ \eng{Feature Loss} یا \eng{Embedding Loss}، خروجی \eng{Teacher} و \eng{Student} را می‌توان خانه‌به‌خانه مقایسه کرد. اگر بخشی از خروجی‌ها چنین باشد:

\[
\text{\eng{Teacher}} = [1,\ 2,\ 0]
\]

\[
\text{\eng{Student}} = [0.5,\ 0.3,\ 1]
\]

یک شکل سادهٔ خطا چنین است:

\[
L = (1 - 0.5)^2 + (2 - 0.3)^2 + (0 - 1)^2
\]

همین منطق می‌تواند برای \eng{embedding loss} و \eng{feature loss} در لایه‌های مختلف استفاده شود. البته لازم نیست همهٔ لایه‌ها با \eng{Teacher} مقایسه شوند. انتخاب لایه‌هایی که در آن‌ها هم‌ترازسازی انجام می‌شود، می‌تواند بخشی از طراحی مدل باشد.

\subsection{دو مرحلهٔ آموزش در \eng{SpikeBERT}}
آموزش \eng{SpikeBERT} را می‌توان در دو مرحله توضیح داد. مرحلهٔ اول انتقال دانش عمومی از \eng{Teacher} به \eng{Student} است و مرحلهٔ دوم آموزش مخصوص یک وظیفهٔ مشخص.

\subsubsection{انتقال دانش عمومی}
در مرحلهٔ اول، هدف این است که \eng{Student} دانش عمومی \eng{Teacher} را یاد بگیرد. این دانش شامل \eng{embedding}ها و ویژگی‌های میانی شبکه است. در این بخش، \eng{embedding loss} و \eng{hidden feature loss} اهمیت بیشتری دارند، چون مدل اسپایکی باید یاد بگیرد نمایش‌های داخلی خود را به نمایش‌های \eng{BERT} معمولی نزدیک کند.

\subsubsection{آموزش مخصوص وظیفه}
در مرحلهٔ دوم، مدل برای یک وظیفهٔ مشخص، مانند \eng{Sentiment Analysis}، تنظیم می‌شود. در اینجا خروجی‌های نهایی یا \eng{logit}های \eng{Teacher} و \eng{Student} نیز با هم مقایسه می‌شوند.

برای مثال، اگر \eng{Teacher} برای دو کلاس خروجی زیر را بدهد:

\[
[0.2,\ 0.1]
\]

و \eng{Student} چنین خروجی‌ای تولید کند:

\[
[0.5,\ 0]
\]

\eng{Logit Loss} می‌تواند به شکل زیر محاسبه شود:

\[
L_{\text{logit}} = (0.2 - 0.5)^2 + (0.1 - 0)^2
\]

\subsection{\eng{Cross Entropy Loss}}
در کنار مقایسه با \eng{Teacher}، خروجی \eng{Student} با برچسب واقعی یا \eng{Ground Truth} نیز مقایسه می‌شود. اگر برچسب واقعی کلاس مثبت باشد، می‌توان آن را به شکل زیر نمایش داد:

\[
[1,\ 0]
\]

اگر خروجی \eng{Student} مثلاً چنین باشد:

\[
[0.6,\ 0.3]
\]

برای خطای طبقه‌بندی از \eng{Cross Entropy Loss} استفاده می‌شود. فرم کلی آن برای کلاس درست چنین است:

\[
L_{\text{CE}} =
-\log\left(
\frac{e^{z_y}}{\sum_j e^{z_j}}
\right)
\]

در این رابطه، \(z_y\) \eng{logit} مربوط به کلاس درست است.

در نهایت، همهٔ خطاها با ضرایب مشخص جمع می‌شوند:

\[
\begin{aligned}
L_{\text{total}} ={}&
\lambda_1 L_{\text{embedding}}
+ \lambda_2 L_{\text{feature}} \\
&+ \lambda_3 L_{\text{logit}}
+ \lambda_4 L_{\text{CE}}
\end{aligned}
\]

سپس بر اساس این خطای نهایی، \eng{backpropagation} انجام می‌شود. چون شبکه اسپایکی است، برای عبور گرادیان از اسپایک‌ها از \eng{Surrogate Gradient} استفاده می‌شود. نکتهٔ مهم این است که حتی \eng{embedding} اولیهٔ \eng{SpikeBERT} نیز در طول آموزش \eng{fine-tune} می‌شود.

\coursefigure{0.95}{spikebert-knowledge-distillation.png}{تقطیر دو مرحله‌ای از \eng{BERT} به \eng{SpikeBERT}: هم‌ترازسازی امبدینگ‌ها و ویژگی‌های میانی در کنار زیان لاجیت و زیان طبقه‌بندی.}

\subsection{نقش \eng{Time Step}}
\eng{Time Step} ابرپارامتری است که در شبکه‌های معمولی وجود ندارد، اما در شبکه‌های اسپایکی نقش اصلی دارد:

\[
T = \text{\eng{Time Step}}
\]

پرسش این است که زیاد یا کم کردن \(T\) چه اثری روی دقت، همگرایی، زمان آموزش و مصرف حافظه دارد. در این بحث گفته شد که هرچه \(T\) بزرگ‌تر شود، همگرایی معمولاً دیرتر می‌شود، چون مدل باید گام‌های زمانی بیشتری را پردازش کند.

\subsubsection{رابطهٔ \eng{Time Step} و دقت}
برای توضیح رابطهٔ \eng{Time Step} و دقت، می‌توان از مثال تصمیم‌گیری انسانی استفاده کرد. فرض کنید در خانه نشسته‌ایم و صدایی از پارکینگ می‌آید. بعد صدای آسانسور شنیده می‌شود، آسانسور به طبقهٔ ما نزدیک می‌شود و سپس صدای قفل در می‌آید. در هر مرحله شواهد بیشتری جمع می‌شود و بهتر می‌توانیم حدس بزنیم چه کسی وارد خانه شده است.

این روند شبیه \eng{Evidence Accumulation} در نورون‌های اسپایکی است: شواهد در طول زمان جمع می‌شوند تا از آستانه عبور کنند و تصمیم ساخته شود. بنابراین افزایش \(T\) می‌تواند تا حدی دقت را بیشتر کند، چون اطلاعات زمانی بیشتری در اختیار مدل قرار می‌گیرد.

\subsubsection{محدودیت‌های \eng{Time Step} زیاد}
افزایش \eng{Time Step} همیشه مفید نیست. رابطهٔ کلی \(T\) و دقت معمولاً ابتدا صعودی است، اما از یک نقطه به بعد به حالت اشباع می‌رسد. یعنی ابتدا با افزایش \(T\)، دقت بهتر می‌شود؛ سپس مقدار بهبود کم می‌شود؛ ممکن است دقت ثابت بماند؛ و حتی ممکن است به دلیل توجه به نویز افت کند.

پس هدف این است که کمترین \eng{Time Step}ای پیدا شود که بیشترین دقت عملی را می‌دهد. این نقطه همان دقت زمانی مناسب یا \eng{temporal resolution} مناسب مدل است.

\coursefigure{0.68}{time-steps-and-model-accuracy.png}{اثر تعداد گام‌های زمانی بر دقت: افزایش \eng{Time Step} ابتدا به تجمع شواهد و بهبود دقت کمک می‌کند، اما پس از چند گام به اشباع می‌رسد.}

چند دلیل باعث می‌شود \(T\) خیلی بزرگ مشکل‌ساز شود:

\begin{itemize}
  \item از یک حد به بعد، اطلاعات مفید تمام می‌شود و مدل ممکن است به نویز حساس شود.
  \item با افزایش \(T\)، بار حافظه‌ای مدل بیشتر می‌شود. این نکته با بحث ظرفیت \eng{Working Memory} در جلسات قبل نیز ارتباط دارد.
  \item هرچه \(T\) بزرگ‌تر باشد، آموزش و \eng{inference} کندتر و سنگین‌تر می‌شوند.
\end{itemize}

خلاصهٔ بحث چنین است:

\[
T \uparrow \Rightarrow \text{\eng{Accuracy}} \uparrow
\quad \text{تا نقطهٔ اشباع}
\]

اما پس از آن:

\[
T \uparrow \Rightarrow
\text{هزینهٔ بیشتر، نویز بیشتر، و بهبود کمتر یا حتی افت دقت}
\]

\subsection{بازگشت کوتاه به آزمون \eng{Spiking Attention}}
در ادامه، به آزمون \eng{Attention} اسپایکی برگشت. هدف این آزمون این بود که اگر ورودی پیوسته با \eng{Rate Encoding} به اسپایک تبدیل شود و سپس \eng{Attention} محاسبه شود، همچنان بتوان تغییر معنادار در \eng{embedding} را مشاهده کرد.

مثال مربوط به کلمهٔ «شیر» بود که دو معنا دارد: شیر به‌معنای حیوان و شیر به‌معنای مایع. اگر این کلمه کنار «جنگل» بیاید، انتظار داریم نمایش آن به معنای حیوانی نزدیک‌تر شود.

برای ساده‌سازی، وزن‌های \(\ W_Q\)، \(\ W_K\) و \(\ W_V\) همانی فرض شدند:

\[
Q = K = V = X
\]

مراحل آزمون چنین بود: ابتدا \eng{embedding} پیوسته به اسپایک تبدیل می‌شود؛ سپس امتیازهای توجه محاسبه می‌شوند؛ خروجی دوباره پیوسته می‌شود؛ و در نهایت باید نشان داد که بعد حیوانی کلمهٔ «شیر» افزایش و بعد مایع‌بودن آن کاهش می‌یابد.

اگر \eng{embedding} کلمهٔ «شیر» چنین باشد:

\[
[0.5,\ 0.8]
\]

با \eng{Rate Encoding} و \(T=4\)، مقدار \(0.5\) تقریباً دو اسپایک و مقدار \(0.8\) تقریباً سه اسپایک تولید می‌کند. در مثال، امتیاز توجه «شیر به شیر» برابر \(4\) و امتیاز توجه «شیر به جنگل» برابر \(2\) در نظر گرفته شد. پس:

\[
4 + 2 = 6
\]

\[
\text{\eng{Attention}}(\text{شیر به شیر}) =
\frac{4}{6} = \frac{2}{3}
\]

\[
\text{\eng{Attention}}(\text{شیر به جنگل}) =
\frac{2}{6} = \frac{1}{3}
\]

بنابراین نمایش جدید «شیر» از ترکیب \eng{value}های «شیر» و «جنگل» ساخته می‌شود:

\[
V'_{\text{شیر}} =
\frac{2}{3}V_{\text{شیر}}
+ \frac{1}{3}V_{\text{جنگل}}
\]

نتیجهٔ مورد انتظار این است که زمینهٔ «جنگل» بعد حیوانی «شیر» را تقویت کند و بعد مایع‌بودن آن کاهش یابد.

\section{\eng{SpikeGPT}، \eng{RWKV} و تولید متن اسپایکی}

پس از \eng{SpikeFormer} و \eng{SpikeBERT}، آخرین معماری مهمی که در ادامهٔ بحث شبکه‌های اسپایکی بررسی شد \eng{SpikeGPT} است. ایدهٔ کلی این است که برای بیشتر معماری‌های رایج یادگیری عمیق، دست‌کم یک طرح اولیه برای پیاده‌سازی اسپایکی مطرح شده است. \eng{SpikeGPT} اهمیت ویژه دارد، چون نشان می‌دهد چگونه می‌توان مدل‌های زبانی خودرگرسیو مانند \eng{GPT} را به سمت معماری‌های اسپایکی برد.

\subsection{تفاوت طبقه‌بندی و تولید}
برای فهم \eng{SpikeGPT}، ابتدا باید تفاوت دو نوع وظیفهٔ اصلی را روشن کرد: \eng{classification} و \eng{generation}.

در \eng{classification}، مدل باید کل ورودی را بفهمد و در نهایت یک برچسب به آن اختصاص دهد. برای نمونه، در \eng{BERT}، چون مدل باید هم به کلمه‌های قبل و هم به کلمه‌های بعد توجه کند، معماری معمولاً \eng{bidirectional} است. \eng{BERT} یک معماری \eng{encoder-only} است: کل دنباله را \eng{encode} می‌کند و سپس از خروجی آن برای طبقه‌بندی استفاده می‌شود. به همین دلیل، در \eng{SpikeBERT} نمی‌شد به‌سادگی از \eng{attention} صرف‌نظر کرد، چون طبقه‌بندی به فهم کامل متن از دو جهت وابسته است.

در \eng{generation}، هدف تولید \eng{token} بعدی است. مدل در زمان تولید فقط به گذشته دسترسی دارد، نه آینده. اگر یک \eng{prompt} داده شود، مدل کلمه‌ها را یکی‌یکی تولید می‌کند: کلمه‌های قبلی را می‌بیند، کلمهٔ بعدی را تولید می‌کند، همان کلمه به \eng{context} اضافه می‌شود، و سپس \eng{token} بعدی تولید می‌شود.

پس \eng{GPT} یک مدل علّی و خودرگرسیو است. معمولاً با \eng{masked self-attention} کار می‌کند؛ یعنی \eng{attention} فقط اجازه دارد به گذشته نگاه کند.

در آموزش تولید متن، یک دنبالهٔ ورودی داریم:

\[
x_1,\ x_2,\ x_3,\ \ldots,\ x_s
\]

برچسب مطلوب هر موقعیت، \eng{token} بعدی است. اگر جملهٔ نمونه «آسمان آبی است.» باشد، ورودی «آسمان» باید به برچسب «آبی» برسد، ورودی «آبی» باید به «است» برسد، و ورودی «است» باید نقطه را پیش‌بینی کند. در \eng{SpikeGPT} نیز خروجی مدل با برچسب‌های جابه‌جا شده مقایسه می‌شود و معمولاً از \eng{cross entropy loss} استفاده می‌شود. برای آموزش بخش‌های اسپایکی نیز مانند قبل از \eng{surrogate gradient} کمک گرفته می‌شود.

تفاوت استفاده از \eng{SpikeGPT} در تولید و طبقه‌بندی نیز مهم است. در تولید، خروجی نهایی باید به \eng{vocabulary} وصل شود تا مدل بتواند \eng{logit} مربوط به \eng{token}های مختلف را بسازد و \eng{token} بعدی را انتخاب کند. اما در طبقه‌بندی، الزاماً نیازی به \eng{head} تولید \eng{token} نیست. ورودی از بدنهٔ مدل عبور می‌کند، \eng{hidden state}های میانی گرفته می‌شوند، مثلاً روی آن‌ها \eng{average pooling} انجام می‌شود، و سپس کلاس نهایی پیش‌بینی می‌شود.

\coursefigure{0.82}{spikegpt-training-and-classification.png}{دو حالت استفاده از \eng{SpikeGPT}: آموزش تولیدی با برچسب‌های شیفت‌یافته و استفادهٔ طبقه‌بندی با میانگین‌گیری و سر طبقه‌بندی.}

\subsection{چرا تبدیل مستقیم \eng{Transformer} کافی نیست؟}
فرض پیش‌فرض در این بحث این بود که اگر خود \eng{Transformer} را مستقیم برداریم و فقط آن را اسپایکی کنیم، معمولاً نتیجهٔ خوبی نمی‌گیریم. در \eng{SpikeBERT} هم دیده شد که حتی برای \eng{BERT}، که از نظر تولید متن ساده‌تر از \eng{GPT} است، تبدیل مستقیم کافی نبود و به \eng{Teacher} و \eng{Distillation} نیاز داشت.

بنابراین برای \eng{GPT}، که باید \eng{generation} انجام دهد، مسئله سخت‌تر است. نمی‌توان فقط با کارهایی مانند گذاشتن \eng{encoder} یا \eng{decoder} اسپایکی، تبدیل نمایش پیوسته به اسپایک و اسپایک به نمایش پیوسته، یا حذف و تغییر تابع‌های فعال‌سازی انتظار داشت مدل به‌خوبی کار کند. \eng{SpikeGPT} به همین دلیل مسیر دیگری را انتخاب می‌کند و به‌جای اتکا به \eng{attention} معمولی، از ایدهٔ \eng{RWKV} استفاده می‌کند.

\subsection{شباهت \eng{GPT} و \eng{RNN}}
\eng{GPT} ذاتاً \eng{decoder-only} است و فقط به گذشته نگاه می‌کند. این ویژگی آن را از یک جهت به \eng{RNN} نزدیک می‌کند. \eng{RNN} نیز ترتیبی است و اطلاعات گذشته را در یک \eng{state} نگه می‌دارد.

در \eng{RNN}، مدل در هر لحظه \eng{token} فعلی و \eng{state} قبلی را می‌گیرد، اطلاعات جدید را وارد \eng{state} می‌کند، \eng{state} را به‌روزرسانی می‌کند و خروجی جدید می‌دهد. به‌جای اینکه مانند \eng{attention} همهٔ \eng{token}ها دوبه‌دو با هم مقایسه شوند، یک \eng{state} وجود دارد که خلاصه‌ای از گذشته را نگه می‌دارد.

برای مثال، اگر ابتدا کلمهٔ «شیر» دیده شود، \eng{state} ممکن است هم معنای حیوانی و هم معنای مایع را در خود نگه دارد. اگر بعداً کلمهٔ «جنگل» بیاید، \eng{state} به‌روزرسانی می‌شود و معنای حیوانی تقویت می‌شود.

\subsection{مشکل \eng{RNN}های کلاسیک}
مشکل اصلی \eng{RNN}های کلاسیک \eng{vanishing gradient} و فراموشی اطلاعات بلندمدت است. علت آن تا حدی به استفاده از تابع‌های غیرخطی مانند \(\tanh\) روی \eng{state} برمی‌گردد. وقتی دنباله طولانی می‌شود، گرادیان به‌تدریج کوچک می‌شود و مدل نمی‌تواند وابستگی‌های بلندمدت را خوب حفظ کند.

\eng{Attention} این مشکل را به همان شکل ندارد، چون بخش عمدهٔ آن بر ضرب و جمع خطی متکی است. اما هزینهٔ محاسباتی بالایی دارد. برای \(n\) \eng{token}، ماتریس \eng{attention} معمولاً با پیچیدگی حدودی \(O(n^2)\) ساخته می‌شود. بنابراین برای مدل اسپایکی، بهتر است راهی پیدا شود که هم خاصیت ترتیبی و نگهداری گذشته را داشته باشد و هم هزینهٔ \eng{attention} کامل را تحمیل نکند.

\subsection{معماری \eng{RWKV}}
\eng{RWKV} ترکیبی از ایده‌های \eng{RNN} و \eng{Transformer} است. از یک طرف، مانند \eng{RNN} یک \eng{state} دارد؛ از طرف دیگر، تلاش می‌کند وابستگی‌های بلندمدت را بهتر از \eng{RNN} کلاسیک حفظ کند. \eng{state} بازگشتی آن خطی‌تر است و تابع غیرخطی سنگین مستقیماً روی \eng{state} اعمال نمی‌شود. بنابراین مشکل \eng{vanishing gradient} نسبت به \eng{RNN} کلاسیک کمتر می‌شود.

در \eng{SpikeGPT}، به‌جای استفادهٔ مستقیم از \eng{attention}، از \eng{RWKV} استفاده می‌شود. هدف این است که معماری سبک‌تر و برای اسپایکی‌شدن مناسب‌تر باشد.

\subsubsection{\eng{Receptance}}
\eng{Receptance} یا \(R\) نقش \eng{gate} را دارد و مشخص می‌کند چه مقدار از اطلاعات از \eng{state} عبور کند و وارد خروجی شود.

\subsubsection{\eng{Weight} یا عامل فراموشی}
\eng{Weight} یا \(W\) در این بحث نقش عامل فراموشی را دارد. این مقدار تعیین می‌کند چه بخشی از اطلاعات قبلی در \eng{state} حفظ شود و چه مقدار کاهش پیدا کند.

\subsubsection{\eng{Key}}
\eng{Key} یا \(K\) مشابه \eng{key} در \eng{attention} است و اهمیت \eng{token} فعلی را نشان می‌دهد.

\subsubsection{\eng{Value}}
\eng{Value} یا \(V\) مقدار یا نمایش محتوایی مربوط به \eng{token} فعلی است. در \eng{RWKV} دیگر \eng{query} نداریم، چون قرار نیست برای هر \eng{token} همهٔ \eng{token}های دیگر را جداگانه \eng{query} کنیم. به‌جای آن، یک \eng{state} وجود دارد که گذشته را نگه می‌دارد.

\subsection{ایدهٔ \eng{State} در \eng{RWKV}}
در هر لحظه، مدل یک \eng{state} دارد که خلاصه‌ای از چیزهایی است که تا آن زمان دیده است. برای \eng{token} فعلی، \(K\) اهمیت \eng{token} را مشخص می‌کند، \(V\) محتوای آن را نشان می‌دهد، \(W\) تعیین می‌کند از گذشته چقدر فراموش شود، و \(R\) تعیین می‌کند چه مقدار از \eng{state} فعلی وارد خروجی شود.

برای ساده‌سازی، یک مثال یک‌بعدی در نظر می‌گیریم. فرض کنیم برای \eng{token} اول:

\[
K_1 = 2
\]

\[
V_1 = 10
\]

در نتیجه صورت اولیه چنین می‌شود:

\[
K_1V_1 = 2 \times 10 = 20
\]

و مخرج برابر است با:

\[
K_1 = 2
\]

پس \eng{state} یا مخزن اولیه:

\[
\frac{20}{2} = 10
\]

اکنون \eng{token} دوم دیده می‌شود:

\[
K_2 = 3,\qquad V_2 = 4,\qquad R_2 = 0.9
\]

و عامل فراموشی:

\[
W = 0.5
\]

اطلاعات قبلی با عامل فراموشی کم می‌شود و اطلاعات جدید اضافه می‌شود. صورت جدید:

\[
0.5 \times 20 + 3 \times 4 = 10 + 12 = 22
\]

مخرج جدید:

\[
0.5 \times 2 + 3 = 1 + 3 = 4
\]

پس \eng{state} جدید:

\[
\frac{22}{4} = 5.5
\]

خروجی با \eng{gate} یا همان \(R\) ساخته می‌شود:

\[
\text{\eng{Output}} = R \times \text{\eng{State}}
\]

\[
\text{\eng{Output}} = 0.9 \times 5.5 = 4.95
\]

این خروجی همان نمایش جدیدی است که به لایهٔ بعدی داده می‌شود.

\subsection{تفاوت \eng{Time Step} در \eng{RWKV} و \eng{SNN}}
در فرمول‌های \eng{RWKV}، \(t\) همان گام زمانی دنبالهٔ \eng{token}هاست؛ یعنی مثلاً \eng{token} قبلی و \eng{token} فعلی. این \(t\) نباید با بعد زمانی در شبکه‌های اسپایکی اشتباه گرفته شود.

در شبکهٔ اسپایکی، خود \eng{embedding} ممکن است در چند \eng{time step} اسپایکی کدگذاری شود. اما در \eng{RWKV}، منظور از \(t\)، موقعیت \eng{token} در \eng{sequence} است. بنابراین \eng{SpikeGPT} هم با زمان زبانی یا ترتیبی سروکار دارد و هم با زمان اسپایکی.

\subsection{معماری کلی \eng{SpikeGPT}}
روند کلی در \eng{SpikeGPT} چنین است:

\begin{enumerate}
  \item ابتدا \eng{embedding} ورودی به حالت باینری یا اسپایکی تبدیل می‌شود.
  \item این تبدیل معمولاً با نورون‌هایی مانند \eng{LIF} انجام می‌شود.
  \item ورودی اسپایکی وارد بلوک \eng{RWKV} می‌شود.
  \item خروجی \eng{RWKV} دوباره از یک لایهٔ اسپایکی عبور می‌کند.
  \item سپس خروجی وارد بخش \eng{feed-forward} اصلاح‌شده، یعنی \eng{R-FFN}، می‌شود.
  \item خروجی \eng{R-FFN} نیز دوباره اسپایکی می‌شود.
  \item بسته به وظیفه، خروجی برای تولید متن یا طبقه‌بندی استفاده می‌شود.
\end{enumerate}

بنابراین در ابتدا و انتهای هر بلوک تلاش می‌شود نمایش دوباره به حالت اسپایکی برگردد.

\coursefigure{0.95}{spikegpt-architecture.png}{ساختار \eng{SpikeGPT}: بلوک‌های \eng{Spiking RWKV} و \eng{Spiking RFFN} جایگزین توجه پرهزینه می‌شوند و محاسبات را رویدادمحور و تنک‌تر می‌کنند.}

\subsubsection{\eng{Time-Mix}}
در \eng{Transformer} کلاسیک، بخش \eng{attention} مسئول فهم \eng{context} است؛ یعنی مدل می‌فهمد \eng{token} فعلی در کنار \eng{token}های قبلی چه معنایی دارد. در \eng{SpikeGPT} این نقش با \eng{RWKV} جایگزین می‌شود.

به بلوک \eng{RWKV} معمولاً \eng{Time-Mix} گفته می‌شود، چون اطلاعات زمانی و \eng{token}های قبلی را با \eng{token} فعلی ترکیب می‌کند.

\subsubsection{\eng{Channel-Mix}}
در \eng{Transformer} کلاسیک، بخش \eng{feed-forward} یا \eng{MLP} روی خود \eng{token} فعلی عمیق‌تر می‌شود و ویژگی‌های پیچیده‌تر استخراج می‌کند. در \eng{SpikeGPT} این بخش با \eng{R-FFN} جایگزین می‌شود.

به این بخش \eng{Channel-Mix} گفته می‌شود، چون تمرکز اصلی آن دیگر روی زمان نیست، بلکه روی ویژگی‌ها و کانال‌های خود \eng{token} فعلی است.

\subsubsection{\eng{R-FFN}}
\eng{R-FFN} از نظر نقش کلی شبیه \eng{feed-forward} در \eng{Transformer} است. در \eng{Transformer} کلاسیک، معمولاً ابتدا بُعد نمایش بزرگ‌تر می‌شود، سپس یک تابع فعال‌سازی مانند \eng{ReLU} اعمال می‌شود، بعد بُعد دوباره فشرده می‌شود و نمایش بهتر ساخته می‌شود.

در \eng{R-FFN} نیز همین ایده وجود دارد. ورودی ابتدا \eng{expand} می‌شود تا مدل بتواند الگوهای پیچیده‌تری استخراج کند، سپس از تابع فعال‌سازی عبور می‌کند و دوباره \eng{compress} می‌شود. در این بخش سه ماتریس مطرح می‌شود: یک ماتریس برای \eng{gate} یا \(G\)، یک ماتریس برای مقدار اصلی یا \(P\)، و یک ماتریس برای \eng{projection} نهایی.

بخشی از مسیر بزرگ‌تر می‌شود، از \eng{ReLU} عبور می‌کند و دوباره فشرده می‌شود. سپس با مقدار اولیه به‌صورت ضرب عنصر‌به‌عنصر یا \eng{Hadamard product} ترکیب می‌شود:

\[
[a_1,\ a_2,\ a_3] \odot [b_1,\ b_2,\ b_3]
=
[a_1b_1,\ a_2b_2,\ a_3b_3]
\]

هدف این است که به ویژگی‌های مختلف وزن داده شود.

\subsection{پیوستگی محاسبات میانی و اسپایکی بودن مدل}
یک پرسش مهم این است: اگر پس از اسپایک دوباره ضرب ماتریسی انجام می‌شود و خروجی پیوسته می‌شود، چه تفاوتی با یک شبکهٔ کاملاً غیر اسپایکی دارد؟

پاسخ این است که در شبکه‌های اسپایکی لازم نیست همهٔ نقاط معماری همیشه صفر و یک باشند. در بسیاری از بخش‌ها مقدارهای پیوسته دیده می‌شوند، اما این پیوستگی با پیوستگی در شبکه‌های معمولی یکسان نیست.

نخست، خروجی لایه‌های اسپایکی بسیار \eng{sparse} است. بسیاری از مقدارها صفرند و فقط وقتی اسپایک رخ می‌دهد، محاسبه فعال می‌شود. همین \eng{sparsity} می‌تواند محاسبات و مصرف انرژی را کاهش دهد.

دوم، مقدارهای پیوستهٔ میانی در اینجا بیشتر به پتانسیل غشایی یا \eng{membrane potential} مربوط‌اند، نه به همان معنای معمول ویژگی‌های پیوسته در \eng{ANN}. اطلاعات روی زمان کدگذاری شده است. بنابراین اگرچه از نظر ظاهری عدد پیوسته دیده می‌شود، ماهیت آن با مقدار پیوسته در شبکهٔ غیر اسپایکی تفاوت دارد.

\subsection{مصرف انرژی و محدودیت اجرای \eng{SNN} روی \eng{GPU}}
\eng{SNN}ها ذاتاً برای سخت‌افزارهای نورومورفیک طراحی شده‌اند، در حالی که \eng{ANN}ها معمولاً روی \eng{GPU} یا \eng{TPU} اجرا می‌شوند. اگر \eng{SNN} روی \eng{GPU} اجرا شود، ممکن است بخشی از مزیت اصلی آن دیده نشود؛ چون \eng{GPU} برای محاسبات متراکم و موازی طراحی شده است، نه برای محاسبات \eng{sparse} و \eng{event-driven}.

شبکهٔ اسپایکی را می‌توان شبیه یک گراف \eng{sparse} دید: گره‌های زیادی وجود دارد، اما فقط بعضی ارتباط‌ها در زمان‌های خاص فعال می‌شوند. شبکهٔ غیر اسپایکی بیشتر شبیه یک گراف متراکم است که ارتباط‌های زیادی همیشه فعال‌اند. بنابراین اجرای \eng{SNN} روی \eng{GPU} الزاماً مزیت واقعی آن را نشان نمی‌دهد؛ مزیت اصلی زمانی روشن‌تر می‌شود که مدل روی سخت‌افزار نورومورفیک اجرا شود.

برای مقایسهٔ منصفانهٔ مصرف انرژی، معمولاً به‌جای اجرای مستقیم روی دو سخت‌افزار متفاوت، تعداد عملیات‌ها شمرده می‌شود. در شبکهٔ غیر اسپایکی، عملیات اصلی معمولاً \eng{MAC} است:

\[
\text{\eng{MAC}} = \text{\eng{Multiply}} + \text{\eng{Accumulate}}
\]

اما در شبکهٔ اسپایکی، چون ورودی‌ها صفر و یکی هستند، بسیاری از عملیات‌ها به جمع کاهش پیدا می‌کنند:

\[
\text{\eng{AC}} = \text{\eng{Accumulate}}
\]

\coursefigure{0.95}{gpt-vs-spikegpt-energy-use.png}{مقایسهٔ مصرف انرژی \eng{Vanilla GPT} و \eng{SpikeGPT}: تبدیل عملیات \eng{MAC} به \eng{AC} در بخش‌های اصلی مدل، کاهش انرژی کلی در این برآورد را نشان می‌دهد.}

انرژی \eng{MAC} از \eng{AC} بیشتر است. بنابراین اگر تعداد \eng{MAC}ها در \eng{Vanilla GPT} و تعداد \eng{AC}ها در \eng{SpikeGPT} شمرده شود، می‌توان مصرف انرژی تقریبی دو مدل را مقایسه کرد. طبق گزارش مقاله، مقالهٔ \eng{SpikeGPT} کاهش انرژی حدوداً ۳۲ برابری را گزارش می‌کند.

جمع‌بندی این بخش چنین است: \eng{GPT} مدلی \eng{decoder-only} و خودرگرسیو است؛ برای تولید متن فقط به گذشته نیاز دارد؛ \eng{masked self-attention} کامل می‌تواند برای نسخهٔ اسپایکی ناکارآمد باشد؛ تبدیل مستقیم \eng{Transformer} به \eng{SNN} معمولاً کافی نیست؛ و \eng{SpikeGPT} به همین دلیل از \eng{RWKV} استفاده می‌کند. در این معماری، \eng{Time-Mix} نقش فهم زمینه را دارد، \eng{Channel-Mix} یا \eng{R-FFN} ویژگی‌های عمیق‌تر را از \eng{token} فعلی استخراج می‌کند، و خروجی بلوک‌ها با \eng{LIF} یا لایه‌های مشابه دوباره اسپایکی می‌شود.

\chapter{هم‌ترازی مغز، کدگذاری عصبی و استدلال}

\section{کدگذاری و رمزگشایی عصبی}

\subsection{\eng{Neural Encoding} و \eng{Neural Decoding}}
پس از پایان بحث شبکه‌های اسپایکی، مسیر بحث به سمت \eng{Neural Encoding} و \eng{Neural Decoding} رفت. هدف این بخش این است که روشن شود اطلاعات چگونه در مغز بازنمایی می‌شوند و چگونه می‌توان از روی فعالیت مغزی، رفتار، فرمان یا متن را بازسازی کرد.

در علوم اعصاب، \eng{Encoding} به این معناست که یک محرک خارجی وارد سیستم عصبی می‌شود و به نوعی فعالیت یا فرمان در مغز تبدیل می‌گردد. صدا، تصویر، لمس یا هر محرک بیرونی دیگر می‌تواند چنین مسیری را آغاز کند.

در مقابل، \eng{Decoding} یعنی از روی فعالیت‌های مغزی بتوانیم یک خروجی خارجی را بازسازی کنیم؛ برای مثال یک فرمان حرکتی، یک رفتار، یک متن یا حتی نیت فرد را حدس بزنیم. پس مسیر کلی را می‌توان چنین خلاصه کرد:

\[
\text{محرک خارجی} \longrightarrow \text{فعالیت مغزی} \longrightarrow \text{فرمان یا رفتار}
\]

در \eng{Encoding} از محرک به فعالیت مغزی می‌رسیم، و در \eng{Decoding} مسیر را برعکس طی می‌کنیم: از فعالیت مغزی تلاش می‌کنیم محرک، نیت یا فرمان را بازسازی کنیم.

\subsection{نظریه‌های اصلی \eng{Neural Coding}}
سه نظریهٔ اصلی دربارهٔ کدگذاری عصبی در این بحث به ترتیب تاریخی بررسی شدند:

\begin{enumerate}
  \item \eng{Grandmother Cell Theory}
  \item \eng{Population Coding}
  \item \eng{Sparse Coding}
\end{enumerate}

هیچ‌کدام از این نظریه‌ها در این بحث به‌عنوان پاسخ مطلق معرفی نشدند. هرکدام بخشی از شهود لازم برای فهم مغز و مدل‌های محاسباتی امروزی را فراهم می‌کنند.

\subsubsection{\eng{Grandmother Cell}}
نظریهٔ \eng{Grandmother Cell} حدود دههٔ ۱۹۶۰ مطرح شد. طبق این دیدگاه، برای هر مفهوم، شیء یا موجودیت خارجی، یک سلول یا یک گروه بسیار محدود و اختصاصی از نورون‌ها وجود دارد. مثال کلاسیک این است که برای مفهوم «مادربزرگ»، یک نورون یا گروهی خاص از نورون‌ها مسئول تشخیص همان مفهوم باشد.

نکتهٔ اصلی این نظریه انحصار است. منظور لزوماً این نیست که دقیقاً فقط یک نورون وجود دارد، بلکه ایده این است که مسئولیت شناخت یک مفهوم به‌شکل اختصاصی در اختیار یک سلول یا گروه کوچک قرار می‌گیرد.

این نظریه چند اشکال روشن دارد. نخست، اگر آن سلول یا گروه خاص آسیب ببیند، باید کل مفهوم از دست برود؛ درحالی‌که حافظه و شناخت معمولاً تا این اندازه به یک واحد منفرد وابسته نیستند. دوم، اگر برای هر چهره، صدا، شیء، مفهوم و حالت جهان به یک نورون یا گروه اختصاصی نیاز باشد، ظرفیت لازم بسیار بزرگ و از نظر زیستی و محاسباتی نامعقول می‌شود.

با این حال، ایدهٔ \eng{Grandmother Cell} کاملاً بی‌فایده نیست. در برخی سطوح مدل‌های عصبی، به‌ویژه در لایهٔ خروجی شبکه‌های طبقه‌بندی، شباهتی با این ایده دیده می‌شود. اگر یک شبکه سه کلاس داشته باشد و برای هر کلاس یک نود خروجی تعریف شود، هر نود تا حدی نقشی شبیه یک بازنمایی اختصاصی برای آن کلاس پیدا می‌کند. این تشبیه، نسخهٔ زیستی و مطلق \eng{Grandmother Cell} را اثبات نمی‌کند، اما نشان می‌دهد ایدهٔ «نمایندهٔ اختصاصی» در برخی معماری‌ها ظاهر می‌شود.

\subsubsection{\eng{Population Coding}}
\eng{Population Coding} در برابر دیدگاه اختصاصی \eng{Grandmother Cell} قرار می‌گیرد و حدوداً از دههٔ ۱۹۸۰ برجسته‌تر شد. در این نظریه، یک سلول یا گروه کوچک مسئول یک مفهوم نیستند؛ بلکه کل جمعیت نورون‌ها در تصمیم‌گیری نقش دارند.

این ایده را می‌توان به یک نظرسنجی جمعی تشبیه کرد. خروجی نهایی از مشارکت تعداد زیادی نورون ساخته می‌شود و هر نورون، بسته به فعالیت و وزن خود، سهمی در تصمیم دارد. در مدل‌های عصبی مصنوعی نیز شکل سادهٔ این ایده را در جمع وزن‌دار می‌بینیم:

\[
\sum_i W_i X_i
\]

در اینجا \(X_i\) می‌تواند فعالیت یا اسپایک نورون باشد و \(W_i\) وزن سیناپسی یا وزن یادگرفته‌شدهٔ آن. همهٔ نورون‌ها به یک اندازه تعیین‌کننده نیستند، اما تصمیم نهایی حاصل مشارکت جمعی آن‌هاست.

برای توضیح این نظریه، می‌توان به آزمایش‌هایی روی حرکت میمون‌ها اشاره کرد. در این نوع آزمایش‌ها، میمون برای حرکت یا پرتاب به جهت‌های مختلف مانند چپ، راست، بالا و پایین آموزش داده می‌شود و سپس فعالیت نورون‌های مغزی هنگام اجرای حرکت بررسی می‌گردد. مشاهدهٔ مهم این است که هنگام حرکت به یک جهت خاص، فقط یک نورون یا یک ناحیهٔ کاملاً اختصاصی فعال نمی‌شود؛ بلکه تعداد زیادی نورون با شدت‌های متفاوت درگیر می‌شوند. هنگام حرکت به چپ، نورون‌هایی که با جهت چپ سازگارترند فعالیت بیشتری دارند، اما نورون‌های دیگر هم می‌توانند با شدت کمتر فعال باشند.

مزیت نخست \eng{Population Coding} پایداری است. اگر یک نورون آسیب ببیند، کل مفهوم یا فرمان از بین نمی‌رود، چون اطلاعات در یک جمعیت بزرگ‌تر توزیع شده است. مزیت دوم دقت بالاتر است؛ چون تصمیم جمعی می‌تواند از تصمیم یک واحد منفرد پایدارتر باشد، هرچند این گزاره مطلق نیست. مزیت سوم ظرفیت محاسباتی بیشتر است: به‌جای اینکه برای هر زاویه، رنگ یا شیء یک نورون اختصاصی داشته باشیم، ترکیب فعالیت نورون‌ها می‌تواند حالت‌های بسیار متنوعی را بسازد.

ضعف اصلی این نظریه مصرف انرژی و تفسیرپذیری پایین‌تر است. اگر جمعیت بزرگی از نورون‌ها در بسیاری از تصمیم‌ها مشارکت کنند، هزینهٔ انرژی بالا می‌رود. از طرف دیگر، دیگر نمی‌توان گفت یک نورون مشخص مسئول یک خروجی مشخص است. به همین دلیل، \eng{Population Coding} از نظر تفسیرپذیری به یک \eng{black box} نزدیک می‌شود.

\subsubsection{\eng{Sparse Coding}}
\eng{Sparse Coding} را می‌توان حالتی میانی میان دو نظریهٔ قبلی دانست. این دیدگاه حدوداً از دههٔ ۱۹۹۰ پررنگ‌تر شد. در \eng{Sparse Coding}، نه همهٔ اطلاعات در اختیار یک سلول یا گروه بسیار کوچک است، و نه همهٔ نورون‌ها در همهٔ وظایف فعال می‌شوند. برای هر وظیفه، فقط زیرمجموعه‌ای نسبتاً کوچک و مرتبط از نورون‌ها فعال می‌شود.

در این حالت، ماتریس فعالیت نورونی تنک است. ممکن است برای یک کار خاص فقط ۲۰ درصد نورون‌ها، یا حتی ۵ درصد آن‌ها، فعال باشند و بقیه خاموش بمانند. این روش هم مصرف انرژی را کاهش می‌دهد و هم تفسیرپذیری را نسبت به \eng{Population Coding} بهتر می‌کند.

\eng{Sparse Coding} با مفهوم سلسله‌مراتب نیز ارتباط دارد. در پردازش تصویر، لایه‌های ابتدایی ممکن است لبه‌ها را تشخیص دهند، لایه‌های میانی رنگ، شکل یا ترکیب‌های پیچیده‌تر را بازنمایی کنند، و لایه‌های انتهایی به مفهوم یا شیء نهایی نزدیک شوند. در چنین معماری‌ای، گروه‌های مختلف نورونی در سطح‌های متفاوتی از انتزاع فعال می‌شوند.

مزیت \eng{Sparse Coding} این است که از یک طرف مانند \eng{Population Coding} به یک نورون منفرد وابسته نیست و از طرف دیگر مانند آن پرمصرف و کم‌تفسیر نیست. به همین دلیل، می‌توان آن را ترکیبی از پایداری نسبی و تفسیرپذیری نسبی دانست.

\subsection{ترکیب نظریه‌های کدگذاری در مدل‌های امروزی}
در مدل‌های محاسباتی امروز، بهتر است این سه نظریه را رقیب‌های مطلق ندانیم. در بسیاری از معماری‌ها می‌توان هر سه سطح را هم‌زمان دید:

\begin{align*}
\text{لایه‌های ابتدایی} &\longrightarrow \text{\eng{Population Coding}}\\
\text{لایه‌های میانی} &\longrightarrow \text{\eng{Sparse Coding}}\\
\text{لایه‌های انتهایی} &\longrightarrow \text{\eng{Grandmother-like Representation}}
\end{align*}

در لایه‌های ابتدایی، اطلاعات خام‌تر و کم‌انتزاع‌تر است. در پردازش صدا، این بخش می‌تواند به ویژگی‌هایی مانند واج، فرکانس یا الگوهای آوایی مربوط باشد. در بینایی نیز نواحی اولیه مانند \eng{V1} تا \eng{V4} ویژگی‌هایی مثل لبه، جهت، رنگ و فرم‌های ابتدایی را پردازش می‌کنند. این نمایش‌ها قدرتمندند، اما تفسیر هر بعد آن‌ها دشوار است؛ شبیه \eng{embedding}های مدل‌های یادگیری عمیق.

در لایه‌های میانی، ویژگی‌ها مشخص‌تر می‌شوند. ابتدا لبه‌ها و ویژگی‌های ساده‌تر دیده می‌شوند، سپس شکل‌ها و ترکیب‌ها و در ادامه ویژگی‌هایی که به اشیا یا مفاهیم نزدیک‌ترند. این سطح به \eng{Sparse Coding} نزدیک‌تر است، چون همهٔ واحدها برای همهٔ ورودی‌ها فعال نیستند.

در لایه‌های انتهایی، خروجی معمولاً به کلاس‌ها یا مفاهیم مشخص نگاشت می‌شود. اگر مدل چند کلاس داشته باشد، هر نود خروجی تا حدی مسئول یک کلاس است. این وضعیت شبیه یک بازنمایی \eng{Grandmother-like} است، با این تأکید که این شباهت به معنای پیاده‌سازی دقیق نظریهٔ زیستی \eng{Grandmother Cell} نیست.

\subsection{نسبت \eng{Symbolic AI} با نظریه‌های کدگذاری}
در این بحث این پرسش مطرح شد که اگر بخواهیم \eng{Symbolic AI} را با نظریه‌های کدگذاری عصبی مقایسه کنیم، به کدام نظریه نزدیک‌تر است. پاسخ اولیهٔ یکی از خوانندگان \eng{Grandmother Cell} بود، اما توضیح این بود که \eng{Symbolic AI} بیشتر به \eng{Sparse Coding} نزدیک است.

دلیل این مقایسه آن است که در \eng{Symbolic AI} معمولاً با سیستم‌های قانون‌محور، ویژگی‌های استخراج‌شده و ساختارهای نمادین سروکار داریم. تصمیم نهایی نه کاملاً توزیع‌شده و همه‌گیر مانند \eng{Population Coding} است، و نه کاملاً انحصاری مانند \eng{Grandmother Cell}. مجموعه‌ای از قواعد یا ویژگی‌های مشخص در تصمیم دخالت می‌کنند؛ بنابراین شباهت آن با \eng{Sparse Coding} بیشتر است.

\subsection{\eng{Brain-to-Text}}
بخش بعدی به مقاله‌ای از \eng{Meta} در سال ۲۰۲۵ دربارهٔ \eng{Brain-to-QWERTY} یا \eng{Brain-to-Text} اختصاص داشت. در چارچوب کلی \eng{Brain-to-Text} هدف این است که از فعالیت مغزی به متن برسیم، اما در \eng{Brain-to-QWERTY} دامنه محدودتر می‌شود: سیستم تلاش می‌کند بفهمد فرد هنگام تایپ چه کلیدهایی را روی صفحه‌کلید \eng{QWERTY} فشار داده یا قصد فشار دادن آن‌ها را داشته است.

پس هدف این مقاله خواندن تمام افکار فرد نیست. دامنهٔ مسئله به سناریوی تایپ محدود شده است، چون هرچه دامنه محدودتر باشد، \eng{decoding} فعالیت مغزی عملی‌تر می‌شود.

\subsubsection{مرحلهٔ آموزش}
در مرحلهٔ آموزش، فرد در دستگاه قرار می‌گیرد و فعالیت مغزی او هنگام خواندن و تایپ ثبت می‌شود. در این توضیح، روش‌هایی مانند \eng{EEG} و \eng{MEG} یا سیگنال‌های مشابه به‌عنوان مسیرهای ممکن ثبت فعالیت مغزی مطرح شدند.

روند کلی چنین است:

\begin{enumerate}
  \item متنی به فرد نشان داده می‌شود.
  \item فرد ابتدا متن را می‌خواند.
  \item سپس از او خواسته می‌شود متن را با صفحه‌کلید \eng{QWERTY} تایپ کند.
  \item هنگام تایپ، سیگنال‌های مغزی ثبت می‌شوند.
  \item سیستم می‌داند متن ورودی چه بوده و فرد قرار بوده چه چیزی تایپ کند.
  \item مدل یاد می‌گیرد فعالیت مغزی را به کلیدها یا متن تایپ‌شده نگاشت کند.
\end{enumerate}

در این مرحله، خروجی مدل با متن اصلی یا کلیدهای مورد انتظار مقایسه می‌شود و \eng{loss} محاسبه می‌گردد. سپس وزن‌های بخش‌های مختلف مدل بر اساس همین خطا به‌روزرسانی می‌شوند. این فرایند معمولاً برای همان فرد تکرار می‌شود تا مدل روی سیگنال‌های مغزی او کالیبره شود.

معماری کلی که در این بحث توصیف شد چنین بود:

\begin{align*}
\text{سیگنال مغزی هنگام تایپ} &\longrightarrow \text{\eng{CNN}}\\
&\longrightarrow \text{ویژگی‌های میانی}\\
&\longrightarrow \text{\eng{Transformer Decoder-only}}\\
&\longrightarrow \text{\eng{Language Model}}\\
&\longrightarrow \text{متن بازسازی‌شده}
\end{align*}

\coursefigure{0.95}{brain-to-qwerty-pipeline.png}{پایپ‌لاین \eng{Brain-to-QWERTY}: فرد جمله را می‌خواند و تایپ می‌کند، سیگنال‌های \eng{MEG/EEG} ثبت می‌شوند و مدل شامل ماژول کانولوشنی، ترنسفورمر و مدل زبان، متن هدف را بازسازی می‌کند.}

\subsubsection{نقش \eng{CNN}}
سیگنال مغزی خام ابتدا به یک \eng{Convolutional Neural Network} داده می‌شود. نقش \eng{CNN} استخراج ویژگی از سیگنال خام است. یعنی دادهٔ اولیهٔ مغزی به یک \eng{representation} یا مجموعه‌ای از ویژگی‌های میانی تبدیل می‌شود تا بخش‌های بعدی مدل بتوانند با آن کار کنند.

\subsubsection{نقش \eng{Transformer}}
ویژگی‌های استخراج‌شده به یک \eng{Transformer} داده می‌شوند. این \eng{Transformer} در این توضیح از نوع \eng{decoder-only} معرفی شد. وظیفهٔ آن مدل‌کردن توالی زبانی یا توالی کلیدهای تایپ‌شده از روی ویژگی‌های مغزی است؛ برای مثال تشخیص دهد فرد در حال تایپ چه حرف‌ها یا چه کلمه‌هایی بوده است.

\subsubsection{نقش \eng{Language Model}}
\eng{Language Model} در انتهای مسیر می‌تواند خطاهای خروجی را اصلاح کند. اگر بخش‌های قبلی مدل خروجی‌ای شبیه «صلام» تولید کنند، مدل زبانی با توجه به زمینه می‌تواند آن را به «سلام» نزدیک کند. بنابراین نقش آن شبیه یک تصحیح‌گر کانتکستی یا \eng{autocorrect} است.

\subsubsection{فاز آزمون}
در فاز آزمون، متن ورودی برای مدل معلوم نیست. فرد چیزی را تایپ می‌کند یا قصد تایپ آن را دارد و مدل فقط از روی فعالیت مغزی تلاش می‌کند بفهمد او چه نوشته یا قصد نوشتن چه چیزی داشته است. هدف نهایی این است که از سیگنال مغز بتوان به متن، فرمان یا نیت فرد رسید.

\subsection{کاربرد \eng{Brain-to-Text} و مقایسه با \eng{Neuralink}}
کاربرد مثبت چنین سیستم‌هایی کمک به افرادی است که ذهن فعال دارند اما توانایی حرکت یا ارتباط فیزیکی ندارند؛ برای مثال افرادی که نمی‌توانند تایپ کنند، دچار آسیب نخاعی‌اند یا محدودیت شدید حرکتی دارند. اگر فعالیت مغزی آن‌ها \eng{decode} شود، نیت یا دستور ذهنی می‌تواند به متن یا فرمان خارجی تبدیل شود و اجرای فرمان به سیستم‌های دیگر سپرده شود.

در همین زمینه، \eng{Neuralink} نیز نمونهٔ مهمی از ایمپلنت مغزی است که با جراحی در مغز قرار می‌گیرد و هدف آن تبدیل فعالیت مغزی، به‌ویژه نیت حرکتی، به فرمان خارجی است. برای مثال فرد به حرکت نشانگر ماوس فکر می‌کند و سیستم تلاش می‌کند همان نیت را به حرکت ماوس تبدیل کند.

تفاوت مهم این است که \eng{Brain-to-QWERTY} دامنهٔ محدودتری دارد و به تایپ روی صفحه‌کلید مربوط است، اما در \eng{Neuralink} هدف می‌تواند مستقیم‌تر باشد: تبدیل فکر یا نیت حرکتی به حرکت ماوس، کلیک یا کنترل یک ابزار.

\subsection{مقایسهٔ روش‌های ثبت فعالیت مغزی}
یکی از چالش‌های اصلی \eng{Brain-to-Text} محدودیت سخت‌افزار و روش ثبت فعالیت مغزی است. در این بحث سه دسته روش مقایسه شد: \eng{fMRI}، \eng{EEG} و روش‌های تهاجمی مانند ایمپلنت‌های مغزی.

\subsubsection{\eng{fMRI}}
\eng{fMRI} بر اساس تغییرات جریان خون کار می‌کند. وقتی بخشی از مغز فعال می‌شود، جریان خون در آن بخش تغییر می‌کند و \eng{fMRI} این تغییر را ثبت می‌کند.

مزیت \eng{fMRI} رزولوشن مکانی خوب آن است؛ یعنی می‌توان نسبتاً دقیق‌تر دید کدام ناحیهٔ مغز فعال شده است. ضعف آن رزولوشن زمانی پایین است، چون تغییرات جریان خون با تأخیر رخ می‌دهد. نکتهٔ مهم این است که \eng{fMRI} در مقیاس ثانیه معنا دارد، نه میلی‌ثانیه.

\subsubsection{\eng{EEG}}
\eng{EEG} فعالیت الکتریکی مغز را از روی پوست سر ثبت می‌کند. مزیت آن رزولوشن زمانی خوب است، چون تغییرات سریع الکتریکی را می‌تواند در حد میلی‌ثانیه ثبت کند. ضعف آن رزولوشن مکانی پایین‌تر است؛ سیگنال باید از بافت‌ها، جمجمه و پوست سر عبور کند و بنابراین تشخیص دقیق محل تولید سیگنال دشوارتر می‌شود.

\subsubsection{ایمپلنت‌های مغزی}
در روش‌های تهاجمی، حسگر یا تراشه با جراحی وارد مغز می‌شود. چون سنسور به نورون‌ها نزدیک‌تر است، هم رزولوشن زمانی خوب است و هم رزولوشن مکانی نسبت به روش‌های غیرتهاجمی بهتر می‌شود. اما مشکل اصلی، تهاجمی بودن این روش‌هاست: نیاز به جراحی، مسائل ایمنی، پایداری بلندمدت و ملاحظات اخلاقی.

بنابراین مقایسهٔ کلی چنین است:

\begin{itemize}
  \item \eng{fMRI}: غیرتهاجمی، رزولوشن مکانی خوب، رزولوشن زمانی پایین، مبتنی بر جریان خون.
  \item \eng{EEG}: غیرتهاجمی، رزولوشن زمانی خوب، رزولوشن مکانی پایین‌تر، مبتنی بر فعالیت الکتریکی از روی پوست سر.
  \item ایمپلنت‌های مغزی: تهاجمی، با رزولوشن زمانی و مکانی بهتر. نمونه‌ای مانند \eng{Neuralink} نیازمند جراحی است و چالش‌های پزشکی و اخلاقی دارد.
\end{itemize}

\subsection{چالش‌های \eng{Brain-to-Text}}
سه چالش اصلی که در این بحث برجسته شد عبارت‌اند از \eng{plasticity} مغز، تعمیم‌پذیری بین افراد و چندوجهی‌بودن فعالیت مغزی. این چالش‌ها باعث می‌شوند ساخت یک مدل عمومی و پایدار برای \eng{Brain-to-Text} بسیار دشوارتر از یک مسئلهٔ سادهٔ طبقه‌بندی باشد.

\subsubsection{\eng{Plasticity}}
مغز انسان ثابت نیست. هر فرد در طول زمان یاد می‌گیرد، فراموش می‌کند، خسته می‌شود و وضعیت ذهنی‌اش تغییر می‌کند. بنابراین مدلی که امروز روی سیگنال‌های مغزی یک فرد کالیبره شده، الزاماً در روزهای بعد دقیقاً همان ورودی‌ها را نمی‌بیند.

مسئله را می‌توان چنین توضیح داد: ما مدلی را روی مغز یک فرد آموزش می‌دهیم، اما خود آن مغز هم هم‌زمان در حال آموزش دیدن و تغییر کردن است. بنابراین مدل باید بتواند با \eng{plasticity} مغز کنار بیاید.

\subsubsection{تعمیم‌پذیری بین افراد}
نقشهٔ مغزی هر فرد با فرد دیگر متفاوت است و می‌توان آن را به اثر انگشت تشبیه کرد. مدلی که روی یک فرد آموزش دیده الزاماً روی فرد دیگر خوب کار نمی‌کند. در سامانه‌هایی مانند \eng{Neuralink} یا \eng{Brain-to-QWERTY} معمولاً مدل برای یک فرد خاص آموزش و آزمون می‌شود.

چالش اصلی این است که آیا می‌توان یک \eng{foundation model} برای مغز ساخت که روی افراد مختلف تعمیم پیدا کند یا نه. در مدل‌های زبانی تا حدی به چنین مدل‌هایی رسیده‌ایم، اما در مغز مسئله پیچیده‌تر است، چون باید فعالیت مغزی به زبان، حرکت، تصویر یا فرمان نگاشت شود.

\subsubsection{\eng{Multimodality}}
فعالیت مغزی ذاتاً چندوجهی است. حتی اگر آزمایش فقط روی متن تمرکز کند، مغز هم‌زمان از تصویر، صدا، حافظه، توجه، خستگی، احساسات و عوامل دیگر تأثیر می‌گیرد. به همین دلیل در آزمایش‌ها تلاش می‌شود دامنه محدود شود؛ برای مثال فقط متن نشان داده شود و فرد فقط تایپ کند.

اما هرچه سیستم به شرایط واقعی نزدیک‌تر شود، چندوجهی‌تر می‌شود و \eng{decoding} دشوارتر خواهد بود.

\subsection{مسیر آینده: سنسورهای پوشیدنی و سخت‌افزار \eng{Neuromorphic}}
مسیر ایده‌آل این حوزه حرکت به سمت سنسورهایی است که غیرتهاجمی، پوشیدنی، کم‌مصرف و دقیق باشند؛ یعنی هم رزولوشن زمانی خوب داشته باشند، هم رزولوشن مکانی مناسب، هم برای بدن قابل تحمل باشند و هم گرمای زیادی تولید نکنند. در این بحث، سنسورهای کوانتومی به‌عنوان نمونه‌ای از مسیرهای مورد توجه مطرح شدند.

مسئلهٔ انرژی در اینجا حیاتی است. اگر سخت‌افزاری نزدیک بدن یا داخل بدن قرار بگیرد، مصرف انرژی زیاد می‌تواند گرما تولید کند و برای بافت مغز خطرناک باشد. بنابراین سخت‌افزار باید هم دقیق و هم کم‌مصرف باشد.

یکی از ایده‌های مهم برای این هدف، سخت‌افزار \eng{neuromorphic} است. این سخت‌افزارها از ساختار و عملکرد مغز الهام می‌گیرند و معمولاً \eng{event-driven} هستند. در سیستم‌های \eng{time-based}، محاسبه بر اساس کلاک انجام می‌شود؛ چه رویدادی رخ داده باشد و چه رخ نداده باشد. اما در سیستم‌های \eng{event-driven}، فقط وقتی رویدادی اتفاق بیفتد پردازش انجام می‌شود. در شبکه‌های اسپایکی، این رویداد همان اسپایک است:

\[
\text{وقتی \eng{spike} رخ دهد} \longrightarrow \text{پردازش انجام می‌شود}
\]

چون مغز نیز تا حد زیادی رویدادمحور عمل می‌کند، سخت‌افزارهای \eng{neuromorphic} برای این حوزه گزینه‌ای طبیعی به نظر می‌رسند. این نکته به‌ویژه در \eng{Brain-to-Text} و ایمپلنت‌های مغزی مهم است.

\subsection{ترجیح مکانی، قانون \eng{Hebb} و مقداردهی اولیه}
در پایان این بخش، نکتهٔ مهم این است که در مدل‌های عصبی خروجی معمولاً از ترکیب وزن‌ها و فعالیت‌ها ساخته می‌شود:

\[
\sum_i W_i X_i
\]

اما این وزن‌ها فقط از یک عامل نمی‌آیند. سه عامل مهم در شکل‌گیری وزن نهایی مطرح شد: \eng{Hebbian Learning}، ترجیح مکانی و مقداردهی اولیه.

طبق قانون \eng{Hebb}، نورون‌هایی که با هم فعال می‌شوند، ارتباطشان تقویت می‌شود. این قانون در شکل‌گیری وزن‌های سیناپسی نقش دارد. علاوه بر آن، موقعیت مکانی یا ترجیح عملکردی نورون‌ها نیز مهم است. برای مثال، در آزمایش حرکت میمون، هنگام حرکت به سمت چپ، نورون‌هایی که از نظر مکانی یا عملکردی به جهت چپ نزدیک‌ترند بیشتر فعال می‌شوند.

عامل سوم مقداردهی اولیه است. افراد از ابتدا نیز تفاوت‌هایی در وزن‌ها یا تنظیمات اولیهٔ مغزی دارند و همین تفاوت‌ها در پاسخ نهایی نقش پیدا می‌کند. بنابراین وزن نهایی را می‌توان حاصل ترکیب این عوامل دانست:

\[
\text{وزن نهایی}
=
\text{\eng{Hebbian Learning}}
+
\text{ترجیح مکانی}
+
\text{مقداردهی اولیه}
\]

جمع‌بندی این بخش این است که سه نظریهٔ \eng{Grandmother Cell}، \eng{Population Coding} و \eng{Sparse Coding} سه نگاه متفاوت به بازنمایی عصبی ارائه می‌کنند. مدل‌های امروزی را می‌توان ترکیبی از این سه سطح دید: \eng{embedding}های اولیه به \eng{Population Coding} نزدیک‌اند، ویژگی‌ها و \eng{hidden state}های میانی به \eng{Sparse Coding}، و لایهٔ خروجی به بازنمایی \eng{Grandmother-like}. در ادامه، \eng{Brain-to-QWERTY} به‌عنوان یک نمونهٔ عملی از \eng{decoding} مغزی بررسی شد و چالش‌های آن، از تغییرپذیری مغز تا سخت‌افزار و انرژی، روشن شد.

\section{\eng{Brain Alignment} و \eng{Brain Score}}

\subsection{مسئلهٔ بازنمایی}
بحث \eng{Brain Alignment} پس از مباحث \eng{Brain-to-QWERTY} و \eng{Brain Decoding} مطرح شد. در بخش قبل دیدیم که می‌توان از روی سیگنال‌های مغزی، تا حدی قصد تایپ یا فرمان ذهنی فرد را بازسازی کرد. این کاربردها از یک طرف برای کمک به افراد دارای محدودیت حرکتی بسیار مهم‌اند، و از طرف دیگر نگرانی‌های جدی دربارهٔ ذهن‌خوانی، سوءاستفاده از داده‌های ذهنی و افشای ناخواستهٔ افکار ایجاد می‌کنند.

برای ورود به \eng{Brain Alignment}، مسئله را می‌توان از زاویهٔ بازنمایی توضیح داد. در یک نگاه فلسفی، می‌توان میان ماده، صورت یا \eng{Form}، و سطحی مجردتر تمایز گذاشت. آنچه در جهان مادی می‌بینیم، بازنمایی یا \eng{representation} چیزی عمیق‌تر است.

در همین چارچوب، خروجی یک مدل زبانی، سیگنال \eng{fMRI} هنگام فکر کردن، تصویر، متن، صدا یا رفتار انسان را می‌توان بازنمایی‌هایی متفاوت از یک معنا یا حقیقت واحد دانست. ممکن است چند بازنمایی متفاوت داشته باشیم، اما همهٔ آن‌ها به یک مفهوم مشترک اشاره کنند. مسئلهٔ اصلی این است که آیا بازنمایی‌های مدل و بازنمایی‌های مغز انسان به یکدیگر نزدیک‌اند یا نه.

\coursefigure{0.72}{platonic-representation-hypothesis.png}{فرضیهٔ بازنمایی افلاطونی: تصویر، متن و فعالیت مغزی را می‌توان به‌عنوان تصویرهایی متفاوت از یک واقعیت یا مفهوم زیرین مشترک در نظر گرفت.}

\subsection{آیا مدل واقعاً می‌فهمد؟}
نکتهٔ مهم این است که نمی‌توان به‌صورت صددرصدی و الگوریتمی ثابت کرد یک مدل «حقیقت» را فهمیده است یا نه. این مسئله تا حدی شبیه مسئلهٔ \eng{hallucination} است: هیچ الگوریتم قطعی‌ای وجود ندارد که تضمین کند مدل هیچ‌وقت دچار \eng{hallucination} نمی‌شود، چون درست یا غلط بودن پاسخ به زمینه وابسته است.

یک جمله ممکن است در زمینه‌ای تخیلی درست باشد، در زمینه‌ای پزشکی غلط باشد، یا از نگاه افراد مختلف تفسیرهای متفاوتی پیدا کند. بنابراین «فهمیدن» یا «نفهمیدن» مدل در معنای مطلق، مسئله‌ای کاملاً حل‌شده نیست.

با این حال، می‌توان بررسی کرد که مدل تا چه حد از پوستهٔ ظاهری داده فاصله گرفته و به معنا نزدیک‌تر شده است. مثال، سیب بود. اگر مدل فقط یاد گرفته باشد سیب معمولاً گرد، قرمز یا زرد است، ممکن است با دیدن یک سیب فاسد یا کپک‌زده دچار مشکل شود. اما انسان حتی اگر قبلاً سیب کپک‌زده ندیده باشد، مفهوم «فاسد شدن» را می‌فهمد و می‌تواند بگوید هنوز با یک سیب، اما یک سیب فاسد، روبه‌روست.

پس فهم انسانی صرفاً حفظ ظاهر نیست. مدل نیز اگر واقعاً به معنا نزدیک شده باشد، باید بتواند از ویژگی‌های سطحی عبور کند و ساختار انتزاعی‌تری از مفهوم را بازنمایی کند.

\subsection{ایدهٔ اصلی \eng{Brain Alignment}}
هدف \eng{Brain Alignment} این است که بررسی کنیم آیا بازنمایی‌های مدل با بازنمایی‌های مغز انسان هم‌راستا هستند یا نه. برای قضاوت دربارهٔ کیفیت یک مدل، فقط مقایسهٔ خروجی نهایی کافی نیست. باید ببینیم لایه‌های میانی مدل چه چیزی را \eng{encode} کرده‌اند، مغز انسان هنگام همان \eng{task} چه بازنمایی‌هایی دارد، و آیا این دو مجموعه بازنمایی قابل نگاشت به هم هستند یا نه.

اگر مدلی فقط خروجی درست بدهد، ممکن است شبیه طوطی الگوهایی را حفظ کرده باشد. اما اگر بازنمایی‌های میانی‌اش با بازنمایی‌های انسانی هم‌راستا باشند، احتمال بیشتری دارد که مدل به سطحی عمیق‌تر از معنا نزدیک شده باشد.

همچنین از روایت‌های تجربه‌های نزدیک به مرگ برای توضیح دشواری نگاشت میان سطح‌های مختلف بازنمایی استفاده کرد. در چنین روایت‌هایی، فرد گاهی تجربه‌ای غیرمادی را با مفاهیم مادی توضیح می‌دهد؛ مثلاً می‌گوید «درخت‌هایی بودند که حرف می‌زدند». یعنی تجربه را با واژگان جهان مادی بیان می‌کند، اما همان واژه‌ها ویژگی‌هایی پیدا می‌کنند که در مادهٔ عادی وجود ندارد. این مثال نشان می‌دهد نگاشت میان بازنمایی‌ها همیشه ساده و یک‌به‌یک نیست.

\subsection{طراحی آزمایش هم‌ترازی}
سؤال عملی این است: چگونه بفهمیم مدل از سطح ظاهر عبور کرده و به سطح معنا نزدیک شده است؟

مقایسهٔ مدل با انسان یک مسیر طبیعی است، اما اگر انسان را \eng{ground truth} بگیریم، باید مشخص کنیم چه چیزی را از انسان جمع‌آوری می‌کنیم. در بسیاری از آزمایش‌ها تا امروز، خروجی انسان و مدل مقایسه شده است: انسان \eng{task} را انجام می‌دهد، مدل هم همان \eng{task} را انجام می‌دهد، و خروجی‌ها مقایسه می‌شوند. اما برای \eng{Brain Alignment} فقط خروجی کافی نیست؛ بازنمایی‌های میانی نیز باید بررسی شوند.

یک مسیر دیگر استفاده از بازنمایی چندوجهی است. اگر حقیقت واحدی مانند مفهوم سیب از مسیرهای مختلفی مثل تصویر، متن، صدا و رفتار انسانی بررسی شود، و مدل بتواند در همهٔ این \eng{modality}ها همان مفهوم را تشخیص دهد، احتمالاً بازنمایی عمیق‌تری ساخته است.

مسیر سوم انتخاب \eng{task}هایی است که معنایی‌تر و دورتر از پوستهٔ مادی‌اند؛ مثل \eng{task}های فلسفی، استدلالی، علی و معلولی، فیزیکی، کشف قانون، استنتاج و قضاوت در وضعیت‌های جدید.

مثال نیوتون و افتادن سیب این پرسش را روشن می‌کند: اگر مدلی را زیر درخت بگذاریم و قبلاً قانون جاذبه را به آن نگفته باشیم، آیا می‌تواند از افتادن سیب تعجب کند و به وجود یک قانون برسد؟ در چنین \eng{task}هایی مطلوب است که مدل پاسخ را از قبل ندیده باشد، قانون به‌صورت صریح به آن داده نشده باشد، و مجبور شود رابطهٔ علی و معلولی بسازد.

هدف نهایی \eng{Brain Alignment} اثبات قطعی فهم مدل نیست. هدف عملی‌تر این است که بفهمیم مدل کجا را نمی‌فهمد تا بتوانیم آن را اصلاح کنیم. مسئلهٔ سخت «آیا مدل حقیقت را فهمیده است؟» به مسئلهٔ قابل‌بررسی‌تری کاهش داده می‌شود: دو جعبهٔ سیاه داریم، مغز انسان و مدل هوش مصنوعی؛ خروجی‌هایشان تا حدی مشابه است. حالا می‌خواهیم ببینیم لایه‌های میانی‌شان نیز با هم هم‌راستا هستند یا نه.

\subsubsection{تعریف لایه‌های میانی مغز}
در مدل عصبی مصنوعی، لایه‌ها و \eng{hidden state}ها مشخص‌اند. اما در مغز باید تصمیم بگیریم که «لایهٔ میانی» را چگونه تعریف کنیم. این پرسش به مسئلهٔ \eng{Brain Parcellation} می‌رسد: مغز را چگونه به بخش‌ها یا \eng{parcel}های مختلف تقسیم کنیم تا بتوان بازنمایی‌های آن را با مدل مقایسه کرد؟

\subsubsection{انتخاب \eng{Task}}
باید مشخص شود چه \eng{task}هایی برای سنجش هم‌ترازی مناسب‌اند. \eng{task}های زبانی، پیش‌بینی، استدلالی، معنایی و چندوجهی می‌توانند گزینه باشند. هرچه \eng{task} فقط به خروجی سطحی وابسته نباشد و مدل را مجبور به ساختن بازنمایی معنایی کند، برای بررسی \eng{alignment} مناسب‌تر است.

\subsubsection{تعریف \eng{Score}}
پس از جمع‌آوری داده از مغز و مدل، باید مشخص شود هم‌ترازی چگونه اندازه‌گیری می‌شود. معیار می‌تواند \eng{similarity}، \eng{correlation}، \eng{loss} یا ترکیبی از چند معیار باشد. مهم این است که یک معیار عملی داشته باشیم که نشان دهد بازنمایی مدل تا چه حد می‌تواند دادهٔ مغزی یا رفتاری انسان را پیش‌بینی کند.

\subsection{مقالهٔ \eng{Brain Score}}
\eng{Brain Score} به‌عنوان نمونه‌ای از این خط پژوهشی مطرح شد. تمرکز این مقاله بیشتر روی سؤال سوم است: چگونه برای \eng{alignment} یک \eng{benchmark} و \eng{score} تعریف کنیم؟

ایدهٔ کلی آزمایش چنین است: انسان داخل دستگاه \eng{fMRI} قرار می‌گیرد، یک محرک مثل متن به او داده می‌شود، همان متن به مدل نیز داده می‌شود، از مغز انسان سیگنال \eng{fMRI} گرفته می‌شود، و از مدل بازنمایی‌های میانی یا \eng{hidden vectors} استخراج می‌گردد. سپس یک مدل ساده یاد می‌گیرد بازنمایی مدل را به فعالیت مغزی نگاشت کند. در مرحلهٔ آزمون بررسی می‌شود این نگاشت تا چه حد می‌تواند فعالیت مغز را پیش‌بینی کند.

\subsubsection{اجزای آزمایش}
در این آزمایش دو جعبهٔ سیاه داریم: مغز انسان و مدل، مثلاً یک \eng{LLM}. ورودی یا \eng{stimulus} می‌تواند متن باشد.

برای مدل، متن وارد شبکه می‌شود و از لایه‌های میانی بردارهایی مانند \(v_1\)، \(v_2\) و \(v_3\) گرفته می‌شود. مدل اصلی \eng{freeze} می‌ماند، چون هدف ارزیابی آن است، نه آموزش دوبارهٔ آن.

برای انسان، همان متن نمایش داده می‌شود و از مغز او \eng{fMRI} گرفته می‌شود. ممکن است رفتارهایی مثل سرعت خواندن نیز ثبت شوند. سپس یک مدل کوچک و ساده آموزش داده می‌شود تا نگاشت میان بازنمایی‌های مدل و سیگنال‌های مغزی یا رفتاری را یاد بگیرد.

\subsubsection{مدل \eng{Mapping}}
مدل اصلی، مثلاً \eng{LLM}، دست‌کاری نمی‌شود. اما یک مدل کمکی آموزش داده می‌شود. وظیفهٔ این مدل کمکی این است که اگر بردار میانی مدل را دریافت کند، پیش‌بینی کند کدام الگوی \eng{fMRI} متناظر با آن است؛ یا اگر \eng{confidence} مدل در \eng{next token prediction} را دریافت کند، آن را به رفتار انسانی مثل سرعت خواندن نگاشت کند.

این مدل کمکی معمولاً می‌تواند خطی باشد. هدف ساختن یک مدل پیچیدهٔ جدید نیست؛ هدف این است که ببینیم آیا بازنمایی‌های موجود مدل با مغز قابل نگاشت هستند یا نه.

یکی از ایده‌های مهم، مقایسهٔ رفتار انسان با \eng{confidence} مدل است. اگر جمله برای انسان قابل پیش‌بینی و روان باشد، سریع‌تر خوانده می‌شود. اگر جمله غیرمنتظره باشد، انسان مکث می‌کند. برای مدل هم اگر \eng{next token} قابل پیش‌بینی باشد، \eng{confidence} بالاتر است؛ و اگر عبارت غیرمنتظره باشد، \eng{confidence} پایین‌تر می‌آید. بنابراین می‌توان \eng{confidence} مدل را با سرعت خواندن انسان مرتبط کرد.

نمونه این بود که «پردازش زبان طبیعی» عبارتی آشناست و ادامهٔ آن آسان‌تر پیش‌بینی می‌شود، اما «پردازش زبان فیزیک» عجیب‌تر است و می‌تواند باعث مکث شود.

\subsubsection{مرحلهٔ آموزش و آزمون}
در مرحلهٔ آموزش، متن‌هایی به انسان و مدل داده می‌شود. بازنمایی‌های مدل استخراج می‌گردد، \eng{fMRI} و رفتار انسان ثبت می‌شود، و مدل \eng{mapping} یاد می‌گیرد این‌ها را به هم وصل کند.

در مرحلهٔ آزمون، متن جدید به مدل داده می‌شود. مدل \eng{mapping} از روی بازنمایی مدل پیش‌بینی می‌کند که \eng{fMRI} یا رفتار انسانی چه خواهد بود. سپس این پیش‌بینی با دادهٔ واقعی انسان مقایسه می‌شود و از این مقایسه \eng{score} نهایی ساخته می‌شود.

\eng{Brain Score} می‌تواند از چند مؤلفه تشکیل شود؛ برای مثال هم‌ترازی با \eng{fMRI}، هم‌ترازی با خروجی یا پیش‌بینی، و هم‌ترازی با رفتارهایی مثل \eng{reading time}. در این توضیح، معیار اصلی برای مقایسه می‌تواند \eng{correlation} باشد.

\coursefigure{0.95}{brain-score-framework.png}{چارچوب \eng{Brain Score}: محرک زبانی مشترک به مدل و انسان داده می‌شود، سپس بازنمایی‌های داخلی مدل با پاسخ‌های مغزی و رفتاری انسان مقایسه می‌شوند.}

\subsubsection{اهمیت \eng{Next Word Prediction}}
در \eng{Brain Score} چند \eng{task} بررسی شده‌اند؛ از جمله \eng{question answering}، \eng{coherence}، \eng{entailment}، \eng{sentence similarity} و \eng{next word prediction}. نکتهٔ جالب این است که همبستگی مثبت و معنادار بیشتر در \eng{task} مربوط به \eng{next word prediction} دیده شده است.

این نتیجه با مباحث قبلی ترم سازگار است: مغز انسان دائماً در حال پیش‌بینی است. بنابراین \eng{task}هایی که حالت \eng{predictive} دارند، احتمالاً برای \eng{Brain Alignment} مناسب‌ترند. \eng{Task}هایی که صرفاً به خروجی نهایی مربوط‌اند، ممکن است بیشتر لایهٔ آخر را درگیر کنند؛ اما \eng{task}های پیش‌بینی‌کننده بازنمایی‌های میانی را بهتر آشکار می‌کنند.

\subsection{\eng{Brain Parcellation}}
\eng{Brain Parcellation} پاسخی به این پرسش است که لایه‌ها یا بخش‌های میانی مغز را چگونه تعریف کنیم. در این روش، مغز یا یک ناحیهٔ مورد علاقه از مغز به واحدهای کوچک‌تر و مجزا تقسیم می‌شود.

\subsubsection{\eng{Voxel}}
در \eng{fMRI}، مغز به واحدهای کوچک مکعبی تقسیم می‌شود که به هرکدام \eng{voxel} گفته می‌شود. \eng{Voxel}ها مرزهای طبیعی و زیستی مغز نیستند؛ آن‌ها واحدهایی‌اند که ما برای اندازه‌گیری تعریف می‌کنیم.

\subsubsection{\eng{Parcel}}
چند \eng{voxel} که در یک گروه قرار می‌گیرند، یک \eng{parcel} می‌سازند:

\[
\text{چند \eng{voxel}} \longrightarrow \text{یک \eng{parcel}}
\qquad
\text{چند \eng{parcel}} \longrightarrow \text{یک \eng{atlas} یا \eng{parcellation}}
\]

\subsubsection{\eng{Region of Interest}}
گاهی لازم نیست کل مغز بررسی شود. بسته به \eng{task}، فقط یک ناحیهٔ خاص انتخاب می‌شود. برای مثال اگر \eng{task} زبانی باشد، می‌توان فقط ناحیه‌های مرتبط با زبان را بررسی کرد. به چنین ناحیه‌ای \eng{Region of Interest} یا \eng{ROI} می‌گوییم.

\subsubsection{سیگنال \eng{BOLD}}
در \eng{fMRI} معمولاً از شاخصی به نام \eng{BOLD} استفاده می‌شود:

\[
\text{\eng{BOLD}} = \text{\eng{Blood Oxygen Level Dependent}}
\]

این شاخص تغییرات سطح اکسیژن خون را نشان می‌دهد. چون فعالیت عصبی با مصرف اکسیژن و جریان خون مرتبط است، تغییرات \eng{BOLD} به‌عنوان نشانه‌ای از فعالیت مغزی استفاده می‌شود.

برای هر \eng{voxel}، در طول زمان یک \eng{time series} داریم. اگر یک \eng{parcel} شامل چند \eng{voxel} باشد، می‌توان برای آن \eng{parcel} یک بردار میانگین یا بازنمایی کلی ساخت. برای مثال، اگر یک \eng{parcel} سه \eng{voxel} داشته باشد و در چهار زمان اندازه‌گیری انجام شود، دادهٔ اولیه یک ماتریس \(3 \times 4\) است که می‌توان آن را به بردار میانگین \eng{parcel} تبدیل کرد.

\subsection{روش‌های \eng{Parcellation}}
دو رویکرد کلی برای \eng{parcellation} مطرح شد: تقسیم‌بندی بر اساس آناتومی و تقسیم‌بندی بر اساس کارکرد.

\subsubsection{بر اساس آناتومی}
در این روش، مغز بر اساس ساختار آناتومیک تقسیم می‌شود. در بینایی، نمونه‌هایی مانند \eng{V1}، \eng{V2}، \eng{V4} و \eng{IT} مطرح‌اند. در اینجا تقسیم‌بندی بر اساس محل و ساختار فیزیکی مغز است.

\subsubsection{\eng{Functional Parcellation}}
در بحث \eng{Brain Alignment}، \eng{functional parcellation} معمولاً مهم‌تر است. در این روش، مهم نیست بخش‌ها از نظر فیزیکی دقیقاً کجا هستند؛ مهم این است که کدام بخش‌ها با هم کار می‌کنند. ممکن است دو نقطه از مغز از نظر مکانی دور باشند، اما چون هنگام یک کارکرد خاص با هم فعال می‌شوند، در یک شبکه قرار بگیرند.

این وضعیت را به کارمندانی تشبیه کرد که در طبقه‌های مختلف یک سازمان هستند، اما برای یک کار مشخص مدام با هم ارتباط دارند. از روی ارتباط کاری می‌فهمیم در یک تیم‌اند، نه صرفاً از روی محل فیزیکی‌شان.

\subsubsection{\eng{Localizer} و \eng{Atlas}}
برای انتخاب ناحیه‌های مغزی دو راه وجود دارد. راه اول استفاده از \eng{atlas}های موجود است؛ یعنی \eng{parcellation}هایی که قبلاً در مقاله‌های دیگر ساخته شده‌اند. پژوهشگر می‌گوید مبنای من فلان \eng{atlas} معروف است.

راه دوم ساختن \eng{localizer} مخصوص همان \eng{task} است. برای مثال اگر \eng{task} زبانی باشد، محرک زبانی به فرد داده می‌شود و دیده می‌شود کدام بخش‌های مغز فعال می‌شوند. سپس همان بخش‌ها بررسی می‌گردند. در این حالت پژوهشگر لزوماً از قبل نمی‌گوید ناحیهٔ زبان دقیقاً کجاست؛ بلکه بر اساس دادهٔ همان \eng{task} مشخص می‌کند کدام بخش‌ها فعال شده‌اند.

\subsection{شبکه‌های هفت‌گانهٔ \eng{Yeo}}
\eng{Yeo 7-Network Parcellation} یکی از تقسیم‌بندی‌های معروف کارکردی است. در این تقسیم‌بندی، مغز به هفت شبکهٔ کارکردی تقسیم می‌شود. این تقسیم‌بندی \eng{functional} است؛ یعنی بر اساس کارکرد و هم‌فعال‌شدن نواحی مغز ساخته شده، نه صرفاً آناتومی.

\coursefigure{0.82}{yeo-seven-network-parcellation.png}{شبکه‌های هفت‌گانهٔ \eng{Yeo}: افرازبندی کارکردی قشر مغز به شبکه‌های بینایی، حس‌حرکتی، پیش‌فرض، لیمبیک، توجه پشتی، توجه شکمی و فرونتوپاریتال.}

\subsubsection{\eng{Visual Network}}
\eng{Visual Network} مربوط به پردازش دیداری است. محل اصلی آن در بخش پس‌سری مغز است؛ یعنی تقریباً پشت سر، در راستای پردازش ورودی چشم‌ها.

\subsubsection{\eng{Somatomotor Network}}
\eng{Somatomotor Network} به حرکت و حس بدنی مربوط است؛ از جمله حرکت دست و پا، هماهنگی عضلات، انقباض عضلات و کنترل حرکتی بدن.

\subsubsection{\eng{Dorsal Attention Network}}
\eng{Dorsal Attention Network} به توجه درونی و ارادی مربوط است. وقتی خودمان تصمیم می‌گیریم روی چیزی تمرکز کنیم، مثلاً می‌خواهیم ساکت باشیم و متنی را بخوانیم، این شبکه درگیر می‌شود. این نوع توجه از درون هدایت می‌شود.

\subsubsection{\eng{Ventral Attention Network}}
\eng{Ventral Attention Network} به توجه بیرونی و ناگهانی مربوط است. وقتی کسی در می‌زند، صدایی ناگهانی می‌شنویم، یا چیزی از بیرون توجه ما را می‌گیرد، این شبکه فعال می‌شود. این توجه از تصمیم درونی نمی‌آید، بلکه با یک محرک بیرونی تحریک می‌شود.

برای توضیح \eng{dorsal} و \eng{ventral}، در این چارچوب می‌توان گفت اگر بدن یا مغز را در وضعیت جنینی تصور کنیم، \eng{dorsal} به سمت پشت یا بالا و \eng{ventral} به سمت شکم یا پایین اشاره دارد. در مغز نیز \eng{dorsal} معمولاً به بخش‌های بالاتر و پشتی‌تر و \eng{ventral} به بخش‌های پایین‌تر و جلویی‌تر اشاره می‌کند.

\subsubsection{\eng{Frontoparietal Network}}
\eng{Frontoparietal Network} در بخش‌هایی از ناحیهٔ پیشانی و آهیانه‌ای فعال است. کارکرد اصلی آن مدیریت و هماهنگی چند \eng{task}، تصمیم‌گیری، مرتب‌سازی و کنترل رفتار است. نام \eng{frontoparietal} به این معنا نیست که کل بخش پیشانی و کل بخش آهیانه‌ای درگیر است؛ بلکه فقط بخش‌هایی از این نواحی که در این کارکرد فعال می‌شوند، در شبکه قرار می‌گیرند.

\subsubsection{\eng{Default Mode Network}}
\eng{Default Mode Network} زمانی فعال‌تر می‌شود که \eng{task} مشخصی انجام نمی‌دهیم؛ مثل استراحت، تفکر آزاد، خیال‌پردازی، رؤیاپردازی یا ایده‌پردازی آزاد. گاهی وقتی از یک \eng{task} خسته می‌شویم و آن را رها می‌کنیم، ناگهان ایده‌های بهتری به ذهنمان می‌آید. این می‌تواند با فعال‌شدن \eng{Default Mode Network} مرتبط باشد، چون مغز خاموش نشده و وارد حالت پردازش آزادتر شده است.

\subsubsection{\eng{Limbic Network}}
\eng{Limbic Network} با احساسات، هیجان‌ها و پردازش عاطفی مرتبط است.

\subsection{فعالیت شبکه‌ها و \eng{Resting-State fMRI}}
نکتهٔ مهم این است که همهٔ شبکه‌های مغز نمی‌توانند هم‌زمان با حداکثر شدت فعال باشند. مغز درون جمجمه‌ای محدود قرار دارد و جریان خون نیز محدود است. فعالیت مغزی شبیه پینگ‌پنگ میان شبکه‌هاست: بعضی شبکه‌ها فعال‌تر می‌شوند، بعضی پایین می‌آیند، و بسته به \eng{task} الگوی فعالیت تغییر می‌کند. برای مثال نمی‌توان هم‌زمان هم درگیر تفکر آزاد بود و هم با تمرکز کامل یک \eng{task} مشخص را انجام داد.

برای توضیح محدودیت فیزیکی مغز، می‌توان به جراحی مغز اشاره کرد. در برخی جراحی‌ها، بعد از باز کردن جمجمه، ممکن است نتوان بلافاصله استخوان را مثل قبل سر جای خود گذاشت، چون مغز دچار تورم می‌شود. این مثال نشان می‌دهد محدودیت فیزیکی جمجمه و جریان خون در فعالیت مغز نقش واقعی دارند.

در \eng{parcellation} معروف \eng{Yeo} از \eng{resting-state fMRI} استفاده شده است. در \eng{resting-state fMRI} فرد \eng{task} خاصی انجام نمی‌دهد، اما این به معنای خاموش بودن مغز نیست. مغز حتی هنگام استراحت یا خواب فعال است، انرژی زیادی مصرف می‌کند، همچنان پیش‌بینی انجام می‌دهد و شبکه‌های آن کاملاً خاموش نمی‌شوند.

از نظر انرژی نیز همین نکته مهم است: مغز انسان با اینکه فقط بخش کوچکی از وزن بدن را تشکیل می‌دهد، حدود ۲۰ درصد انرژی بدن را مصرف می‌کند. انجام کار شناختی سخت می‌تواند الگوی فعالیت شبکه‌ها را عوض کند، اما افزایش مصرف انرژی نسبت به فعالیت پایه معمولاً کوچک است و در حد چند درصد گزارش می‌شود. بنابراین \eng{resting state} حالت خاموشی نیست؛ بیشتر یک خط پایهٔ فعال و پرهزینه است که \eng{task}ها روی آن تغییرات نسبی ایجاد می‌کنند.

\subsection{همبستگی در \eng{Parcellation}}
برای \eng{parcellation}، \eng{time series} سیگنال \eng{BOLD} وکسل‌ها بررسی می‌شود. اگر دو \eng{voxel} یا دو ناحیهٔ مغزی در طول زمان با هم بالا و پایین شوند، یعنی همبستگی مثبت بالایی داشته باشند، احتمالاً در یک شبکهٔ کارکردی قرار دارند.

روند ساده چنین است:

\begin{enumerate}
  \item برای هر \eng{voxel} یک \eng{time series} گرفته می‌شود.
  \item همبستگی میان \eng{time series}ها محاسبه می‌شود.
  \item \eng{voxel}هایی که همبستگی بالایی دارند \eng{cluster} می‌شوند.
  \item هر \eng{cluster} به‌عنوان یک \eng{parcel} یا \eng{network} در نظر گرفته می‌شود.
\end{enumerate}

در اینجا \eng{clustering} همان \eng{parcellation} است.

\subsection{\eng{Task-Positive} و \eng{Task-Negative}}
در پایان بحث \eng{parcellation}، شبکه‌ها را می‌توان با دو مفهوم توضیح داد. \eng{Task-positive networks} هنگام انجام یک \eng{task} فعال‌تر می‌شوند. شبکه‌های \eng{Dorsal Attention}، \eng{Ventral Attention} و \eng{Frontoparietal} معمولاً در این دسته‌اند.

در مقابل، \eng{task-negative networks} شبکه‌هایی‌اند که هنگام انجام \eng{task} کاهش فعالیت دارند و در حالت عدم انجام \eng{task} فعال‌ترند. این دسته به‌ویژه با \eng{Default Mode Network} مرتبط است.

\coursefigure{0.82}{task-positive-and-task-negative-networks.png}{نمونه‌ای از شبکه‌های وظیفه‌مثبت و وظیفه‌منفی و نیز شبکه‌های توجه پشتی و شکمی؛ شکل رابطهٔ رقابتی و همپوشانی محدود شبکه‌های کارکردی را نشان می‌دهد.}

\subsection{بازگشت به \eng{Brain Alignment}}
پس از \eng{parcellation} می‌توان بازنمایی‌های مغزی را با بازنمایی‌های مدل مقایسه کرد. روند کلی چنین است:

\begin{enumerate}
  \item مغز به \eng{parcel}های کارکردی تقسیم می‌شود.
  \item برای هر \eng{parcel} سیگنال \eng{BOLD} یا بازنمایی زمانی استخراج می‌شود.
  \item از مدل، \eng{hidden state}ها یا بردارهای لایه‌های میانی گرفته می‌شود.
  \item میان بازنمایی‌های مدل و مغز \eng{correlation} یا \eng{similarity} محاسبه می‌شود.
  \item هرچه این ارتباط بیشتر باشد، می‌توان گفت هم‌ترازی بیشتر است.
\end{enumerate}

با این حال، این همچنان اثبات فهم حقیقی نیست؛ فقط راهی عملی برای نزدیک‌شدن به مسئله است.

\subsection{خط پژوهشی آینده: استنتاج و مغالطه}
در پایان، در این چارچوب می‌توان گفت اگر بخواهیم واقعاً بررسی کنیم مدل به معنا و فهم انسانی نزدیک شده یا نه، باید از \eng{task}های سخت‌تر استفاده کنیم: طراحی \eng{task}های استنتاجی، ساختن شرایطی که مدل در آن به اشتباه بیفتد، بررسی مغالطه‌ها، و سنجش \eng{reasoning} انسانی و خطاهای انسانی در مدل.

مغالطه پدیده‌ای انسانی است. اگر مدلی بتواند مغالطه را تشخیص دهد یا حتی ساختار مغالطه را بازتولید کند، یعنی تا حدی ساختار استدلال را فهمیده است. این مسیر یعنی چالش‌کردن مدل از نظر استنتاجی، ادامهٔ طبیعی بحث هم‌ترازی و استدلال است.

\section{استدلال، قیاس و هم‌ترازی خطای انسان و مدل}

\subsection{چرا پاسخ درست کافی نیست؟}
بخش قبل نشان داد که در پس داده‌ها می‌توان یک حقیقت یا مفهوم پنهان را در نظر گرفت؛ چیزی که در این بحث با \(Z\) نشان داده شد. خود \(Z\) را مستقیماً در اختیار نداریم، اما بازنمایی‌هایی از آن داریم: بازنمایی در مغز انسان، بازنمایی در لایه‌های میانی مدل، و خروجی نهایی مدل.

در \eng{Brain Alignment} هدف این است که بفهمیم آیا مدل واقعاً به آن مفهوم پنهان نزدیک شده است یا فقط از نظر آماری خروجی‌هایی شبیه داده‌های آموزشی تولید می‌کند. چون \(Z\) را مستقیم نداریم، مغز انسان به‌عنوان مرجع عملی استفاده می‌شود. \eng{Brain Score} نیز در همین چارچوب می‌پرسد لایه‌های میانی مدل تا چه حد با لایه‌های میانی مغز انسان هم‌بستگی دارند.

اما نکتهٔ مهم این بخش این بود که هم‌ترازی نباید فقط بر اساس پاسخ درست سنجیده شود. اگر انسان‌ها در یک موقعیت خاص خطا می‌کنند، مدلی که واقعاً شبیه انسان می‌فهمد باید در الگوی خطا نیز تا حدی شبیه انسان باشد. بنابراین سؤال فقط این نیست که مدل چند درصد پاسخ درست داده است؛ سؤال‌های مهم‌تر این‌اند:

\begin{itemize}
  \item آیا مدل همان جاهایی اشتباه می‌کند که انسان‌ها اشتباه می‌کنند؟
  \item آیا خطاهای شایع انسانی در مدل نیز شایع‌اند؟
  \item آیا ترجیحات، زمان پاسخ، و توزیع خطاهای مدل به انسان شبیه است؟
\end{itemize}

پس هم‌ترازی می‌تواند چند سطح داشته باشد: هم‌بستگی لایه‌های میانی مدل با مغز، هم‌بستگی خروجی مدل با خروجی انسان، شباهت در زمان پاسخ، شباهت در ترجیحات، و شباهت در خطاها و بایاس‌های انسانی.

\subsection{یادگیری از خطای انسان}
در نگاه اول شاید عجیب باشد که بخواهیم مدل خطاهای انسان را نیز یاد بگیرد. در برخی کاربردها، مثل جراحی رباتیک، هدف دقیقاً این است که مدل خطای انسان را پوشش دهد و بهتر از انسان عمل کند. اما موضوع این بخش ساختن چنین سیستم‌هایی نبود؛ موضوع این بود که آیا مدل واقعاً مفهوم پشت داده را فهمیده است یا نه.

اگر مغز انسان نزدیک‌ترین مرجع ما به فهم مفهوم باشد، خطاهای انسانی نیز بخشی از همان مرجع می‌شوند. از این زاویه، سوگیری شناختی الزاماً باگ نیست. بسیاری از این خطاها در طول تکامل برای بقا مفید بوده‌اند.

برای مثال، انسانی که در موقعیت خطرناک بیش از حد صبر کند تا همه‌چیز را منطقی اثبات کند، ممکن است زنده نماند. در مقابل، تصمیم‌گیری سریع بر اساس پیش‌فرض‌های ذهنی می‌تواند شانس بقا را افزایش دهد. گاهی خطا داشتن یعنی ساده‌سازی و تصمیم‌گیری سریع‌تر.

نمونهٔ دیگر مصرف انرژی مغز است. مغز نمی‌خواهد همهٔ اطلاعات محیط را با جزئیات کامل پردازش کند؛ بنابراین فیلتر می‌گذارد و فقط بخشی از اطلاعات را برجسته می‌کند. وقتی فرد به یک محیط آموزشی نگاه می‌کند، به همهٔ صندلی‌ها و جزئیات محیط توجه نمی‌کند، بلکه بیشتر روی چهره‌ها، واکنش‌ها و صداها تمرکز می‌کند. این فیلترکردن از یک نظر خطاست، اما برای عملکرد طبیعی شناخت لازم است.

حتی طیف بینایی انسان نیز همین منطق را نشان می‌دهد. انسان فقط نور مرئی را می‌بیند، نه کل طیف الکترومغناطیسی. اگر چشم انسان مادون قرمز را نیز می‌دید، شاید دائماً جریان گرما، تغییرات دمای بدن افراد و حرکت هوای گرم و سرد را می‌دید و توجهش مختل می‌شد. حذف برخی اطلاعات ضعف مطلق نیست؛ نوعی بهینه‌سازی است.

\subsection{\eng{Pareidolia} و تشخیص الگو}
\eng{Pareidolia} یعنی دیدن چهره یا الگوی آشنا در چیزهایی که واقعاً چنین الگوی معناداری ندارند؛ مثل ابرها، نقش مبل، گلدوزی یا اشیای بی‌جان. برای مثال، کودکان خردسال گاهی روی نقش مبل یا اشیای دیگر چهره می‌بینند و آن را به انسان نسبت می‌دهند.

این رفتار الزاماً آموخته‌شده نیست، بلکه بخشی از سیستم ادراکی انسان است. از نظر منطقی، دیدن چهره در نقش بی‌جان خطاست؛ اما از نظر تکاملی می‌تواند مفید باشد، چون انسان باید چهره‌ها، نگاه‌ها و واکنش‌های دیگران را سریع تشخیص دهد. اگر کسی وارد یک جمع شود و اصلاً به چهرهٔ افراد واکنش نشان ندهد، برای ما ناخوشایند و غیرطبیعی به نظر می‌رسد.

پس \eng{Pareidolia} نمونه‌ای از خطای ادراکی است که در سطح بقا و تعامل اجتماعی می‌تواند معنا داشته باشد.

\subsection{\eng{System 1}، \eng{System 2} و تفکر انتقادی}
در بحث استدلال دوباره به تفاوت \eng{System 1} و \eng{System 2} اشاره شد. \eng{System 1} سریع، کم‌هزینه و مبتنی بر \eng{prior knowledge} یا دانش قبلی است. این سیستم بر اساس تجربه‌های پیشین و الگوهای آماری سریع تصمیم می‌گیرد.

\eng{System 2} کندتر، پرهزینه‌تر و تحلیلی‌تر است. وقتی فعال می‌شود، تلاش می‌کند مسئله را واقعاً بررسی و اثبات کند. تفکر انتقادی بیشتر به \eng{System 2} مربوط است؛ یعنی فرد در موقعیت‌های لازم می‌تواند از پاسخ سریع و پیش‌فرض ذهنی عبور کند و وارد تحلیل عمیق‌تر شود.

\subsection{انواع \eng{Reasoning}}
\eng{Reasoning} یا استدلال انواع مختلفی دارد. در این بحث سه نوع اصلی نام برده شد:

\begin{itemize}
  \item \eng{Deductive Reasoning}
  \item \eng{Inductive Reasoning}
  \item \eng{Abductive Reasoning}
\end{itemize}

تمرکز این بخش روی شکل خاصی از \eng{Deductive Reasoning} بود: \eng{Syllogistic Reasoning}. در فارسی می‌توان آن را قیاس یا «صغرا و کبرا چیدن» نامید.

\subsection{\eng{Syllogism}}
\eng{Syllogism} ساختاری صوری برای استدلال قیاسی است. در ساده‌ترین حالت، دو مقدمه و یک نتیجه داریم:

\begin{enumerate}
  \item \eng{Premise 1}
  \item \eng{Premise 2}
  \item \eng{Conclusion}
\end{enumerate}

\subsubsection{ساختار قیاس}
در این ساختار یک \eng{Middle Term} یا ترم میانی وجود دارد که در هر دو مقدمه مشترک است و اجازه می‌دهد میان دو گزاره ارتباط برقرار کنیم. اگر ترم مشترک وجود نداشته باشد، اصولاً قیاسی شکل نمی‌گیرد، چون چیزی وجود ندارد که دو مقدمه را به هم وصل کند.

\subsubsection{\eng{Mode A}}
\eng{Mode A} رابطهٔ کلی مثبت است:

\[
\text{\eng{All A are B}}
\]

یعنی همهٔ \(A\)ها \(B\) هستند.

\subsubsection{\eng{Mode I}}
\eng{Mode I} رابطهٔ جزئی مثبت است:

\[
\text{\eng{Some A are B}}
\]

یعنی بعضی از \(A\)ها \(B\) هستند.

\subsubsection{\eng{Mode E}}
\eng{Mode E} رابطهٔ کلی منفی است:

\[
\text{\eng{No A are B}}
\]

یعنی هیچ \(A\)ای \(B\) نیست.

\subsubsection{\eng{Mode O}}
\eng{Mode O} رابطهٔ جزئی منفی است:

\[
\text{\eng{Some A are not B}}
\]

یعنی بعضی از \(A\)ها \(B\) نیستند.

\subsubsection{\eng{Figure}}
علاوه بر \eng{mode}ها، در \eng{Syllogism} چند \eng{figure} مختلف نیز داریم. \eng{Figure} یعنی جایگاه ترم‌ها، به‌ویژه ترم میانی، در دو مقدمه تغییر کند. با چهار \eng{mode} و چهار \eng{figure}، در مجموع \(4 \times 4 \times 4 = 64\) فرم ممکن برای \eng{Syllogism} ساخته می‌شود. بعضی از این فرم‌ها نتیجهٔ معتبر دارند و بعضی ندارند.

در صورت‌بندی رایج مسئله‌های قیاسی، از این ۶۴ فرم، ۲۷ فرم نتیجهٔ معتبر دارند و ۳۷ فرم به \eng{No Valid Conclusion} یا \eng{NVC} می‌رسند. بنابراین ارزیابی مدل فقط این نیست که در چند فرم به پاسخ درست برسد؛ باید دید آیا می‌تواند فرم‌هایی را که اساساً نتیجهٔ معتبر ندارند نیز از فرم‌های معتبر جدا کند یا نه.

چهار شکل اصلی را می‌توان به‌صورت زیر خلاصه کرد:
\[
\begin{array}{c@{\qquad}c@{\qquad}c@{\qquad}c}
\text{\eng{Figure 1}} & \text{\eng{Figure 2}} & \text{\eng{Figure 3}} & \text{\eng{Figure 4}}\\
A-B & B-A & A-B & B-A\\
B-C & C-B & C-B & B-C
\end{array}
\]

\subsubsection{\eng{Valid Conclusion}}
یک نتیجه زمانی \eng{valid} است که در همهٔ حالت‌های ممکن درست باشد و هیچ مثال نقضی نداشته باشد. فرم معروف \eng{Barbara} یا \((AAA)-1\) از ساده‌ترین فرم‌های معتبر است:

\begin{align*}
\text{\eng{All A are B}}\\
\text{\eng{All B are C}}\\
\therefore \quad \text{\eng{All A are C}}
\end{align*}

این نتیجه همیشه معتبر است. اما بسیاری از فرم‌ها نتیجهٔ معتبر ندارند. ممکن است در بعضی حالت‌ها یک نتیجه درست از آب دربیاید، اما اگر همیشه درست نباشد، \eng{valid} محسوب نمی‌شود. چنین مواردی می‌توانند منشأ مغالطه باشند.

\subsection{استفاده از قیاس برای \eng{Brain Alignment}}
پرسش اصلی این است که ساختارهای قیاسی چه ربطی به \eng{Brain Alignment} و مدل‌های زبانی دارند. می‌توان ۶۴ فرم \eng{Syllogism} را به انسان‌ها داد و دید در کدام فرم‌ها درست پاسخ می‌دهند و در کدام فرم‌ها خطا می‌کنند. سپس همان فرم‌ها را به مدل زبانی داد و خروجی مدل را با انسان مقایسه کرد.

در اینجا دقت کلی کافی نیست. مهم‌تر این است که توزیع خطاهای مدل شبیه توزیع خطاهای انسان باشد. اگر انسان‌ها در یک فرم خاص زیاد اشتباه می‌کنند، ولی مدل همیشه درست جواب می‌دهد، این الزاماً نشانهٔ هم‌ترازی نیست. ممکن است مدل فقط ساختارهای صوری را از دادهٔ آموزشی یاد گرفته باشد، بدون اینکه شبیه انسان استدلال کند.

\subsection{نتایج انسانی در \eng{Syllogistic Reasoning}}
در یک مطالعهٔ مرتبط، همین ۶۴ فرم \eng{Syllogism} به انسان‌ها داده شده بود تا مشخص شود انسان‌ها در فرم‌های دارای نتیجهٔ معتبر و فرم‌های \eng{NVC} چه الگوی پاسخی دارند.

\coursefigure{0.95}{analogy-accuracy-by-model-scale.png}{دقت مدل‌ها در فرم‌های مختلف قیاس با مقیاس‌های گوناگون؛ نمودار نشان می‌دهد افزایش اندازهٔ مدل می‌تواند دقت را از خط پایهٔ انسانی بالاتر ببرد.}

برای هر فرم، توزیع پاسخ انسان‌ها بررسی شده بود. حتی در فرم‌های ساده‌ای مثل \eng{Barbara} هم همهٔ افراد پاسخ درست نمی‌دهند: تعداد زیادی درست جواب می‌دهند، اما همچنان بخشی از افراد پاسخ اشتباه تولید می‌کنند. این نشان می‌دهد که حتی در قیاس‌های صوری ساده نیز انسان‌ها همیشه کاملاً منطقی عمل نمی‌کنند.

در این بحث این پرسش مطرح شد که چرا با وجود شناخته‌شدن این فرم‌ها در طول قرن‌ها، انسان‌ها همچنان در بسیاری از آن‌ها اشتباه می‌کنند. چند پاسخ ممکن مطرح شد: انسان در زندگی روزمره به همهٔ این فرم‌ها نیاز نداشته، برخی فرم‌ها کمتر در تجربهٔ واقعی استفاده شده‌اند، بعضی فرم‌های غلط یا مغالطه‌آمیز ممکن است در تعاملات اجتماعی کاربرد داشته باشند، و مغز انسان برای حل همهٔ فرم‌های منطقی بهینه نشده، بلکه برای تصمیم‌گیری سریع در جهان واقعی بهینه شده است.

بنابراین ممکن است انسان بتواند با تمرین همهٔ این فرم‌ها را یاد بگیرد، اما مغز به‌صورت طبیعی خودش را درگیر تمام حالت‌های ممکن نکرده است.

در یک فرم سخت‌تر، چند پاسخ مختلف ممکن است پیشنهاد شود، اما پاسخ معتبر ابتدا به‌راحتی پیدا نشد. نمودار ون نشان می‌دهد که برای \eng{valid} بودن، نتیجه باید در همهٔ حالت‌ها درست باشد. اگر فقط در بعضی حالت‌ها درست باشد، نتیجه معتبر نیست.

در \eng{Syllogism}های صوری، معمولاً بحث مجموعهٔ تهی وارد نمی‌شود، چون هدف بررسی قیاس رسمی است، نه همهٔ پیچیدگی‌های زبان طبیعی یا هستی‌شناسی. اما وقتی این ساختارها به زبان طبیعی برده می‌شوند، مسائلی مانند تهی بودن یک مجموعه می‌توانند مهم شوند.

\subsection{از فرم صوری به زبان طبیعی}
در حالت صوری، جمله‌ها به شکل‌هایی مانند \(\text{\eng{All A are B}}\) و \(\text{\eng{Some B are C}}\) نوشته می‌شوند. اما برای مقایسهٔ واقعی با انسان و مدل زبانی، بهتر است این‌ها به \eng{Natural Language} تبدیل شوند؛ یعنی به‌جای \(A\)، \(B\) و \(C\)، جمله‌های طبیعی و معنادار داشته باشیم.

در این حالت کار سخت‌تر می‌شود، چون مدل باید ترم میانی را تشخیص دهد، رابطهٔ منطقی میان دو مقدمه را پیدا کند، مغالطه‌های زبانی را تشخیص دهد، و نتیجهٔ معتبر یا نبود نتیجهٔ معتبر را تولید کند.

\subsection{دیتاست \eng{Avicenna}}
دیتاستی به نام \eng{Avicenna} حدود ۶۰۰۰ نمونهٔ \eng{Syllogism} در زبان طبیعی دارد.

در هر نمونه معمولاً این اجزا وجود دارد:

\begin{itemize}
  \item \eng{Premise 1}
  \item \eng{Premise 2}
  \item \eng{Syllogistic Relation}
  \item \eng{Conclusion} یا \eng{No Valid Conclusion}
\end{itemize}

\eng{Syllogistic Relation} مشخص می‌کند آیا واقعاً بین دو مقدمه ترم میانی معناداری وجود دارد یا نه. گاهی دو جمله از نظر واژگانی کلمات مشترک دارند، اما این اشتراک برای استدلال کافی نیست. مثلاً ممکن است در یک جمله «شیر جنگل» و در جملهٔ دیگر «شیر آب» آمده باشد. واژه مشترک است، اما ترم میانی واقعی وجود ندارد؛ پس استدلال معتبر شکل نمی‌گیرد.

\subsection{تفاوت \eng{Syllogism} و \eng{Entailment}}
در \eng{Entailment} ممکن است فقط یک گزاره داشته باشیم و با کمک دانش عمومی نتیجه‌ای بگیریم. برای مثال، اگر جمله‌ای دربارهٔ سیگار و آرسنیک باشد، ممکن است با دانش قبلی بدانیم آرسنیک ماده‌ای خاص است و از آن نتیجه‌ای بگیریم.

اما در \eng{Syllogism}، دو مقدمه باید به‌واسطهٔ یک ترم میانی به هم وصل شوند و نتیجه از ترکیب آن دو مقدمه به دست بیاید. پس هر نتیجه‌گیری زبانی الزاماً \eng{Syllogism} نیست.

\subsection{مغالطه‌ها و استدلال‌های روزمره}
برای نشان دادن نقش \eng{prior knowledge} و احساسات در استدلال، دو مثال روزمره مطرح شد.

مثال اول:

\begin{enumerate}
  \item دولت‌های فاسد تورم ایجاد می‌کنند.
  \item کشور ایران تورم دارد.
  \item پس دولت ایران فاسد است.
\end{enumerate}

این نتیجه از نظر منطقی معتبر نیست، چون تورم می‌تواند علت‌های دیگری هم داشته باشد. این ساختار شبیه مغالطهٔ تأیید تالی است. اما اگر فرد از قبل نسبت به دولت نارضایتی داشته باشد، نتیجه برای او پذیرفتنی به نظر می‌رسد. در اینجا \eng{prior knowledge} و احساسات قبلی باعث می‌شوند \eng{System 1} فعال شود و فرد وارد بررسی دقیق \eng{System 2} نشود.

مثال دوم:

\begin{enumerate}
  \item هر کشوری که تولید ناخالص داخلی‌اش رشد کند، اقتصادش بزرگ‌تر شده است.
  \item طبق آمار جهانی، تولید ناخالص داخلی ایران رشد کرده است.
  \item پس اقتصاد ایران بزرگ‌تر شده است.
\end{enumerate}

از نظر ساختار قیاسی، این نتیجه معتبر است. اما ممکن است فرد به دلیل تجربهٔ شخصی یا وضعیت اقتصادی اطرافش آن را نپذیرد. اینجا نتیجهٔ منطقی معتبر است، اما \eng{prior knowledge} یا تجربهٔ زیسته باعث مقاومت در برابر پذیرش آن می‌شود.

\subsection{مقایسهٔ مدل‌های زبانی و انسان در قیاس}
در بخش پایانی، مقاله‌ای بررسی شد که مدل‌های زبانی را روی فرم‌های صوری \eng{Syllogism} با انسان‌ها مقایسه کرده بود. مدل‌ها با فرم‌هایی مانند \(\text{\eng{All A are B}}\)، \(\text{\eng{Some B are C}}\) و حالت‌های مشابه آزمون شده بودند، نه الزاماً با زبان طبیعی پیچیده.

نتیجهٔ مهم این بود که با بزرگ‌تر شدن مدل‌ها، دقت آن‌ها در برخی موارد از انسان بیشتر می‌شود. اما این به معنای هم‌ترازی کامل با انسان نیست. مدل‌های بزرگ ممکن است پاسخ درست‌تری بدهند، ولی توزیع خطاهایشان شبیه انسان نباشد. یعنی جاهایی که انسان‌ها اشتباه می‌کنند، مدل الزاماً اشتباه نمی‌کند، و جاهایی که مدل خطا دارد، الزاماً با خطاهای انسانی منطبق نیست.

\coursefigure{0.95}{analogy-response-pattern-correlation.png}{همبستگی الگوی پاسخ‌ها در فرم‌های قیاسی مختلف؛ این نمودار برای تحلیل هم‌ترازی خطاها مهم‌تر از دقت خام است، چون شباهت رفتار مدل با الگوهای انسانی را نشان می‌دهد.}

پس عملکرد بهتر از انسان لزوماً به معنای فهم انسانی یا هم‌ترازی با انسان نیست. ممکن است مدل به دلیل حافظهٔ بهتر، دادهٔ آموزشی بیشتر یا دیدن نمونه‌های مشابه در آموزش، پاسخ درست بدهد. این با فهم انسانی یکی نیست.

مدل‌های زبانی معمولاً با متن‌های درست، استدلال‌های رسمی و داده‌های گسترده آموزش می‌بینند. در نتیجه ممکن است فرم‌های درست را زیاد دیده باشند. اما انسان در جهان واقعی با مغالطه، ابهام، تجربهٔ شخصی، بایاس و تصمیم‌گیری سریع روبه‌روست. بنابراین اگر مدل فقط در فرم صوری بهتر عمل کند، هنوز نمی‌توان گفت مثل انسان استدلال می‌کند.

\subsection{حافظه، آمار و استدلال}
یک پرسش پایانی این است که آیا مدل فقط حفظ کرده است یا واقعاً استدلال می‌کند. نکتهٔ مهم این است که «حفظ کردن» تعبیر دقیقی نیست. هم انسان و هم مدل به‌نوعی آماری یاد می‌گیرند. انسان نیز بر اساس تجربه‌های قبلی، \eng{prior}ها و توزیع‌های آماری جهان تصمیم می‌گیرد.

اما آمار مدل با آمار تجربهٔ انسانی یکی نیست. اگر بخواهیم مدل واقعاً به انسان نزدیک‌تر شود، باید توزیع دادهٔ آموزشی و خطاهای آن را به سمت خطاها و الگوهای انسانی تنظیم کنیم.

یکی از راه‌ها \eng{Adversarial Data Augmentation} است: ببینیم مدل در چه نوع نمونه‌هایی شکست می‌خورد، سپس همان نوع نمونه‌ها را بیشتر به داده اضافه کنیم تا مدل روی آن‌ها بهتر شود. در دیتاست \eng{Avicenna} نیز چنین کاری انجام شده بود: فرم‌هایی که مدل در آن‌ها مشکل داشت شناسایی شدند و داده‌های مشابه آن‌ها تقویت شدند. در نتیجه عملکرد مدل بهتر شد، اما این بهبود صرفاً با افزایش تصادفی داده حاصل نشده بود؛ بلکه بر اساس تحلیل خطا انجام شده بود.

جمع‌بندی این بخش این است که برای فهم واقعی مدل، نباید فقط دقت نهایی را سنجید. باید دید آیا مدل در مسیر رسیدن به پاسخ، شبیه انسان عمل می‌کند یا نه. در استدلال قیاسی، انسان‌ها الگوی خطای خاصی دارند و این خطاها تصادفی نیستند؛ بلکه به \eng{prior knowledge}، \eng{System 1}، \eng{System 2}، تجربه، بایاس و ساختار شناختی انسان مربوط‌اند.

اگر مدل در پاسخ نهایی از انسان بهتر باشد اما توزیع خطایش با انسان هم‌تراز نباشد، هنوز نمی‌توان گفت واقعاً با انسان \eng{aligned} است. بنابراین در \eng{Brain Alignment} علاوه بر پاسخ درست، باید به خطاهای انسان، خطاهای مدل، توزیع خطاها، توضیح یا مسیر استدلال انسان، لایه‌های میانی مدل، اثر \eng{prior knowledge} و تفاوت میان \eng{reasoning} صوری و \eng{reasoning} در جهان واقعی توجه کرد.


\end{document}
